

\documentclass[bahasa, 12pt, a4paper, onecolumn, oneside, final]{report}



\special{papersize=210mm,297mm}

\usepackage[top=2cm,bottom=2cm,left=4cm,right=3cm]{geometry}

\usepackage{mathptmx}
\usepackage{amsmath}
\usepackage{amssymb}

\usepackage{titlesec}
\newcommand{\subsubsubsection}[1]{\paragraph{#1}\mbox{}\\}

\titleformat{\subsubsubsection}
{\normalfont\normalsize\bfseries}{\theparagraph}{1em}{}
\titlespacing*{\subsubsubsection}
{0pt}{3.25ex plus 1ex minus .2ex}{1.5ex plus .2ex}

\usepackage[utf8]{inputenc}

\usepackage[bahasa]{babel}
\usepackage{csquotes}
\DeclareQuoteAlias{american}{bahasa}

\usepackage[style=apa,backend=biber]{biblatex}
\DefineBibliographyStrings{bahasa}{number = {Nomor},
	jourarticle = {artikel jurnal},
}

\usepackage{longtable}
\usepackage{multirow}
\usepackage{graphicx}
\usepackage{subfig}
\usepackage{titling}
\usepackage{booktabs}
\usepackage{tabularx}
\usepackage{placeins}
\usepackage{blindtext}
\usepackage{sectsty}
\usepackage{chngcntr}
\usepackage{etoolbox}
\usepackage{array}
\usepackage[hidelinks]{hyperref}       \usepackage{titlesec}       \usepackage{titletoc}       \usepackage{tocbibind}      

\usepackage{silence}
\WarningFilter{scrwfile}{THIS IS AN EXPERIMENTAL PACKAGE!}
\WarningFilter{sectsty}{Command \underbar  has changed}
\WarningFilter{sectsty}{Command \underline  has changed}
\usepackage{scrwfile}       \usepackage{parskip}
\usepackage{afterpage}
\usepackage{relsize}
\usepackage{xcolor, colortbl}
\usepackage{setspace}
\usepackage{enumitem}
\usepackage{listings}
\usepackage{listings}
\usepackage{xcolor}

\definecolor{codegreen}{rgb}{0,0.6,0}
\definecolor{codegray}{rgb}{0.5,0.5,0.5}
\definecolor{codepurple}{rgb}{0.58,0,0.82}
\definecolor{backcolour}{rgb}{0.98,0.98,0.98}
\definecolor{codeblue}{rgb}{0.13,0.13,1}

\lstdefinelanguage{Move}{
keywords={public, struct, has, store, key, drop, copy, fun, let, mut, return, abort, module, use, entry, acquires, as, break, const, continue, else, false, friend, if, loop, macro, native, phantom, spec, true, while, inline},
	keywordstyle=\color{codeblue}\bfseries,
ndkeywords={address, bool, u8, u16, u32, u64, u128, u256, vector, signer, String, Option},
	ndkeywordstyle=\color{codepurple}\bfseries,
	sensitive=true,
	comment=[l]{//},
	morecomment=[s]{/*}{*/},
	commentstyle=\color{codegreen}\itshape,
	stringstyle=\color{red}\ttfamily,
	morestring=[b]"
}

\lstdefinelanguage{Rust}{
keywords={as, break, const, continue, crate, else, enum, extern, false, fn, for, if, impl, in, let, loop, match, mod, move, mut, pub, ref, return, self, Self, static, struct, super, trait, true, type, unsafe, use, where, while, async, await, dyn, union},
	keywordstyle=\color{codeblue}\bfseries,
ndkeywords={i8, i16, i32, i64, i128, isize, u8, u16, u32, u64, u128, usize, f32, f64, bool, char, str, String, Vec, Option, Result, Box, Rc, Arc, Cell, RefCell, HashMap},
	ndkeywordstyle=\color{codepurple}\bfseries,
	sensitive=true,
	comment=[l]{//},
	morecomment=[s]{/*}{*/},
morecomment=[l]{///},
	commentstyle=\color{codegreen}\itshape,
	stringstyle=\color{red}\ttfamily,
	morestring=[b]"
}

\lstset{
	backgroundcolor=\color{backcolour}, basicstyle=\ttfamily\small,         columns=fullflexible,               breakatwhitespace=false,
	breaklines=true,                    captionpos=b,                       keepspaces=true,                    frame=single,                       rulecolor=\color{black!30},         numbers=left,                       numbersep=10pt,                     numberstyle=\tiny\color{codegray},  showspaces=false,
	showstringspaces=false,             showtabs=false,
	tabsize=4,                          xleftmargin=2em,                    framexleftmargin=1.5em              } 
\graphicspath{{resources/}}   

\newcommand*{\lopname}{DAFTAR LAMPIRAN}
\TOCclone[\lopname]{toc}{atoc}
\addtocontents{atoc}{\protect\value{tocdepth}=-1}
\newcommand\listofappendices{
	\cleardoublepage
	\phantomsection
	\listofatoc
	\addcontentsline{toc}{chapter}{\lopname}
}

\newcommand*\savedtocdepth{}
\AtBeginDocument{\edef\savedtocdepth{\the\value{tocdepth}}}

\let\originalappendix\appendix
\renewcommand\appendix{\originalappendix
	\cleardoublepage
	\addtocontents{toc}{\protect\value{tocdepth}=-1}\addtocontents{atoc}{\protect\value{tocdepth}=\savedtocdepth}

	\titlecontents{chapter}
	[0pt]
	{\bfseries}
	{Lampiran \thecontentslabel.\quad}
	{}
	{\hfill\contentspage}

	\titleformat{\chapter}[block]
	{\bfseries}
	{\chaptertitlename\ \thechapter.\quad}{0pt}
	{\bfseries}
}

\makeatletter
\renewcommand{\@dotsep}{1}
\makeatother

\titlecontents{chapter}
[0pt]
{\bfseries}
{\MakeUppercase{Bab} \thecontentslabel\quad\uppercase}
{}
{\mdseries\titlerule*[0.35em]{.}\bfseries\contentspage}
\titlecontents{figure}
[0pt]
{}
{Gambar \thecontentslabel.\quad}
{}
{\mdseries\titlerule*[0.35em]{.}\bfseries\contentspage}
\titlecontents{table}
[0pt]
{}
{Tabel \thecontentslabel.\quad}
{}
{\mdseries\titlerule*[0.35em]{.}\bfseries\contentspage}

\let\originalprintbibliography\printbibliography
\renewcommand\printbibliography{\phantomsection
	\cleardoublepage
	\originalprintbibliography
	\addcontentsline{toc}{chapter}{\bibname}
}

\renewcommand{\baselinestretch}{1.5}

\chapterfont{\centering \large}
\titleformat{\chapter}[display]
{\Large\centering\bfseries}
{\chaptertitlename\ \thechapter}{0pt}
{\Large\bfseries\uppercase}

\setcounter{tocdepth}{3}
\setcounter{secnumdepth}{4}

\makeatletter

\makeatother

\counterwithin{figure}{chapter}
\counterwithin{table}{chapter}

\newcommand*{\blankpage}{\afterpage{\null\newpage}}

\renewcommand*{\equationautorefname}{Persamaan}\renewcommand*{\footnoteautorefname}{catatan kaki}\renewcommand*{\itemautorefname}{item}\renewcommand*{\figureautorefname}{Gambar}\renewcommand*{\tableautorefname}{Tabel}\renewcommand*{\partautorefname}{Bagian}\renewcommand*{\appendixautorefname}{Lampiran}\renewcommand*{\chapterautorefname}{Bab}\renewcommand*{\sectionautorefname}{Subbab}\renewcommand*{\subsectionautorefname}{Subsubbab}\renewcommand*{\subsubsectionautorefname}{Subsubsubbab}\renewcommand*{\paragraphautorefname}{paragraf}\renewcommand*{\subparagraphautorefname}{subparagraf}\renewcommand*{\FancyVerbLineautorefname}{garis}\renewcommand*{\theoremautorefname}{Teorema}\renewcommand*{\pageautorefname}{halaman}

\hbadness=99999
\hyphenpenalty=10000
\exhyphenpenalty=10000
\tolerance=9999
\setlength{\emergencystretch}{8em} 

\hyphenation {
a-na-li-sa
	a-pli-ka-si
	a-lo-ka-si
	an-ta-ra
	ada-nya
	akan
be-be-ra-pa
	ber-ge-rak
	be-ri-kut
	ber-ko-mu-ni-ka-si
	buah
	Bouvet
	ber-fo-kus
	ber-fung-si
	ber-ja-lan
ca-ri
	Carzaniga
	cloud
	CloudFormation
	con-tai-ner
	ClusterIP
da-e-rah
	di-nya-ta-kan
	de-fi-ni-si
	di-bu-tuh-kan
	di-gu-na-kan
	di-tam-bah-kan-nya
	di-tem-pat-kan
	di-la-ku-kan
	di-kem-bang-kan
	di-im-ple-men-ta-si-kan
	da-pat
	di-ka-te-go-ri-kan
	de-ngan
	di-kem-bang-kan
	da-ta
e-ner-gi
	eks-klu-sif
	eks-ter-nal
fa-si-li-tas
ga-bung-an
ha-lang-an
	ha-sil
	hell
i-nduk
	in-for-ma-si
	im-ple-men-tasi
kom-po-si-si
	ka-re-na
	ke-sa-ba-ran-nya
	ka-me-ra
	kua-li-tas
	ke-na-ngan
	kom-plek-si-tas
	ke-ti-ka
	ke-le-bi-han-nya
	ke-gi-a-tan
	ko-mu-ni-ka-si
	ke-cil
la-ya-nan
me-ngu-ra-ngi
	meng-eva-lu-a-si
	me-nge-lo-la
	men-da-lam
	men-ja-lan-kan
	mak-si-mal
	me-nye-le-sai-kan
	me-ngunjungi
	men-du-kung
	me-nu-rut
	me-la-ku-kan
	mem-buat
	men-daftar-kan
	me-ngu-sul-kan
	me-mi-li-ki
	meng-gu-na-kan
	men-ja-di
	me-ru-pa-kan
	men-ja-ga
	me-mu-dah-kan
	me-ne-rus-
	mem-pro-ses
Na-mun
ob-so-lete
	or-kes-tra-si
	oto-ma-ti-sa-si
pro-vi-der
	pe-ru-sa-ha-an
	pe-rang-kat
	pro-ses
	plat-form
	pro-duk-si
	pe-ne-li-tian
	pe-ru-ba-han
	pa-ra-dig-ma
	pe-man-tau-an
	pe-ngum-pu-lan
	pa-ckage
quality
se-la-in
	stan-dar-di-sasi
	se-cu-ri-ty
	so-lu-si
	se-lu-ruh
	Soft-ware
	soft-ware
	se-buah
	se-ca-ra
ter-li-bat
	ter-pi-sah-kan
	ter-mi-nal
	ter-ba-tas
un-tuk
}
 

\hyphenation {
con-strained
de-ploy-ment
	de-activation
re-source
	re-mote
zoo-keeper
}
 \usepackage{pdfpages}
\usepackage{pdflscape}

\makeatletter
\let\thesis@defaultfigure\figure
\let\endthesis@defaultfigure\endfigure
\let\thesis@defaulttable\table
\let\endthesis@defaulttable\endtable
\newcommand{\usemainfloatplacements}{\renewenvironment{figure}[1][!b]{\@float{figure}[!b]}{\end@float}\renewenvironment{table}[1][!t]{\@float{table}[!t]}{\end@float}}
\newcommand{\restoredefaultfloatplacements}{\let\figure\thesis@defaultfigure
  \let\endfigure\endthesis@defaultfigure
  \let\table\thesis@defaulttable
  \let\endtable\endthesis@defaulttable
}

\makeatother

\addbibresource{references.ta.bib}

\begin{document}

\title{Pengembangan Kontrol Akses Granular dan Delegasi \textit{Offline} pada Sistem Rekam Medis Elektronik DecMed Menggunakan Macaroons}
\date{}
\author{
	Tazkia Nizami \\
	NIM: 13522032
}
\newcommand\tanggalpengesahan{2 Juli 2026}

\pagenumbering{roman}
\setcounter{page}{1}

\clearpage
\pagestyle{empty}

\begin{center}
	\smallskip

	\Large \bfseries \MakeUppercase{\thetitle}
	\vfill



	\large Disusun sebagai syarat kelulusan tingkat sarjana
	\vfill

	\large Oleh

	\Large \theauthor

	\vfill
	\begin{figure}[ht]
		\centering
		
\includegraphics[width=0.15\textwidth]{cover-ganesha.jpg}
	\end{figure}
	\vfill

	\large
	\uppercase{
		Program Studi Teknik Informatika \\
		Sekolah Teknik Elektro \& Informatika \\
		Institut Teknologi Bandung
	}

	Juli 2026

\end{center}

\clearpage
 \clearpage
\pagestyle{empty}

\begin{center}
	\smallskip

	\Large \bfseries \MakeUppercase{\thetitle}
	\vfill



	\large Oleh

	\Large \theauthor

	\large Program Studi Teknik Informatika \\

	\normalsize \normalfont
	Sekolah Teknik Elektro dan Informatika \\
	Institut Teknologi Bandung \\

	\vfill
	\normalsize \normalfont
	Telah disetujui dan disahkan sebagai Tugas Akhir \\
di Bandung, pada tanggal \tanggalpengesahan

	\vspace{0.5cm}
	Pembimbing,

	\vfill
	\underline{Dr. Phil. Eng. Hari Purnama, S.Si., M.Si.} \\
	NIP. 118110072

\end{center}
\clearpage
 \chapter*{Lembar Pernyataan}

Dengan ini saya menyatakan bahwa:

\begin{enumerate}

	\item Pengerjaan dan penulisan Laporan Tugas Akhir ini dilakukan tanpa menggunakan bantuan yang tidak dibenarkan.
	\item Segala bentuk kutipan dan acuan terhadap tulisan orang lain yang digunakan di dalam penyusunan laporan tugas akhir ini telah dituliskan dengan baik dan benar.
	\item Laporan Tugas Akhir ini belum pernah diajukan pada program pendidikan di perguruan tinggi mana pun.

\end{enumerate}

Jika terbukti melanggar hal-hal di atas, saya bersedia dikenakan sanksi sesuai dengan Peraturan Akademik dan Kemahasiswaan Institut Teknologi Bandung bagian Penegakan Norma Akademik dan Kemahasiswaan khususnya Pasal 2.1 dan Pasal 2.2.
\vspace{15mm}

Bandung, \tanggalpengesahan

\vspace{0.5cm}
\includegraphics[width=4cm]{resources/TTD.png}\\[-0.2cm]
\theauthor
 
\pagestyle{plain}

\clearpage
\chapter*{ABSTRAK}
\addcontentsline{toc}{chapter}{ABSTRAK}
\begin{center}
	\center
	\begin{singlespace}
		\large\bfseries\MakeUppercase{\thetitle}

		\normalfont\normalsize
		Oleh:

		\bfseries \theauthor
	\end{singlespace}
\end{center}

\begin{singlespace}
	\small
	Mekanisme delegasi pada DecMed sebelumnya masih bergantung pada keterlibatan langsung pasien untuk setiap penerbitan token, sementara data rekam medis elektronik belum dipisah ke unit akses kecil sehingga kontrol akses granular sulit ditegakkan pada alur klinis kolaboratif. Tugas akhir ini mengembangkan pembaruan DecMed melalui integrasi Macaroons, segmentasi data \textit{off-chain}, dan verifikasi otoritas pada PRE Server agar delegasi dapat berlangsung dengan atenuasi akses yang tetap terikat konteks layanan. Solusi dirancang dengan \textit{caveat} yang konsisten untuk \texttt{patient\_address}, \texttt{related\_rme\_id}, \texttt{dataset\_category}, \texttt{function\_category}, \texttt{purpose}, \texttt{expires\_before}, serta rantai delegasi (\texttt{root\_subject}, \texttt{delegated\_by}, \texttt{delegated\_to}) dan pembatasan \texttt{max\_delegation\_depth}, lalu dihubungkan dengan mekanisme revokasi melalui \texttt{parent\_token\_hash}. Implementasi mencakup \textit{smart contract} IOTA Move untuk metadata, delegasi, audit, dan revokasi, \textit{Hospital Client} serta \textit{Patient Client}, dan penegakan kebijakan di PRE Server sebelum pengambilan data terenkripsi serta proses \textit{re-encryption}. Evaluasi dilakukan melalui pemetaan kebutuhan--implementasi--pengujian pada uji unit dan uji skenario; dari 17 kebutuhan sistem, 14 berstatus terpenuhi dan 3 terpenuhi sebagian, terutama pada cakupan skenario klinis dan batas prototipe. Hasil tersebut menunjukkan bahwa kontrol akses granular, delegasi \textit{offline}, atenuasi akses, revokasi, dan verifikasi otoritas dapat dioperasionalkan secara konsisten pada DecMed tanpa memperluas hak akses token turunan melebihi token induk.

	\textbf{\textit{Kata kunci: kontrol akses granular, delegasi offline, atenuasi akses, Macaroons, rekam medis elektronik, DecMed, PRE Server, revokasi, verifikasi otoritas}}

\end{singlespace}
\clearpage \clearpage
\chapter*{ABSTRACT}
\addcontentsline{toc}{chapter}{ABSTRACT}

\begin{center}
	\center
	\begin{singlespace}
		\large\bfseries\MakeUppercase{\thetitle}

		\normalfont\normalsize
		By:

		\bfseries \theauthor
	\end{singlespace}
\end{center}

\begin{singlespace}
	\small
	DecMed's prior delegation flow required direct patient involvement whenever a new access token had to be issued, while electronic medical record data remained too coarse as a single encrypted payload, making granular enforcement difficult in collaborative clinical workflows. This thesis extends DecMed by integrating Macaroons, off-chain data segmentation, and authority verification on the PRE Server so delegation can proceed through constrained offline attenuation without detaching from clinical context. The design consistently defines caveats for \texttt{patient\_address}, \texttt{related\_rme\_id}, \texttt{dataset\_category}, \texttt{function\_category}, \texttt{purpose}, \texttt{expires\_before}, delegation chain attributes (\texttt{root\_subject}, \texttt{delegated\_by}, \texttt{delegated\_to}), and \texttt{max\_delegation\_depth}, then binds descendant tokens to revocation handling through \texttt{parent\_token\_hash}. Implementation covers IOTA Move smart contracts for metadata, delegation, auditing, and revocation, along with Hospital Client and Patient Client adaptations and PRE-side enforcement before encrypted data retrieval and re-encryption. Evaluation is conducted through requirements-to-implementation-to-test traceability using unit and scenario testing; out of 17 system requirements, 14 are fulfilled and 3 are partially fulfilled, mainly due to prototype scope and limited clinical scenario coverage. These results indicate that granular access control, offline delegation, access attenuation, revocation, and authority verification can be implemented coherently in DecMed while preserving the rule that delegated tokens cannot exceed the privilege scope of their parent tokens.

	\textbf{\textit{Keywords: granular access control, offline delegation, access attenuation, Macaroons, electronic medical records, DecMed, PRE Server, revocation, authority verification}}
\end{singlespace}
\clearpage

\clearpage
 
\chapter*{Kata Pengantar}
\addcontentsline{toc}{chapter}{KATA PENGANTAR}

Puji dan syukur penulis panjatkan kepada Tuhan Yang Maha Esa atas berkat dan rahmatnya, laporan tugas akhir yang berjudul "\thetitle" dapat diselesaikan dalam rangka memenuhi syarat kelulusan tingkat sarjana. Perlu diakui pengerjaan tugas akhir ini didukung oleh banyak pihak. Khususnya, penulis ingin mengucapkan terima kasih kepada:

\begin{enumerate}
	\item Ani Rochani Rochmawati dan Widodo selaku orang tua penulis dan Fathiya Alisa Zahrandika selaku kakak penulis atas dukungan materi dan emosi selama ini hingga penulis berhasil menyelesaikan pengerjaan tugas akhir.
	\item Bapak Dr. Phil. Eng. Hari Purnama, S.Si., M.Si., selaku dosen pembimbing atas segala bentuk dukungan yang telah diberikan dan kesabarannya dalam membimbing penulis serta memberikan saran dalam pengerjaan tugas akhir.
	\item Seluruh pihak lain yang tidak disebutkan disini yang telah membantu dan mendukung dalam proses pengerjaan tugas akhir.
\end{enumerate}

Akhir kata, penulis mengucapkan terima kasih kepada semua pihak yang telah terlibat dalam pengerjaan tugas akhir ini. Penulis juga ingin menyampaikan mohon maaf apabila terdapat kesalahan maupun kekurangan dalam laporan tugas akhir ini. Penulis berharap semoga tugas akhir ini dapat bermanfaat bagi pembaca dan riset-riset kedepannya.

\begin{flushright}
	\vspace{0.5cm}
	Bandung, \tanggalpengesahan

	\includegraphics[width=4cm]{resources/TTD.png}\\[-0.2cm]

	Tazkia Nizami
\end{flushright} 
\titleformat*{\section}{\centering\bfseries\Large\MakeUpperCase}
\titlespacing*{\chapter}{0pt}{0pt}{4pt}

\renewcommand{\contentsname}{DAFTAR ISI}
\renewcommand{\listfigurename}{DAFTAR GAMBAR}
\renewcommand{\listtablename}{DAFTAR TABEL}
\renewcommand{\lstlistingname}{Kode}
\renewcommand{\lstlistlistingname}{DAFTAR KODE PROGRAM}
\renewcommand{\bibname}{DAFTAR REFERENSI}

\tableofcontents
\listofappendices
\listoffigures
\listoftables
\lstlistoflistings

\newpage

\titleformat*{\section}{\bfseries\large}
\pagenumbering{arabic}

\setcounter{page}{1}
\renewcommand{\chaptername}{BAB}
\renewcommand{\thechapter}{\Roman{chapter}}
\usemainfloatplacements


\chapter{Pendahuluan}
\label{chap:pendahuluan}

\section{Latar Belakang}
\label{sec:latarbelakang}

Transformasi sistem rekam medis semakin didorong secara oleh regulasi pemerintah. Peraturan Menteri Kesehatan Republik Indonesia (Permenkes) Nomor 24 Tahun 2022 mewajibkan fasilitas layanan kesehatan (fasyankes) untuk mengimplementasikan Rekam Medis Elektronik (RME) yang terintegrasi \autocite{permenkes24_2022}. Integrasi ini krusial untuk menjamin interoperabilitas data kesehatan nasional sesuai amanat Undang-Undang Nomor 29 Tahun 2004 \autocite{uu_29_2004}.

Perkembangan dari sistem RME memiliki tantangan dalam beberapa aspek, salah satunya adalah terkait kepemilikan data \autocite{liu2025accesscontrol}


Kepemilikan data merupakan elemen kritis dalam manajemen RME, yang menuntut pergeseran paradigma kontrol akses menjadi sistem yang \textit{patient-centric} \autocite{patil2024blockchain, peng2023patientcentric}. Dalam model ini, pasien memiliki kedaulatan untuk menentukan siapa saja yang dapat mengakses informasi kesehatan mereka. Untuk mewujudkan \textit{Data Ownership}, diperlukan sebuah lapisan kontrol akses yang memberikan kontrol yang presisi terhadap data RME \autocite{peng2023patientcentric,patil2024blockchain}.

Model kontrol akses yang paling umum digunakan pada sistem RME adalah \textit{role‑based access control} (RBAC) \autocite{hegde2025safeguarding}. Namun, RBAC cenderung memberikan hak akses berdasarkan peran umum sehingga berisiko menimbulkan \textit{overprivilege}. Di sisi lain, jika granularitas akses dipaksakan melalui penambahan banyak \textit{role}, sistem dapat mengalami \textit{role explosion} \autocite{singh2025granular}.

Dalam praktik klinis, perawatan pasien di rumah sakit bersifat kolaboratif dan melibatkan banyak tenaga kesehatan lintas unit (dokter, petugas laboratorium, dan sebagainya) \autocite{ramaswamy2023multidisciplinarycare, kmk1128_2022}. NHS England mengakui bahwa untuk memberikan perawatan yang baik, informasi yang relevan sebaiknya dibagikan antar tenaga kesehatan yang terlibat dengan dasar \textit{implied consent}, selama menjaga prinsip \textit{need‑to‑know} \autocite{nhs_consent_2025}.
Studi terkait melaporkan bahwa pasien mendukung berbagi data klinis mereka di dalam sistem kesehatan untuk tujuan perawatan langsung tanpa perlu persetujuan eksplisit setiap kali, asalkan mereka diberi informasi yang jelas dan dapat menolak bila tidak setuju \autocite{jones2022public}. Kondisi ini menunjukkan bahwa mekanisme kontrol akses pada RME perlu menyeimbangkan kontrol pasien dengan kebutuhan operasional klinis, terutama ketika akses perlu diteruskan secara terbatas kepada tenaga kesehatan yang terlibat dalam episode perawatan yang sama.

Salah satu solusi sistem RME terdesentralisasi adalah DecMed \autocite{sadewa2025decmed}. Sistem ini memanfaatkan DLT IOTA, \textit{InterPlanetary File System} (IPFS), dan \textit{Proxy Re-Encryption} (PRE) dengan menerapkan \textit{capability-based access control}. Namun, mekanisme pemberian dan pendelegasian akses pada DecMed masih bertumpu pada pasien sebagai \textit{resource owner} untuk menerbitkan token hak akses baru. Dalam konteks layanan klinis yang kolaboratif, asumsi tersebut dapat menjadi hambatan karena kebutuhan akses sering kali muncul secara dinamis sesuai peran tenaga kesehatan dalam episode perawatan. Selain itu, kontrol akses granular tidak cukup hanya direpresentasikan pada token hak akses. Data RME yang dilindungi juga masih disimpan sebagai satu \textit{entry} besar dan perlu dipisahkan ke dalam unit yang lebih kecil agar data dapat diakses sesuai dengan kebutuhan layanan. Dengan demikian, terdapat dua keterbatasan utama yang perlu diperbaiki, yaitu mekanisme delegasi akses yang masih bergantung pada keterlibatan langsung pasien dan struktur data RME yang perlu mendukung pemisahan informasi medis sesuai prinsip \textit{need-to-know}.

Untuk menjawab keterbatasan tersebut, Tugas Akhir ini memperbarui mekanisme kontrol akses pada DecMed melalui integrasi Macaroons, segmentasi data medis, dan fitur \textit{offline attenuation}. Macaroons digunakan untuk membawa batasan akses melalui \textit{caveats} \autocite{birgisson2014macaroons}. Fitur \textit{offline attenuation} memungkinkan pasien menerbitkan token awal pada proses administrasi, lalu pemegang token menurunkan hak akses secara lokal kepada tenaga kesehatan yang terlibat. Integrasi ini diarahkan untuk memperbaiki mekanisme delegasi akses pada DecMed sebelumnya dan kontrol akses yang granular dengan segmentasi data medis dengan menerapkan prinsip \textit{need-to-know}. 
\section{Rumusan Masalah}
\label{sec:rumusan-masalah}

Berdasarkan latar belakang yang telah diuraikan pada subbab \ref{sec:latarbelakang}, rumusan masalah yang dibahas dalam tugas akhir ini adalah sebagai berikut:

\begin{enumerate}
	\item Bagaimana rancangan dan implementasi mekanisme kontrol akses granular pada kerangka kerja DecMed dengan memanfaatkan Macaroons?
	\item Bagaimana rancangan dan implementasi segmentasi data RME yang mendukung penegakan kontrol akses granular pada DecMed?
	\item Bagaimana rancangan dan implementasi fitur \textit{offline attenuation} pada token otorisasi untuk memfasilitasi pendelegasian hak akses antar tenaga kesehatan yang terlibat dalam layanan klinis tanpa memerlukan intervensi langsung dari pasien?
\end{enumerate}

\section{Tujuan}

Tujuan dari Tugas Akhir ini adalah:
\begin{enumerate}
	\item Merancang dan menerapkan kontrol akses granular pada kerangka kerja DecMed dengan memanfaatkan teknologi Macaroons.
	\item Merancang dan menerapkan segmentasi data yang mendukung penegakan kontrol akses granular pada DecMed.
	\item Merancang dan menerapkan fitur atenuasi \textit{offline} untuk memfasilitasi pendelegasian hak akses antar tenaga medis tanpa memerlukan intervensi langsung dari pasien.
\end{enumerate}

\section{Batasan Masalah}

Terdapat beberapa batasan masalah dari Tugas Akhir ini. Batasan masalah tersebut antara lain:

\begin{enumerate}
	\item Kasus skenario delegasi layanan klinis terbatas hanya pada aktor dokter, perawat, petugas laboratorium, dan apoteker pada alur Rawat Jalan dan Rawat Inap.
	\item Data RME pada dataset Instalasi Gawat Darurat (IGD) tidak termasuk dalam kajian rancangan dan implementasi.
	\item Segmentasi data RME hanya dilakukan untuk dokumentasi klinis. Dokumentasi administratif tidak termasuk ke dalam fokus Tugas Akhir ini dan hanya diimplementasikan sesuai dengan implementasi DecMed sebelumnya.
\end{enumerate}

\section{Metodologi}

Terdapat beberapa tahapan untuk melaksanakan tugas akhir secara sistematis,
berikut merupakan tahapan-tahapan yang dilakukan:

\begin{enumerate}
	\item \textbf{Studi Literatur}
	      \label{subsec:studi-literatur}

	      Tahap awal berfokus pada pengumpulan fondasi teoretis dan kontekstual melalui studi literatur.

	\item \textbf{Analisis Kebutuhan}
	      \label{subsec:analisis-kebutuhan}

	      Tahap ini membahas analisis untuk mendefinisikan kekurangan pada kerangka kerja DecMed. Kekurangan tersebut diinterpretasikan sebagai masalah dan didefinisikan menjadi kebutuhan. Tahap ini diakhiri dengan perumusan kebutuhan fungsional dan nonfungsional sistem.

	\item \textbf{Analisis dan Rancangan Solusi}
	      \label{subsec:perancangan-arsitektur-dan-model-kontrol-akses}

	      Tahap ini membahas alternatif solusi yang dapat digunakan untuk mengatasi masalah pada kerangka kerja DecMed. Solusi yang terpilih kemudian direncanakan menjadi rancangan solusi.

	\item \textbf{Implementasi Prototipe}
	      \label{subsec:implementasi-prototipe}

	      Tahap ini merupakan realisasi dari rancangan ke dalam bentuk pengembangan dari prototipe DecMed yang sudah ada.

	\item \textbf{Pengujian}
	      \label{subsec:pengujian}

	      Tahap ini merupakan tahapan untuk menguji prototipe dapat memenuhi kebutuhan sistem sesuai dengan skenario pengujian yang ditentukan.

	\item \textbf{Evaluasi}
	      \label{subsec:evaluasi}

	      Tahap ini merupakan tahapan untuk mengevaluasi prototipe berdasarkan hasil pengujian, yang menjadi dasar dari kesimpulan dan memberikan saran untuk penelitian selanjutnya.
\end{enumerate} 
\section{Sistematika Pembahasan}
\label{sec:sistematika-pembahasan}

Berikut merupakan sistematika pembahasan dari Tugas Akhir ini:

\begin{description}
	\item[Bab I - Pendahuluan]
	      Bab ini menjelaskan gagasan utama dari tugas akhir, yang meliputi latar belakang, rumusan masalah, tujuan, batasan, metodologi, hingga sistematika pembahasan mengenai proses pengembangan solusi.

	\item[Bab II - Studi Literatur]
	      Bab ini memaparkan studi literatur yang telah dilakukan terkait sistem yang dikembangkan. Pembahasan difokuskan pada uraian teoritis berbagai komponen yang digunakan, termasuk regulasi dan teknologi yang digunakan dalam sistem.

	\item[Bab III - Analisis dan Rancangan Solusi]
	      Bab ini dijelaskan analisis terhadap masalah yang ditemukan serta berbagai alternatif solusi yang dapat digunakan untuk mengatasinya. Selain itu, bab ini juga membahas rancangan solusi yang dikembangkan, meliputi pendefinisian kebutuhan fungsional dan nonfungsional, serta rancangan setiap komponen sistem.

	\item[Bab IV - Implementasi dan Pengujian Sistem]
	      Bab ini menjelaskan hasil implementasi dari rancangan solusi yang telah dipaparkan pada Bab III. Selain itu, strategi pengujian serta hasil pengujian yang dilakukan terhadap solusi juga dibahas pada bab ini.

	\item[Bab V - Penutup]
	      Bab ini berisi kesimpulan dari hasil penelitian serta saran untuk pengembangan lebih lanjut.
\end{description}  \chapter{Studi Literatur}
\label{chapter:studi-literatur}

\section{Rekam Medis Elektronik}

\subsection{Definisi RME}
\label{subsec:definisi-rme}

Di Indonesia, definisi dari Rekam Medis dan Rekam Medis Elektronik didefinisikan pada Peraturan Menteri Kesehatan Republik Indonesia Nomor 24 Tahun 2022 (Permenkes 24/2022). "Rekam Medis adalah dokumen yang berisikan identitas pasien, pemeriksaan, pengobatan, tindakan, dan pelayanan lain yang telah diberikan pada pasien" sedangkan "Rekam Medis Elektronik adalah Rekam Medis yang dibuat dengan menggunakan sistem elektronik yang diperuntukkan bagi penyelenggaraan Rekam Medis." \autocite{permenkes24_2022}.

Berdasarkan definisi tersebut, rekam medis adalah salah satu komponen vital dalam sistem kesehatan di Indonesia. Rekam medis bukanlah sebuah catatan biasa yang dibuat lalu disimpan, melainkan sebuah catatan yang terus diperbarui. Keberadaan Rekam Medis bertujuan untuk meningkatkan mutu pelayanan kesehatan dan memberikan kepastian hukum dalam peristiwa medis.

\subsection{Isi RME}
\label{subsec:isi-rme}

Berdasarkan Permenkes 24/2022 Pasal 27 Ayat 1 \autocite{permenkes24_2022}, isi dari Rekam Medis Elektronik terbagi menjadi dua bagian, yaitu:

\begin{enumerate}
	\item Dokumentasi Administratif
	\item Dokumentasi Klinis
\end{enumerate}

Panduan lebih rinci terkait isi dari Rekam Medis Elektronik diatur pada Keputusan Menteri Kesehatan (KMK) No. HK.01.07/MENKES/1423/2022 tentang Pedoman Variabel dan metadata pada Penyelenggaraan Rekam Medis Elektronik \autocite{kemenkes1423-2022}. Pada setiap data RME, terdapat dokumentasi administratif dan dokumentasi klinis. Dokumentasi Klinis kemudian dibagi berdasarkan kategori fungsi klinis. KMK ini mengatur variabel dan metadata yang terdapat pada RME dengan dataset yang terdiri dari:

\begin{enumerate}[label=\arabic*.]
	\item IGD (Instalasi Gawat Darurat).
	\item Rawat Jalan.
	\item Rawat Inap.
	\item Laboratorium.
	\item Apotek.
\end{enumerate}

Detail metadata dataset RME untuk dataset Rawat Jalan, Rawat Inap, Laboratorium, dan Apotek dapat dilihat pada Lampiran \ref{lampiran:metadata-rme}.

\subsection{RME dalam Alur Pelayanan Kesehatan Kolaboratif}
\label{subsec:rme-layanan-kolaboratif}

Penyelenggaraan layanan kesehatan pada praktiknya tidak selalu dilakukan oleh satu aktor tunggal, melainkan melibatkan beberapa tenaga kesehatan dan tenaga lain yang bekerja sesuai kewenangan masing-masing \autocite{who2010_interprofessional_collaborative_practice}. Dalam konteks layanan kolaboratif, pasien dapat berinteraksi dengan petugas rekam medis pada tahap pendaftaran, perawat pada tahap pemeriksaan awal, dokter pada tahap pemeriksaan dan pengambilan keputusan klinis, petugas laboratorium pada tahap pemeriksaan penunjang, serta apoteker atau tenaga kefarmasian pada tahap pelayanan resep dan obat. Pola ini sejalan dengan konsep praktik kolaboratif antarprofesi, yaitu ketika tenaga kesehatan dari latar belakang profesi berbeda bekerja bersama pasien untuk memberikan pelayanan yang bermutu \autocite{ramaswamy2023multidisciplinarycare}.

Struktur RME di Indonesia juga menunjukkan bahwa data kesehatan pasien tidak bersifat tunggal, melainkan terdiri dari dokumentasi administratif dan dokumentasi klinis \autocite{permenkes24_2022}. Pembagian ini menunjukkan bahwa kebutuhan akses terhadap RME dapat berbeda antara aktor administratif dan aktor klinis.

Standar Akreditasi Rumah Sakit memberikan dasar bahwa pelayanan pasien dilakukan melalui kerja kolaboratif antarprofesi. Dalam standar tersebut, DPJP (Dokter Penanggung Jawab Pasien) ditempatkan sebagai \textit{clinical leader}, sedangkan PPA (Profesional Pemberi Asuhan) bekerja sebagai tim interdisiplin dalam proses asuhan pasien \autocite{kmk1128_2022}. Standar yang sama juga menunjukkan keterlibatan perawat dalam pengkajian dan asuhan keperawatan, laboratorium sebagai layanan penunjang untuk memenuhi kebutuhan pasien dan PPA, serta apoteker dalam proses penggunaan obat yang terkoordinasi.

Alur operasional layanan tersebut diperjelas oleh panduan interoperabilitas SATUSEHAT untuk layanan Rawat Jalan dan Rawat Inap. Panduan tersebut memuat tahapan layanan seperti pendaftaran pasien, pendaftaran kunjungan, anamnesis, pemeriksaan, rencana rawat, instruksi medik dan keperawatan, pemeriksaan penunjang, diagnosis, tindakan, peresepan, pengkajian resep, pemberian obat, tindak lanjut, dan resume medis \autocite{satusehat_rme_rawat_jalan,satusehat_rawat_inap_new}. Sementara itu, Pedoman Variabel dan Metadata RME dari Kementerian Kesehatan memetakan dokumentasi RME ke dalam dataset seperti Rawat Jalan, Rawat Inap, Laboratorium, dan Apotek \autocite{kemenkes1423-2022}. Kombinasi antara alur pelayanan, keterlibatan aktor, dan struktur dataset tersebut menunjukkan bahwa kebutuhan akses terhadap RME berbeda menurut tahap layanan, peran aktor, serta jenis data yang digunakan.

\subsection{Hak Akses RME}
\label{subsec:hak-akses-rme}

Karena RME memuat data administratif dan klinis yang termasuk data pribadi spesifik, akses terhadap RME harus dibatasi berdasarkan kewenangan, tindakan, dan kebutuhan pelayanan. Tenaga kesehatan dapat melakukan beberapa aksi terhadap data RME pasien, yaitu \autocite{permenkes24_2022}:

\begin{enumerate}
	\item Menambahkan data baru.
	\item Melakukan perbaikan data.
	\item Melihat atau membaca data.
\end{enumerate}

Undang-Undang Nomor 27 Tahun 2022 tentang Pelindungan Data Pribadi (UU PDP) menjadi landasan yuridis utama yang menetapkan standar keamanan bagi sistem data rekam medis. Pasal 4 ayat 2 huruf a UU PDP mengklasifikasikan data dan informasi kesehatan sebagai data pribadi spesifik \autocite{uu27-2022-pdp}. Pengendali data pribadi diwajibkan untuk melindungi dan memastikan keamanan data yang diprosesnya (Pasal 35), mewajibkan pengendali data pribadi untuk menyediakan \textit{audit trail} pemrosesan data pribadi (Pasal 32), wajib melindungi data dari pemrosesan yang tidak sah (Pasal 38) dan mencegah akses yang tidak sah melalui sistem keamanan yang andal (Pasal 39) \autocite{uu27-2022-pdp}.

Berdasarkan Permenkes No. 73 Tahun 2016 tentang Standar Pelayanan Kefarmasian di Apotek, pelayanan farmasi klinik mencakup pengkajian resep dan \textit{dispensing} \autocite{permenkes73_2016}. Hal ini menunjukkan bahwa apoteker memerlukan akses terhadap informasi yang relevan dengan resep, terapi, dan penggunaan obat pasien untuk memastikan kesesuaian, keamanan, dan ketepatan pelayanan obat.

Kebutuhan akses pada petugas laboratorium dapat diturunkan dari struktur metadata RME pada KMK 1423/2022 \autocite{kemenkes1423-2022}. Pada pedoman tersebut, permintaan pemeriksaan laboratorium yang diminta oleh dokter untuk pasien dicantumkan sebagai bagian dari data pemeriksaan penunjang. Hal ini menunjukkan bahwa dalam alur pelayanan, petugas laboratorium tidak perlu memperoleh akses terhadap seluruh data klinis pasien, tetapi cukup terhadap informasi yang relevan untuk memproses pemeriksaan penunjang terkait laboratorium. \section{IOTA}
\label{sec:iota}

\subsection{Distributed Ledger Technology}
\label{subsec:dlt}

\textit{Distributed Ledger Technology} (DLT) merupakan teknologi pencatatan data terdistribusi yang memungkinkan data direkam, disimpan, dan diverifikasi oleh banyak \textit{node} tanpa bergantung pada satu otoritas terpusat. DLT mencakup berbagai bentuk implementasi, termasuk \textit{blockchain} dan struktur alternatif seperti \textit{Directed Acyclic Graph} (DAG), yang masing-masing memiliki pendekatan berbeda dalam mencapai konsensus, keamanan, dan skalabilitas pada jaringan terdesentralisasi \autocite{savadatti2025survey}.

Karakteristik utama DLT meliputi replikasi data pada banyak \textit{node}, ketahanan terhadap perubahan data melalui mekanisme kriptografi, serta kemampuan verifikasi oleh partisipan jaringan. Karakteristik tersebut membuat DLT banyak digunakan sebagai infrastruktur untuk pencatatan transaksi, penyimpanan referensi data, dan eksekusi aturan terprogram melalui \textit{smart contract} \autocite{savadatti2025survey}. Dalam konteks sistem rekam medis elektronik terdesentralisasi, DLT dapat digunakan untuk mendukung pencatatan data yang dapat diverifikasi tanpa menempatkan seluruh kendali pada satu pihak terpusat.

\subsection{IOTA}
\label{subsec:iota}

IOTA merupakan salah satu implementasi DLT yang dirancang untuk mendukung transaksi pada lingkungan terdesentralisasi, khususnya pada skenario yang membutuhkan skalabilitas tinggi dan biaya transaksi rendah. Berbeda dari \textit{blockchain} konvensional, IOTA menggunakan struktur DAG yang dikenal sebagai \textit{Tangle} sebagai dasar buku besar terdistribusinya \autocite{sedi2025from}. Pada perkembangannya, IOTA mengalami perubahan arsitektur melalui IOTA Rebased. Perubahan ini bertujuan meningkatkan skalabilitas, desentralisasi, dan kemampuan ekosistem IOTA untuk mendukung aplikasi terdesentralisasi, termasuk penggunaan \textit{smart contract} \autocite{sedi2025from}. IOTA menggunakan Move \textit{smart contract} untuk berinteraksi dengan data dan melakukan transaksi.

 \section{\textit{InterPlanetary \textit{File} System} (IPFS)}
\label{sec:ipfs}

IPFS adalah sebuah protokol \textit{peer-to-peer} dan sistem \textit{file} terdistribusi yang dirancang untuk menyimpan dan membagikan data secara desentralisasi \autocite{benet2014ipfs}. Berbeda dengan sistem penyimpanan tradisional yang bergantung pada \textit{server} pusat, IPFS membangun jaringan global dengan setiap pengguna (\textit{node}) dapat berperan sebagai penyimpan dan penyedia konten. Teknologi ini menggunakan \textit{content addressing} daripada \textit{location addressing}, sehingga setiap \textit{file} diidentifikasi dengan unik berdasarkan hash kriptografi dari konten tersebut. Dengan pendekatan desentralisasi ini, IPFS menawarkan solusi yang lebih tahan terhadap sensor, lebih andal, dan tidak bergantung pada satu titik kegagalan pusat.

IPFS beroperasi melalui mekanisme yang mengintegrasikan beberapa komponen kunci. Pertama, menggunakan \textit{Distributed Hash Table} (DHT) yang memungkinkan setiap \textit{file} diberi identitas unik berdasarkan hash konten (\textit{Content Identifier}/CID), memungkinkan \textit{file} diakses dari \textit{node} terdekat yang menyimpannya daripada harus mengakses \textit{server} pusat yang jauh. Kedua, IPFS menerapkan penyimpanan terdistribusi dengan membagi \textit{file} ke dalam blok-blok kecil yang tersebar di berbagai \textit{node} dalam jaringan, sehingga meningkatkan kecepatan akses dan ketahanan terhadap kegagalan. Ketiga, teknologi ini menggunakan sistem \textit{caching} lokal pada setiap \textit{node} yang menyimpan salinan konten yang sering diakses, memungkinkan akses lebih cepat tanpa perlu mengunduh ulang dari internet. Keempat, IPFS menggunakan protokol komunikasi \textit{peer-to-peer} yang memungkinkan setiap \textit{node} bertindak sebagai klien, \textit{server}, atau \textit{node} perantara, sehingga memfasilitasi komunikasi langsung antar \textit{node} tanpa perantara \textit{server} pusat \autocite{benet2014ipfs}. \section{\textit{Proxy Re-Encryption} (PRE)}
\label{sec:proxy-re-encryption}

\textit{Proxy Re-Encryption} adalah skema kriptografi yang memungkinkan pihak ketiga (disebut proxy) mengubah suatu \textit{ciphertext} yang awalnya dienkripsi untuk satu pengguna menjadi \textit{ciphertext} baru yang bisa didekripsi oleh pengguna lain, tanpa pernah membuka \textit{plaintext} dan tanpa memegang kunci privat kedua pihak. Teknik ini dipakai untuk mendelegasikan hak dekripsi secara aman, misalnya pada skenario berbagi data terenkripsi di \textit{cloud}. PRE pertama kali diperkenalkan oleh Blaze et al. \autocite{blaze1998divertible} sebagai solusi untuk mendelegasikan hak dekripsi secara aman.

Gambaran simulasi dari PRE dapat dilihat pada Gambar \ref{fig:basic-reencryption-scheme}. Dalam PRE, pemilik data (misalnya Alice) mengenkripsi datanya dengan kunci publik miliknya dan menyimpannya di server atau cloud. Ketika Alice ingin memberi akses ke Bob, ia membuat \textit{re-encryption key} dari kunci miliknya dan kunci Bob, lalu memberikannya ke proxy agar proxy bisa mengubah \textit{ciphertext} milik Alice menjadi \textit{ciphertext} yang dapat didekripsi Bob, tanpa pernah melihat isi datanya.

\begin{figure}[ht]
	\centering
	\includegraphics[width=0.8\textwidth]{chapter-2/PRE Scheme.png}
	\caption{Gambaran simulasi dari PRE \autocite{ nunez2017umbral}}
	\label{fig:basic-reencryption-scheme}
\end{figure}

\subsection{Umbral PRE}
\label{subsec:umbral-pre}

Umbral adalah skema \textit{threshold proxy re-encryption} yang memungkinkan pemilik data mendelegasikan hak dekripsi kepada penerima akses melalui satu atau beberapa proxy \textit{semi-trusted} \autocite{nunez2017umbral}. Dalam Umbral, pihak yang memberikan akses disebut \textit{delegator}, pihak yang menerima akses disebut \textit{delegatee}, sedangkan proxy bertugas melakukan re-enkripsi tanpa mengetahui isi data. Berbeda dengan PRE sederhana yang dapat menggunakan satu \textit{re-encryption key}, Umbral membagi kemampuan re-enkripsi ke dalam beberapa \textit{key fragments} atau \texttt{kfrag}. Dengan pendekatan ini, proses re-enkripsi dapat dibuat berbasis ambang batas, misalnya minimal $t$ dari total $n$ fragmen harus digunakan agar akses dekripsi dapat berhasil.

Pada tahap awal, data dienkripsi menggunakan kunci publik milik \textit{delegator}. Hasil enkripsi pada Umbral terdiri atas \textit{ciphertext} dan \textit{capsule}. \textit{Ciphertext} berisi data terenkripsi, sedangkan \textit{capsule} memuat informasi kriptografis yang diperlukan untuk membuka kunci enkripsi data. Ketika \textit{delegator} ingin memberi akses kepada \textit{delegatee}, \textit{delegator} membuat sejumlah \texttt{kfrag} berdasarkan kunci privat miliknya, kunci publik milik \textit{delegatee}, dan parameter ambang batas yang ditentukan. Fragmen tersebut kemudian diberikan kepada proxy agar proxy dapat melakukan re-enkripsi terhadap \textit{capsule}.

Ketika \textit{delegatee} ingin mengakses data, \textit{delegatee} mengambil \textit{ciphertext} dan \textit{capsule}, lalu meminta proxy menghasilkan \textit{capsule fragment} atau \texttt{cfrag} dari \texttt{kfrag} yang dimiliki proxy. Proxy hanya melakukan transformasi kriptografis terhadap \textit{capsule}; proxy tidak mendekripsi \textit{ciphertext}, tidak mengetahui \textit{plaintext}, dan tidak membutuhkan kunci privat \textit{delegator} maupun \textit{delegatee}. Setelah \textit{delegatee} memperoleh jumlah \texttt{cfrag} valid yang memenuhi ambang batas $t$, \textit{delegatee} dapat membuka \textit{capsule} menggunakan kunci privatnya, memperoleh kunci simetris, lalu mendekripsi \textit{ciphertext}. Dengan demikian, Umbral memungkinkan delegasi hak dekripsi tanpa membagikan kunci privat dan tanpa memberikan akses langsung terhadap isi data kepada proxy.
 \section{Kontrol Akses}
\label{sec:kontrol-akses}

\subsection{Definisi Kontrol Akses}
\label{subsec:definisi-kontrol-akses}
Kontrol Akses adalah proses pengaturan akses ke sumber daya (data, sistem informasi, perangkat, atau area tertentu) \autocite{sandhu1994access, li2022capability}. Proses ini memutuskan apakah seorang subjek harus diberikan atau ditolak aksesnya berdasarkan aturan tertentu. Akses dapat berarti memakai, memasuki, atau menggunakan. Kontrol Akses bertujuan untuk mencegah aktivitas yang dapat menimbulkan pelanggaran keamanan seperti kebocoran data atau modifikasi tidak sah.







\subsection{Kontrol Akses Granular}
\label{subsec:kontrol-akses-granular}

Kontrol akses granular atau \textit{fine-grained access control} adalah pendekatan kontrol akses yang memungkinkan keputusan otorisasi dilakukan pada tingkat cakupan yang lebih spesifik dibandingkan pemberian izin secara umum atau menyeluruh. Jika kontrol akses kasar (\textit{coarse-grained access control}) hanya membedakan akses berdasarkan kategori besar, seperti \textit{role} pengguna atau jenis akses, maka kontrol akses granular mempertimbangkan detail yang lebih banyak, seperti akses berdasarkan pengguna, per baris/kolom data, konteks ketika melakukan akses, waktu berlaku, dan batasan tambahan lain yang relevan dengan kebijakan sistem \autocite{hu2014abac, odugu2021finegrained}.

Sebuah mekanisme kontrol akses dapat disebut granular apabila sistem mampu membedakan hak akses berdasarkan beberapa dimensi otorisasi secara lebih spesifik. Dimensi tersebut meliputi: (1) subjek yang meminta akses, misalnya identitas pengguna, peran, atau atribut; (2) objek atau sumber daya yang diakses, misalnya pengguna tertentu, sumber data tertentu, \textit{dataset} tertentu, atau bagian tertentu dari data; (3) operasi yang dilakukan, misalnya membaca, menulis, memperbarui, mendelegasikan, dan menghapus; dan (4) kondisi atau konteks akses, misalnya waktu berlaku, tujuan akses, lokasi, atau kondisi layanan tertentu \autocite{hu2014abac}.

Granularitas dalam kontrol akses juga berkaitan erat dengan prinsip \textit{least privilege}, yaitu pemberian hak akses sebatas yang diperlukan untuk menjalankan tugas tertentu \autocite{nist80053r5}. Dengan pendekatan ini, pengguna tidak langsung memperoleh akses penuh terhadap seluruh sumber daya hanya karena memiliki peran tertentu. Sebagai contoh, tenaga medis tidak cukup hanya diberi akses berdasarkan peran umum sebagai ``dokter'', melainkan aksesnya dapat dibatasi hanya pada pasien tertentu, episode layanan tertentu, data tertentu, serta operasi tertentu yang sesuai dengan kebutuhan tugasnya. Dengan demikian, kontrol akses granular dapat mengurangi risiko akses berlebihan dan membatasi dampak apabila kredensial atau token akses disalahgunakan.

Pada Tugas Akhir ini, kontrol akses granular dimaknai sebagai kemampuan sistem untuk membatasi hak akses berdasarkan kombinasi antara subjek, data yang diakses, jenis operasi, waktu berlaku, serta batasan delegasi.

\subsection{Model CapBAC}
\label{subsec:model-capbac}

\textit{Capability-Based Access Control} (CapBAC) adalah model kontrol akses yang hak aksesnya direpresentasikan sebagai \textit{"capability token"} yang tidak dapat dipalsukan \autocite{li2022capability, anagnostopoulou2025bolstering}. Token tersebut menyatakan otoritas untuk melakukan operasi tertentu terhadap objek atau sumber daya tertentu. Dalam implementasi tertentu, capability token juga dapat dibatasi dengan identitas subjek, waktu berlaku, cakupan operasi, atau kondisi tambahan lain untuk mempersempit hak akses yang diberikan.

Secara garis besar, model CapBAC bekerja dalam 2 tahapan, (1) Pengguna meminta capability token dari sebuah \textit{Authorization Server}, kemudian (2) Pengguna menggunakan capability token yang diperoleh dari tahap pertama untuk melakukan akses ke \textit{Resource Server} \autocite{li2022capability}.
\textit{Sequence diagram} dari model CapBAC dapat dilihat pada Gambar \ref{fig:sequence-diagram-capbac}.
CapBAC memungkinkan otorisasi terdesentralisasi karena mengandalkan token kapabilitas yang tidak dapat dipalsukan sebagai bukti.

\begin{figure}[ht]
	\centering
	\includegraphics[width=0.8\textwidth]{resources/chapter-2/capbac-seqdiagram.png}
	\caption{\textit{Sequence diagram} dari model CapBAC. (Dokumentasi penulis)}
	\label{fig:sequence-diagram-capbac}
\end{figure}

 \section{Token Otorisasi}
\label{sec:token-otorisasi}

\subsection{JSON Web Token}
\label{subsec:json-web-token}

JSON Web Token (JWT) adalah format token yang ringkas dan \textit{URL-safe} untuk merepresentasikan klaim yang dipertukarkan antara dua pihak \autocite{jones2015jwt}. Klaim pada JWT direpresentasikan dalam bentuk objek JSON dan dapat diamankan menggunakan JSON Web Signature (JWS) atau JSON Web Encryption (JWE). Dengan demikian, JWT dapat digunakan untuk membawa informasi seperti identitas pengguna, peran, cakupan akses, waktu penerbitan, dan waktu kedaluwarsa token.

Secara umum, JWT terdiri dari tiga bagian utama, yaitu \textit{header}, \textit{payload}, dan \textit{signature}. Ketiga bagian tersebut dikodekan menggunakan \texttt{Base64URL} dan dipisahkan dengan tanda titik. Bentuk umum JWT dapat dituliskan sebagai berikut:

\begin{equation*}
	JWT = header.payload.signature
\end{equation*}

Bagian \textit{header} berisi informasi mengenai tipe token dan algoritma kriptografi yang digunakan. Bagian \textit{payload} berisi klaim yang menyatakan informasi mengenai subjek atau hak akses. Bagian \textit{signature} digunakan untuk memastikan integritas token sehingga perubahan pada \textit{header} atau \textit{payload} dapat terdeteksi oleh pihak yang memverifikasi token.

\subsection{PASETO}
\label{subsec:paseto}

PASETO atau \textit{Platform-Agnostic Security Tokens} adalah format token yang dirancang sebagai alternatif terhadap JWT dengan pendekatan kriptografi yang lebih terarah \autocite{arciszewski2022paseto}. Seperti JWT, PASETO digunakan untuk merepresentasikan klaim dalam bentuk token yang ringkas dan dapat dipertukarkan antar pihak. Namun, PASETO membatasi pilihan algoritma kriptografi berdasarkan versi dan tujuan penggunaan token, sehingga mengurangi risiko kesalahan konfigurasi algoritma yang dapat terjadi pada token berbasis JWT.

Secara umum, token PASETO $P$ dapat direpresentasikan sebagai:

\begin{equation*}
	P = version.purpose.message[.footer]
\end{equation*}

$version$ menunjukkan versi PASETO, $purpose$ menunjukkan tujuan penggunaan token , $message$ adalah \textit{payload} yang telah diproses secara kriptografis, dan $footer$ adalah \textit{footer} opsional. Nilai $purpose$ dapat berupa \texttt{local} atau \texttt{public}. Nilai \texttt{local} digunakan ketika isi token perlu dienkripsi menggunakan kunci simetris, sedangkan nilai \texttt{public} digunakan ketika token perlu ditandatangani menggunakan kriptografi kunci publik.

Keunggulan PASETO terletak pada desainnya yang lebih \textit{opinionated} dibandingkan JWT, terutama dalam pemilihan algoritma kriptografi. Dengan membatasi pilihan algoritma, PASETO berupaya mengurangi fleksibilitas yang berpotensi menimbulkan kesalahan implementasi. Namun, dari sisi model otorisasi, PASETO tetap berperan sebagai token pembawa klaim. Artinya, pembatasan akses seperti peran, cakupan, dan waktu berlaku tetap perlu didefinisikan pada saat token diterbitkan.

\subsection{Macaroons}
\label{sec:macaroons}

Macaroon adalah \textit{credential} otorisasi berbasis \textit{bearer} yang mendukung delegasi terdesentralisasi dalam sistem terdistribusi \autocite{birgisson2014macaroons}. Keunggulan utama Macaroon terletak pada penggunaan \textit{chained-HMAC}, yang memungkinkan proses delegasi bersifat terdesentralisasi secara efisien serta mendukung penambahan syarat \textit{caveat} yang sesuai dengan konteks penggunaan.

Sebuah Macaroon $M$ didefinisikan dengan struktur:

\begin{equation*}
	M = (id, \, \mathrm{location}, \, \vec{c}, \, \mathrm{sig})
\end{equation*}

$id$ adalah pengidentifikasi unik, $location$ adalah alamat layanan target, $\vec{c} = (c_1, c_2, \dots, c_n)$ adalah daftar \textit{caveat}, dan $sig$ adalah tanda tangan \autocite{birgisson2014macaroons}. Tanda tangan ($sig$) dihitung secara rekursif untuk menjamin integritas setiap penambahan batasan (\textit{caveat}):

\begin{align*}
	sig_0 & = \text{HMAC}(RootKey, id)    \\
	sig_n & = \text{HMAC}(sig_{n-1}, c_n)
\end{align*}

$RootKey$ adalah kunci rahasia yang digunakan untuk membangkitkan token Macaroon awal. Struktur ini memungkinkan pemegang Macaroon menambahkan \textit{caveat} baru $c_{new}$ (misalnya: \texttt{time < 24h}) dan menghasilkan $sig_{new}$ untuk diberikan ke pihak lain, tanpa mengetahui $RootKey$ yang asli. Kemampuan untuk menambahkan batasan baru pada token secara lokal tanpa melibatkan penerbit awal ini disebut dengan \textit{offline attenuation}. Mekanisme ini memungkinkan delegasi terdesentralisasi yang aman dan fleksibel. Gambaran mengenai delegasi dan atenuasi kewenangan melalui \textit{caveats} dapat dilihat pada Gambar \ref{fig:delegasi-macaroon}.

\begin{figure}[ht]
	\centering
	\includegraphics
	[width=0.8\textwidth]{resources/chapter-2/macaroons.png}
	\caption{Contoh delegasi dan atenuasi kewenangan pada Macaroon. C memegang Macaroon yang diturunkan dari \textit{service} FS untuk mengakses data pada \textit{service} TS dengan batasan waktu, lokasi IP, dan target data yang spesifik. \autocite{birgisson2014macaroons}}
	\label{fig:delegasi-macaroon}
\end{figure}

\subsection{Biscuit Token}
\label{subsec:biscuit-token}

Biscuit adalah token otorisasi yang dirancang untuk mendukung verifikasi terdesentralisasi, atenuasi \textit{offline}, dan penegakan kebijakan akses berbasis \textit{logic language} \autocite{biscuitsec}. Berbeda dari JWT dan PASETO yang berperan sebagai pembawa klaim, Biscuit dirancang sebagai token otorisasi berbasis \textit{capability}. Token ini memungkinkan suatu layanan memverifikasi hak akses menggunakan informasi publik, sehingga cocok untuk lingkungan terdistribusi.

Biscuit menggunakan kriptografi kunci publik untuk memungkinkan verifikasi token oleh pihak yang memiliki kunci publik dari penerbit token. Biscuit mendukung atenuasi hak akses melalui penambahan blok baru pada token. Pemegang token dapat menambahkan batasan baru untuk mempersempit hak akses sebelum token diberikan ke pihak lain.

Karakteristik utama Biscuit adalah penggunaan \textit{logic language} untuk mengekspresikan \textit{rule}. Informasi otorisasi dapat direpresentasikan dalam bentuk \textit{fact}, \textit{rule}, dan \textit{check}. Pada saat token digunakan, pihak \textit{verifier} dapat mengevaluasi token bersama dengan kebijakan lokal untuk menentukan apakah suatu permintaan akses diperbolehkan. Pendekatan ini membuat Biscuit lebih ekspresif dalam memodelkan kebijakan akses yang kompleks, misalnya kebijakan berbasis relasi antar aktor, sumber daya, aksi, dan konteks permintaan. \section{DecMed}
\label{sec:decmed}

Sistem DecMed adalah sebuah kerangka kerja manajemen RME terdesentralisasi yang dibangun di atas DLT IOTA \autocite{sadewa2025decmed}. Pengembangan sistem ini bertujuan untuk mengatasi tantangan fragmentasi data antar institusi kesehatan dan risiko keamanan dari penyimpanan terpusat yang ada di Indonesia.

\subsection{Arsitektur DecMed}
\label{subsec:arsitektur-decmed}

Arsitektur DecMed membagi penyimpanan menjadi dua lapisan, yaitu penyimpanan \textit{on-chain} dan \textit{off-chain}. Penyimpanan \textit{off-chain} memanfaatkan IPFS untuk menyimpan data klinis pasien yang terenkripsi. Sementara itu, penyimpanan \textit{on-chain} memanfaatkan DLT IOTA untuk menyimpan data administratif dan metadata klinis, termasuk CID dari IPFS, guna menjamin integritas data. Untuk memastikan kedaulatan data pasien, DecMed mengimplementasikan CapBAC yang diintegrasikan dengan PRE. Skema PRE yang diterapkan pada DecMed adalah Umbral PRE yang dijelaskan pada Subbab \ref{subsec:umbral-pre}. Komponen utama pada arsitektur DecMed dan fungsinya dijelaskan pada Tabel~\ref{tab:decmed-komponen}. \textit{Client} DecMed terdiri dari tiga jenis. Masing-masing \textit{client} digunakan oleh aktor yang berbeda, seperti pada Tabel~\ref{tab:decmed-client}.

\begin{figure}[ht]
	\centering
	\includegraphics[width=1\textwidth]{resources/chapter-2/decmed_archi.png}
	\caption{Arsitektur DecMed \autocite{sadewa2025decmed}.}
	\label{fig:decmed-archi}
\end{figure}

\begin{table}[!ht]
	\centering
	\caption{Komponen dan fungsi dalam arsitektur DecMed}
	\label{tab:decmed-komponen}
	\begin{tabularx}{\textwidth}{|>{\raggedright\arraybackslash}p{0.28\textwidth}|>{\raggedright\arraybackslash}X|}
		\hline
		\textbf{Komponen}         & \textbf{Fungsi}                                                                                                                 \\
		\hline
		IOTA DLT                  & Menyimpan metadata, referensi data, dan status akses                                                                            \\
		\hline
		IPFS                      & Menyimpan \textit{encrypted blob} data RME                                                                                      \\
		\hline
		PRE                       & Melakukan verifikasi token akses dan \textit{re-encryption} agar penerima akses dapat membuka data tanpa kunci privat pasien    \\
		\hline
		Redis                     & Mengelola \textit{cache ephemeral}, yang terdiri dari data \texttt{Nonce}, \texttt{Re-Encryption Key}, dan \texttt{Revoked Key}
		\\
		\hline
		\textit{Patient Client}   & Mengambil, mendekripsi, dan memberi akses                                                                                       \\
		\hline
		\textit{Hospital Client}  & Mengambil, menenkripsi, dan mendekripsi data RME sesuai hak akses                                                               \\
		\hline
		\textit{Ministry Client}  & Mengelola fasyankes dan adminnya                                                                                                \\
		\hline
		\textit{IOTA Gas Station} & Mensponsori transaksi \textit{on-chain} agar pengguna tidak perlu membayar \textit{gas} secara langsung                         \\
		\hline
	\end{tabularx}
\end{table}

\begin{table}[!ht]
	\centering
	\caption{Aktor, \textit{client}, dan peran dalam DecMed}
	\label{tab:decmed-client}
	\begin{tabularx}{\textwidth}{|>{\raggedright\arraybackslash}p{0.26\textwidth}|>{\raggedright\arraybackslash}p{0.22\textwidth}|>{\raggedright\arraybackslash}X|}
		\hline
		\textbf{Aktor}               & \textbf{\textit{Client}} & \textbf{Peran dalam DecMed}                                                                    \\
		\hline
		Pasien                       & \textit{Patient Client}  & Mengelola identitas, menyimpan atau memperbarui data, dan memberikan akses kepada tenaga medis \\
		\hline
		Tenaga kesehatan / fasyankes & \textit{Hospital Client} & Mengakses atau memperbarui RME sesuai akses yang diberikan pasien                              \\
		\hline
		Pemerintah                   & \textit{Ministry Client} & Mengelola registrasi atau verifikasi fasyankes dalam sistem                                    \\
		\hline
	\end{tabularx}
\end{table}

DecMed memiliki mekanisme untuk tidak membebankan biaya transaksi / \textit{gas} kepada pengguna. Kerangka kerja DecMed memanfaatkan sebuah komponen tambahan bernama IOTA \textit{Gas Station}. Komponen ini bekerja dengan mekanisme \textit{sponsorship} untuk membayar biaya transaksi \textit{on-chain} pada DecMed. Pada Tugas Akhir ini, komponen \textit{Gas Station IOTA} akan diasumsikan sebagai bagian dari IOTA.

\subsection{Alur Pengelolaan dan Akses RME pada DecMed}

Alur pengelolaan RME pada DecMed dimulai ketika data medis dibuat melalui aplikasi \textit{hospital client}. Data tersebut dienkripsi terlebih dahulu di sisi \textit{client} sebelum disimpan pada IPFS. Setelah data terenkripsi berhasil diunggah, IPFS menghasilkan CID sebagai referensi lokasi data. CID, \textit{capsule key}, dan metadata pendukung kemudian dicatat pada DLT IOTA. Dengan demikian, IOTA tidak menyimpan isi rekam medis secara langsung, melainkan menyimpan metadata yang diperlukan untuk menemukan, memverifikasi, dan membuka data terenkripsi melalui mekanisme PRE.

Ketika pasien memberikan akses kepada tenaga medis, PRE menerbitkan token kapabilitas yang digunakan sebagai bukti otorisasi. Tenaga medis membawa token tersebut ketika meminta akses data melalui \textit{Hospital Client}. \textit{Server} PRE kemudian memverifikasi token, melakukan \textit{re-encryption} terhadap \textit{capsule key}, dan mengembalikan material kriptografis yang diperlukan agar \textit{Hospital Client} dapat membuka data terenkripsi sesuai hak akses yang diberikan. \textit{Sequence Diagram} yang menjelaskan alur pemberian akses dapat dilihat pada Lampiran \ref{lampiran:alur-pemberian-akses-medis} \section{Pemodelan Ancaman STRIDE}
\label{sec:stride}

STRIDE merupakan metode \textit{threat modeling} yang dikembangkan oleh Microsoft untuk membantu mengidentifikasi ancaman keamanan pada sistem perangkat lunak sejak tahap perancangan. Metode ini mengelompokkan ancaman ke dalam enam kategori utama, yaitu \textit{Spoofing}, \textit{Tampering}, \textit{Repudiation}, \textit{Information Disclosure}, \textit{Denial of Service}, dan \textit{Elevation of Privilege} \autocite{shostack2014threatmodeling,microsoft2006msdnstride}.

Dalam praktiknya, STRIDE biasanya digunakan bersama \textit{Data Flow Diagram} (DFD) untuk memetakan kemungkinan ancaman terhadap komponen sistem, aliran data, penyimpanan data, dan batas kepercayaan (\textit{trust boundary}) \autocite{shostack2014threatmodeling}.
Pendekatan ini membantu perancang sistem mengidentifikasi kebutuhan kontrol keamanan seperti otentikasi, otorisasi, enkripsi, integritas data, dan \textit{audit logging} sejak tahap desain \autocite{shostack2014threatmodeling}.

Tabel~\ref{tab:stride-components} merangkum enam kategori ancaman pada STRIDE beserta fokus keamanan yang terkait.

\begin{table}[ht]
	\centering
	\caption{Kategori ancaman pada STRIDE}
	\label{tab:stride-components}
	\begin{tabularx}{\textwidth}{|l|l|X|}
		\hline
		\textbf{Kategori}               & \textbf{Fokus Keamanan} & \textbf{Deskripsi Singkat} \\
		\hline
		\textit{Spoofing}               &
		Otentikasi                      &
		Pemalsuan identitas entitas untuk memperoleh akses tidak sah.                          \\
		\hline

		\textit{Tampering}              &
		Integritas                      &
		Modifikasi data, pesan, atau kode secara tidak sah.                                    \\
		\hline

		\textit{Repudiation}            &
		\textit{Non-repudiation}        &
		Penyangkalan terhadap tindakan atau transaksi karena kurangnya bukti audit.            \\
		\hline

		\textit{Information Disclosure} &
		Kerahasiaan                     &
		Kebocoran informasi sensitif kepada pihak yang tidak berhak.                           \\
		\hline

		\textit{Denial of Service}      &
		Ketersediaan                    &
		Gangguan terhadap ketersediaan layanan atau sumber daya sistem.                        \\
		\hline

		\textit{Elevation of Privilege} &
		Otorisasi                       &
		Perolehan hak akses yang lebih tinggi dari yang seharusnya.                            \\
		\hline
	\end{tabularx}
\end{table}

Melalui STRIDE, pengembang dapat mengidentifikasi ancaman keamanan secara sistematis serta menentukan kontrol mitigasi yang sesuai pada tahap perancangan sistem \autocite{shostack2014threatmodeling}. \section{Penelitian Terkait}
\label{sec:penelitian-terkait}

\noindent Bagian ini meninjau beberapa penelitian yang relevan dengan perancangan kontrol akses granular pada sistem RME berbasis DLT dan implementasi sistem menggunakan Macaroons.

\subsection{Hybrid RBAC dan ABAC pada Sistem RME Berbasis Blockchain}
\label{subsec:hybrid-rbac-abac}

Singh dkk.\ mengusulkan model kontrol akses hibrida untuk sistem RME berbasis Blockchain dengan menggabungkan RBAC dan ABAC \autocite{singh2025granular}. Penelitian ini berangkat dari keterbatasan dua model tersebut ketika digunakan secara terpisah. RBAC dinilai terlalu kaku karena keputusan akses hanya bergantung pada peran pengguna. Pendekatan ini menyulitkan penerapan kontrol akses yang \textit{fine-grained} dan berpotensi menimbulkan \textit{role explosion} ketika kebutuhan akses semakin beragam. Di sisi lain, ABAC menawarkan fleksibilitas yang lebih tinggi karena keputusan akses mempertimbangkan atribut pengguna, objek, dan konteks, namun kompleksitas perancangan dan pengelolaan kebijakannya juga meningkat.

Penelitian tersebut menunjukkan bahwa semakin banyak atribut yang digunakan dalam ABAC, semakin besar pula kemungkinan konflik kebijakan dan beban evaluasi aturan \autocite{singh2025granular}. Oleh karena itu, penulis mengusulkan pendekatan \textit{hybrid}, yaitu RBAC digunakan sebagai lapisan awal untuk menyaring akses berdasarkan peran utama, sedangkan ABAC digunakan untuk menambahkan batasan kontekstual yang lebih rinci. Hasil evaluasi pada \textit{smart contract} menunjukkan bahwa pendekatan \textit{hybrid} ini cukup efisien dari sisi konsumsi \textit{gas} dan tetap relatif sederhana untuk diimplementasikan dibandingkan penggunaan ABAC murni.

\subsection{Model Kontrol Akses Empat Lapis untuk RME Berbasis Blockchain}
\label{subsec:empat-lapis-emr}

Liu dan Chen meneliti model kontrol akses untuk RME berbasis Blockchain dengan menempatkan RME sebagai kasus \textit{big data} \autocite{liu2025accesscontrol}. Peran DLT dalam mengatasi masalah tersebut relevan untuk menjelaskan mengapa pendekatan DLT layak dipertimbangkan sebagai alternatif terhadap sistem RME terpusat, karena DLT menawarkan integritas data yang lebih kuat, jejak audit yang transparan, dan distribusi kendali yang lebih baik.

Untuk menjawab permasalahan tersebut, penelitian ini mengusulkan model otorisasi empat lapis yang terdiri dari:
\begin{enumerate}
	\item \textit{Role Check} menggunakan RBAC,
	\item \textit{Identity Verification} menggunakan \textit{Identity-Based Encryption} (IBE),
	\item \textit{Policy Execution} menggunakan \textit{smart contract}, dan
	\item \textit{Data Volume Check} menggunakan pendekatan \textit{entropy}.
\end{enumerate}





\subsection{Kontrol Akses Granular Menggunakan \textit{Composable NFT}}
\label{subsec:composable-nft-rme}

Amaracitra mengusulkan model \textit{Decentralized Medical Record} (DMR) yang menggunakan \textit{composable NFT} untuk mencapai keamanan data, interoperabilitas, dan privasi melalui akses data yang granular \autocite{amaracitra2025composable}. Nilai utama dari penelitian ini terletak pada cara data RME dipecah menjadi unit-unit yang lebih kecil, sehingga tidak seluruh data pasien harus selalu dibuka dalam satu paket akses.

Dalam model tersebut, struktur data DMR disusun secara hierarkis ke dalam beberapa token, yaitu:
\begin{enumerate}
	\item \textit{DMR Token} sebagai token induk (\textit{root}/\textit{parent}),
	\item \textit{Medical Record Token},
	\item \textit{Doctor Diagnoses Token},
	\item \textit{Laboratory Test Token}, dan
	\item \textit{Pharmacy Data Token}.
\end{enumerate}

Ilustrasi struktur data dari implementasi DMR menggunakan composable NFT ini dapat dilihat pada Gambar \ref{fig:struktur-data-dmr}.

\begin{figure}[ht]
	\centering
	\includegraphics[width=0.8\textwidth]{resources/chapter-2/NFT hierarchical structure.png}
	\caption{Ilustrasi struktur data DMR menggunakan composable NFT. \autocite{amaracitra2025composable}}
	\label{fig:struktur-data-dmr}
\end{figure}

Pemisahan ini penting karena menunjukkan bahwa granularitas akses tidak hanya bergantung pada mekanisme otorisasi, tetapi juga pada cara data medis dimodelkan sejak awal \autocite{amaracitra2025composable}. Dengan struktur seperti ini, akses dapat diberikan hanya pada bagian data yang relevan dengan kebutuhan tenaga kesehatan tertentu. Misalnya, tenaga laboratorium dapat berinteraksi dengan data pemeriksaan laboratorium tanpa harus memperoleh akses ke keseluruhan rekam medis pasien.



\subsection{Implementasi Macaroons untuk Delegasi Akses pada Sistem Nyata}
\label{subsec:macaroons-delegasi-otorisasi}

Selain diperkenalkan secara teoretis sebagai mekanisme otorisasi berbasis \textit{caveats}, Macaroons juga telah digunakan pada beberapa sistem penyimpanan terdistribusi. Salah satu contohnya adalah Storj, yaitu platform penyimpanan objek terdesentralisasi yang menggunakan Macaroons sebagai dasar pengelolaan akses pada jaringan Tardigrade \autocite{storj_macaroons}. Dalam implementasi tersebut, Macaroons diposisikan sebagai \textit{API key} yang berisi aturan akses yang melekat pada token tersebut. Aturan tersebut dapat membatasi operasi yang diperbolehkan, seperti \textit{read}, \textit{write}, \textit{delete}, dan \textit{list}, serta membatasi ruang lingkup akses berdasarkan \textit{bucket}, \textit{path}, \textit{path prefix}, dan jendela waktu akses. Dengan demikian, pengguna dapat membagikan akses terhadap objek data tertentu tanpa memberikan akses penuh terhadap seluruh akun atau seluruh ruang penyimpanan.

Penerapan Macaroons pada Storj menunjukkan bahwa token otorisasi tersebut dapat digunakan pada lingkungan penyimpanan terdesentralisasi. Otoritas akses dibawa oleh token yang secara eksplisit memuat cakupan hak aksesnya. Token tersebut kemudian dapat diturunkan menjadi token dengan cakupan yang lebih sempit melalui penambahan \textit{caveats}.

Contoh implementasi lain dapat ditemukan pada dCache, yaitu sistem penyimpanan data ilmiah berskala besar yang mendukung Macaroons sebagai \textit{bearer token} \autocite{dcache_macaroons}. Dalam dCache, Macaroons dirancang untuk mempermudah proses berbagi akses terhadap file atau direktori tertentu. Pengguna yang memiliki akses terhadap suatu \textit{resource} dapat meminta Macaroon, kemudian token tersebut dapat digunakan oleh pihak lain untuk mengakses sumber daya yang sama sesuai batasan yang tertanam di dalam token. dCache mendefinisikan beberapa \textit{caveats} yang merepresentasikan batasan akses, seperti \texttt{activity} untuk menentukan aktivitas yang diizinkan, \texttt{before} untuk membatasi masa berlaku token, \texttt{ip} untuk membatasi alamat jaringan pengguna, serta \texttt{path} dan \texttt{root} untuk membatasi ruang lingkup akses pada struktur direktori.

Implementasi dCache juga menunjukkan bahwa beberapa \textit{caveats} dapat dikombinasikan untuk menghasilkan hak akses yang lebih terbatas. Sebagai contoh, beberapa \texttt{activity} caveat digabungkan dengan mengambil irisan dari aktivitas yang diizinkan, sehingga penambahan \textit{caveat} baru tidak memperluas hak akses token, melainkan mempersempitnya. Selain itu, pengguna yang telah memiliki Macaroon dapat membuat Macaroon baru dengan tambahan \textit{caveats}, baik secara mandiri menggunakan pustaka Macaroons maupun melalui permintaan ke dCache. Mekanisme ini memperlihatkan pola delegasi akses yang sesuai dengan prinsip atenuasi, yaitu token turunan tidak dapat memiliki hak akses yang lebih besar daripada token asal.

Dari kedua implementasi tersebut, dapat disimpulkan bahwa Macaroons telah digunakan untuk membatasi dan mendelegasikan akses pada sistem penyimpanan data terdistribusi. Storj menekankan penggunaan Macaroons untuk pengelolaan akses objek data pada jaringan penyimpanan terdesentralisasi, sedangkan dCache menunjukkan bagaimana \textit{caveats} dapat digunakan untuk membatasi aktivitas, lokasi sumber daya, alamat jaringan, dan masa berlaku akses.

  \chapter{Analisis dan Rancangan Solusi}
\label{chap:analisis-dan-rancangan}

\section{Analisis Masalah}
\label{sec:analisis-masalah}

Dalam mengembangkan mekanisme kontrol akses granular dan delegasi pada kerangka kerja DecMed, diperlukan analisis terhadap implementasi dan arsitektur sistem yang sudah ada. Hasil analisis tersebut digunakan sebagai dasar untuk mengidentifikasi permasalahan serta merumuskan kebutuhan sistem yang diusulkan pada Tugas Akhir ini.

Analisis masalah pada Tugas Akhir ini diawali dengan analisis terhadap implementasi dan arsitektur DecMed pada Subbab \ref{subsec:analisis-sistem-decmed}, termasuk analisis STRIDE dan skenario layanan kesehatan kolaboratif. Hasil analisis tersebut kemudian digunakan untuk mengidentifikasi permasalahan pada Subbab \ref{subsec:identifikasi-masalah} dan merumuskan kebutuhan sistem pada Subbab \ref{subsec:kebutuhan-sistem}.

\subsection{Analisis Sistem DecMed}
\label{subsec:analisis-sistem-decmed}

Sebagaimana dijelaskan pada Bab \ref{sec:decmed}, DecMed adalah kerangka kerja manajemen RME terdesentralisasi dengan menggunakan DLT IOTA sebagai penyimpanan \textit{on-chain}, IPFS sebagai penyimpanan \textit{off-chain}, serta PRE untuk mendukung distribusi akses data yang aman. Arsitektur ini sudah memberikan fondasi untuk menjaga integritas data dan mengurangi ketergantungan pada penyimpanan terpusat. Terdapat tiga aktor utama pada DecMed, yaitu:

\begin{enumerate}
	\item Pasien
	\item Staf fasyankes, terdiri dari tenaga kesehatan, petugas rekam medis, dan admin fasyankes
	\item Pemerintah, yaitu Kementerian Kesehatan
\end{enumerate}

Pada DecMed, masing-masing aktor memiliki kepentingan dan aplikasi \textit{client} yang berbeda. \textit{Patient Client} digunakan oleh pasien, \textit{Hospital Client} digunakan oleh tenaga kesehatan di sebuah fasyankes, dan \textit{Ministry Client} digunakan oleh pemerintah. Detail kebutuhan fungsional DecMed untuk memenuhi kebutuhan masing-masing aktor dapat dilihat pada Lampiran \ref{lampiran:kebutuhan-fungsional-decmed}.

\subsubsection{Keterbatasan Token Kapabilitas DecMed}
\label{subsubsec:keterbatasan-token-kapabilitas-decmed}

Model kontrol akses yang digunakan pada sistem DecMed adalah model CapBAC. DecMed memanfaatkan token otorisasi berupa JWT sebagai token otorisasi yang merepresentasikan kapabilitas akses. Pada kerangka kerja DecMed, ketika pasien memberikan akses terhadap RME mereka kepada seorang tenaga medis, maka PRE akan menerbitkan token JWT sesuai dengan \textit{request} dari pasien. Token yang diterbitkan oleh pasien membawa klaim seperti pada Kode \ref{lst:decmed-jwt-claims}. Dalam implementasi DecMed sebelumnya, token JWT masih bersifat \textit{coarse-grained}. Penentuan hak akses pengguna dilakukan hanya berdasarkan \textit{role} dan \textit{purpose} yang ada pada token yang disertakan saat melakukan \textit{request} ke PRE. Token tersebut belum dirancang untuk membawa batasan kontekstual, misalnya jenis data yang boleh diakses. Padahal, pada praktik pelayanan kesehatan, setiap tenaga medis tidak selalu membutuhkan keseluruhan isi RME pasien sesuai dengan prinsip \textit{need-to-know}. Kondisi ini menunjukkan bahwa desain \textit{claims} token otorisasi pada DecMed belum cukup untuk menjawab kebutuhan kontrol akses granular.

\begin{lstlisting}[language=Rust, caption={Struktur klaim JWT pada DecMed}, label={lst:decmed-jwt-claims}]
#[derive(Debug, Deserialize, Serialize)]
pub struct JwtClaims {
    pub address: String, // address IOTA pemegang token
    pub role: AuthRole, // AdministrativePersonnel 
    // atau MedicalPersonnel
    pub purpose: ReencryptionPurposeType, // Read atau Update
}
\end{lstlisting}

Rancangan \textit{claims} pada JWT ini menimbulkan keterbatasan karena token JWT pada DecMed masih bersifat \textit{coarse-grained}. Penentuan hak akses pengguna dilakukan hanya berdasarkan \textit{role} dan \textit{purpose} yang ada pada token yang disertakan saat melakukan \textit{request} ke PRE. Token tersebut belum dirancang untuk membawa batasan kontekstual, misalnya jenis data yang boleh diakses. Padahal, pada praktik pelayanan kesehatan, setiap tenaga medis tidak selalu membutuhkan keseluruhan isi RME pasien sesuai dengan prinsip \textit{need-to-know}. Kondisi ini menunjukkan bahwa desain \textit{claims} token otorisasi pada DecMed belum cukup untuk menjawab kebutuhan kontrol akses granular.



\subsubsection{Keterbatasan Segmentasi Data pada DecMed}
\label{subsubsec:keterbatasan-segmentasi-data-decmed}

Pada kerangka kerja DecMed, data RME pasien disimpan dalam dua sistem penyimpanan yang berbeda, yaitu penyimpanan \textit{on-chain} (IOTA) dan \textit{off-chain} (IPFS). Penyimpanan \textit{on-chain} digunakan untuk menyimpan metadata dan CID dari data yang disimpan di \textit{off-chain}. Penyimpanan \textit{off-chain} digunakan untuk menyimpan data administratif dan klinis yang terenkripsi. Pemisahan ini sesuai dengan kebutuhan yang dibahas pada \ref{subsec:isi-rme} terkait dokumentasi administratif dan klinis.

Pada penyimpanan \textit{on-chain}, data RME pasien direpresentasikan sebagai metadata dalam sebuah \textit{struct} \texttt{MedicalMetadata} untuk data klinis. Metadata tersebut menyimpan CID serta informasi kunci terenkripsi yang digunakan dalam proses dekripsi data medis terkait. Struktur penyimpanannya terikat dengan \textit{struct} \texttt{PatientAccount} dan \texttt{PatientMedicalMetadata} pada Kode \ref{lst:decmed-patient-medical-metadata}.

\begin{lstlisting}[language=Move, caption={Struktur metadata medis pasien pada DecMed}, label={lst:decmed-patient-medical-metadata}]
public struct PatientAccount has store {
    ...
    medical_metadata: TableVec<PatientMedicalMetadata>,
}
public struct PatientMedicalMetadata has copy, drop, store {
    index: u64,
    metadata: String,
}
\end{lstlisting}

Karena \textit{field} \texttt{medical\_metadata} pada \textit{struct} \texttt{PatientAccount} menyimpan daftar metadata klinis pasien, akses terhadap struktur ini memungkinkan sistem memperoleh referensi menuju objek data klinis terenkripsi, seperti CID dan \textit{capsule} untuk PRE. CID dan \textit{capsule} tidak secara langsung membuka isi RME karena data klinis tetap disimpan dalam bentuk terenkripsi. Namun, struktur metadata tersebut menunjukkan bahwa unit akses pada DecMed masih mengarah pada satu objek data klinis yang relatif besar, bukan pada segmen data yang spesifik.

Kondisi tersebut menyulitkan penegakan kontrol akses granular. Ketika proses \textit{re-encryption} diberikan terhadap objek data klinis tersebut, penerima akses berpotensi memperoleh cakupan \textit{plaintext} yang lebih luas daripada kebutuhan aktualnya. Hal ini terjadi karena struktur metadata belum menyediakan atribut segmentasi yang cukup untuk membedakan cakupan akses sesuai kebutuhan pelayanan. Variabel \texttt{metadata} pada \textit{struct} \texttt{PatientMedicalMetadata} merupakan serialisasi dari \textit{struct} yang ditunjukkan pada Kode \ref{lst:decmed-medical-metadata}.

\begin{lstlisting}[language=Rust, caption={Struktur serialisasi metadata medis DecMed}, label={lst:decmed-medical-metadata}]
#[derive(Debug, Deserialize, Serialize)]
pub struct MedicalMetadata {
    pub capsule: String, // capsule untuk PRE re-encryption
    pub cid: String, // pointer ke data di IPFS
    pub created_at: String, // timestamp
    pub enc_key_and_nonce: String, // kunci dan nonce yang 
    // terenkripsi
}
\end{lstlisting}

\texttt{cid} pada \textit{struct} \texttt{MedicalMetadata} merupakan \textit{pointer} ke data klinis yang sudah terenkripsi disimpan di penyimpanan \textit{off-chain} sebagai sebuah \textit{Encrypted Blob}. Potongan field relevan pada \textit{struct} \texttt{MedicalData} dapat dilihat pada Kode \ref{lst:decmed-medical-data}.

\begin{lstlisting}[language=Rust, caption={Potongan struktur data medis \textit{off-chain} pada DecMed}, label={lst:decmed-medical-data}]
pub struct MedicalData {
    pub anamnesis: String,
    pub physical_check: String,
    pub psychological_check: String,
    pub diagnose: String,
    pub therapy: String,
}
\end{lstlisting}

Penyimpanan data klinis pada DecMed hanya berfokus pada dataset rawat jalan. Pada \textit{struct} \texttt{MedicalData}, data klinis masih disimpan dengan format yang sederhana, tanpa kategorisasi yang jelas. DecMed memang telah memisahkan penyimpanan \textit{on-chain} dan \textit{off-chain}, tetapi pendekatan tersebut belum secara eksplisit menunjukkan segmentasi data klinis yang cukup rinci untuk menegakkan kebijakan akses granular. Jika data klinis masih dikelola sebagai objek terenkripsi yang terlalu besar, pemberian akses kepada tenaga medis berisiko membuka data di luar kebutuhan klinisnya. Sebagai contoh, seorang perawat tidak membutuhkan akses tulis ke data \texttt{therapy} pada \textit{struct} \texttt{MedicalData} karena yang berhak menentukan terapi adalah dokter.

\subsubsection{Analisis Skenario Layanan Kesehatan Kolaboratif}
\label{subsubsec:analisis-skenario-layanan-klinis}
Kebutuhan akan kontrol akses granular dan delegasi berantai menjadi lebih jelas ketika arsitektur DecMed ditempatkan pada skenario layanan kesehatan yang nyata. Dalam alur pelayanan rumah sakit, pasien tidak hanya berinteraksi dengan satu pihak, tetapi dengan beberapa aktor yang memiliki kebutuhan akses berbeda pada tahap yang berbeda pula sebagaimana dibahas pada Subbab \ref{subsec:rme-layanan-kolaboratif}. Oleh karena itu, subbab ini memodelkan satu skenario layanan kesehatan kolaboratif untuk menunjukkan secara konkret mengapa mekanisme akses pada DecMed saat ini belum memadai.

Skenario yang dianalisis adalah pelayanan pasien rawat jalan dan rawat inap yang melibatkan petugas rekam medis, perawat, dokter penanggung jawab pasien, petugas laboratorium, dan apoteker. Alur tersebut dipilih karena mewakili pola kolaborasi yang umum terjadi di sebuah fasyankes.

\begin{enumerate}[]
	\item \textbf{Tahap pendaftaran pasien.} Saat pasien datang ke rumah sakit, petugas rekam medis hanya membutuhkan akses terhadap data administratif pasien.
	\item \textbf{Tahap pemeriksaan awal oleh perawat.} Setelah proses pendaftaran selesai, perawat memerlukan hak untuk menulis data klinis untuk anamnesis awal, pemeriksaan fisik, dan pemeriksaan psikologis.
	\item \textbf{Tahap pemeriksaan oleh dokter.} Pada tahap ini, dokter penanggung jawab pasien memerlukan hak baca terhadap keseluruhan data RME pasien dan hak tulis untuk pelayanan saat ini.
	\item \textbf{Tahap rujukan internal ke laboratorium.} Apabila dokter membutuhkan pemeriksaan laboratorium, maka pasien perlu diarahkan ke petugas laboratorium. Dalam tahap ini, petugas laboratorium hanya membutuhkan akses terhadap data administratif pasien dan permintaan pemeriksaan laboratorium.
	\item \textbf{Tahap tindak lanjut hasil pemeriksaan.} Setelah hasil laboratorium tersedia, dokter kembali menggunakan data tersebut sebagai bahan pengambilan keputusan klinis.
	\item \textbf{Tahap pemberian resep oleh apotek.} Setelah dokter memberikan resep, apotek memerlukan hak untuk membaca data resep dan mengunggah data obat yang diresepkan.
\end{enumerate}

Berdasarkan skenario tersebut, kebutuhan akses setiap aktor dapat diringkas pada Tabel \ref{tab:skenario-layanan-kolaboratif}.

\begin{table}[!ht]
	\centering
	\caption{Pemetaan Aktor dan Kebutuhan Akses pada Skenario Layanan Kesehatan Kolaboratif}
	\label{tab:skenario-layanan-kolaboratif}
	\small
	\begin{tabular}{|p{0.2\textwidth}|p{0.23\textwidth}|p{0.45\textwidth}|}
		\hline
		\textbf{Aktor}       & \textbf{Jenis Akses}                                & \textbf{Cakupan Data yang Dibutuhkan}                                                                     \\ \hline
		Petugas Rekam Medis  & Baca data administratif                             & Data administratif pasien.                                                                                \\ \hline
		Perawat              & Baca dan tulis data klinis (hanya pemeriksaan awal) & Baca data administratif pasien. Tulis data anamnesis awal, pemeriksaan fisik, dan pemeriksaan psikologis. \\ \hline
		Dokter               & Baca dan tulis data klinis                          & Baca keseluruhan RME pasien. Tulis data sesuai konteks pelayanan saat ini.                                \\ \hline
		Petugas Laboratorium & Baca dan tulis terbatas (hanya data laboratorium)   & Baca data administratif pasien dan permintaan pemeriksaan penunjang. Tulis data laboratorium.             \\ \hline
		Apoteker             & Baca dan tulis terbatas (hanya data apotek)         & Baca data administratif pasien dan data terapi. Tulis data peresepan dan \textit{dispensing}.             \\ \hline
	\end{tabular}
\end{table}

Tabel \ref{tab:skenario-layanan-kolaboratif} juga menunjukkan bahwa kategori tenaga medis tidak dapat diperlakukan sebagai satu kelompok akses tenaga medis. Dokter, perawat, petugas laboratorium, dan apoteker sama-sama terlibat dalam pelayanan klinis, tetapi memiliki kewenangan baca dan tulis yang berbeda. Oleh karena itu, penggunaan role umum seperti \texttt{MedicalPersonnel} yang digunakan pada DecMed saat ini belum cukup untuk menentukan hak akses tulis secara aman, karena sistem tetap perlu membedakan fungsi klinis masing-masing aktor.

Alur pemberian akses terhadap data RME pasien oleh pasien kepada tenaga medis pada DecMed sebelumnya dapat dilihat pada Lampiran \ref{lampiran:alur-pemberian-akses-medis}. Pada alur tersebut, pasien harus memilih satu tenaga medis lalu menerbitkan akses untuk tenaga medis tersebut. Implementasi DecMed sebelumnya tidak memiliki mekanisme delegasi antar tenaga kesehatan. Sehingga ketika akses perlu diberikan kepada lebih dari satu tenaga medis, proses yang sama perlu diulang dari awal untuk masing-masing penerima akses.

Dalam skenario tersebut, keterbatasan DecMed tampak secara operasional. Model token kapabilitas yang statis menyulitkan proses pembatasan akses secara bertahap sesuai perpindahan peran di lapangan. Selain itu, pemberian akses hanya dapat dilakukan secara individual dari pasien kepada satu tenaga medis pada satu waktu, sehingga proses otorisasi menjadi tidak efektif ketika layanan melibatkan beberapa tenaga medis sekaligus. Pasien perlu memberikan akses kepada dokter utama dan petugas laboratorium secara terpisah, padahal proses tersebut dapat dibuat lebih efisien apabila dokter utama dapat mendelegasikan sebagian hak aksesnya kepada petugas laboratorium tanpa perlu intervensi langsung dari pasien. Dokter perlu dapat mendelegasikan sebagian hak akses saja, misalnya akses terhadap permintaan pemeriksaan dan hak tulis hasil laboratorium tanpa ikut membuka seluruh riwayat klinis pasien.

Dengan demikian, skenario layanan kesehatan kolaboratif ini memperlihatkan bahwa tiga persoalan penelitian saling berkaitan: sistem memerlukan kontrol akses yang granular agar setiap aktor hanya memperoleh data yang relevan, segmentasi data \textit{off-chain} agar pembatasan akses dapat ditegakkan secara teknis, dan delegasi berantai dengan atenuasi \textit{offline} agar perpindahan hak akses dapat berlangsung efisien tanpa intervensi ulang dari pasien.
 \subsubsection{Analisis STRIDE pada Sistem DecMed}
\label{subsec:analisis-stride-decmed}

Analisis STRIDE pada subbab ini digunakan untuk mengevaluasi implikasi keamanan dari arsitektur DecMed saat ini ketika ditempatkan pada skenario layanan kesehatan kolaboratif yang menjadi fokus Tugas Akhir ini. Berdasarkan pembahasan pada Subbab \ref{sec:stride}, tahapan awal dalam melakukan pemodelan ancaman menggunakan STRIDE adalah mengidentifikasi komponen dan alur kerja sistem yang menjadi target ancaman. Hal ini dapat dilakukan dengan menggambarkan DFD yang menunjukkan aliran data dan komunikasi antar komponen sistem. Pada analisis ini, dilakukan pemodelan DFD Level 0 dan Level 1. Pemodelan ke level lebih lanjut tidak dilakukan karena pemodelan hingga level 1 sudah cukup untuk mengidentifikasi komponen dari DecMed dan alur data antar komponen. Hasil pemodelan DFD dapat dilihat pada Lampiran \ref{lampiran:dfd-decmed}.

Berdasarkan arsitektur DecMed pada Gambar \ref{fig:decmed-archi} dan DFD pada Lampiran \ref{lampiran:dfd-decmed}, \textit{attack surface} utama sistem mencakup \textit{Patient Client}, \textit{Hospital Client}, \textit{Ministry Client}, data pada lapisan \textit{on-chain}, data terenkripsi pada lapisan \textit{off-chain}, Redis PRE, jaringan IOTA, IOTA \textit{Gas Station}, serta token kapabilitas yang mengatur otorisasi akses. Hasil identifikasi ancaman berdasarkan metode STRIDE terhadap komponen utama pada DecMed dapat dilihat pada Lampiran \ref{lampiran:analisis-stride-decmed}.

Berdasarkan hasil analisis STRIDE pada Tabel \ref{tab:stride-decmed}, ancaman yang paling relevan terhadap fokus Tugas Akhir ini adalah \textit{spoofing} (T01), \textit{tampering} (T02 dan T03), \textit{information disclosure} (T05 dan T06), dan \textit{elevation of privilege} (T08). Ancaman-ancaman tersebut secara langsung berkaitan dengan kebutuhan akan kontrol akses yang lebih granular, segmentasi data \textit{off-chain}, serta mekanisme delegasi yang mampu membatasi hak akses secara bertahap. Ancaman \textit{repudiation} (T04) tetap relevan terhadap kebutuhan audit delegasi, tetapi pada penelitian ini diposisikan sebagai implikasi pendukung karena fokus utama berada pada pembatasan cakupan akses dan delegasi.

Dengan demikian, hasil analisis STRIDE memperkuat bahwa keterbatasan pada model kontrol akses dan delegasi DecMed tidak hanya merupakan persoalan fungsional, tetapi juga memiliki implikasi keamanan yang nyata. Oleh karena itu, diperlukan mekanisme otorisasi yang mampu mendukung pembatasan akses yang lebih rinci dan aman pada sistem rekam medis elektronik terdesentralisasi. \subsubsection{Analisis Mekanisme Penggunaan dan Revokasi Token}
\label{subsubsec:analisis-mekanisme-penggunaan-dan-revokasi-akses}

Pada implementasi DecMed sebelumnya, kapabilitas akses ditentukan oleh kepemilikan token dan klaim yang dibawa oleh token tersebut. Setiap operasi pembacaan maupun penulisan data RME dilakukan melalui \textit{request} yang menyertakan token akses. Komponen PRE kemudian memvalidasi token dengan mengacu pada hak akses yang tercatat pada penyimpanan \textit{on-chain} IOTA. Validasi ini digunakan untuk memastikan bahwa \textit{address} IOTA yang tercantum pada klaim token memiliki hak akses yang sesuai pada penyimpanan \textit{on-chain} dan hak akses tersebut masih berlaku. Struktur klaim token ditunjukkan pada Kode \ref{lst:decmed-jwt-claims}.

Mekanisme tersebut masih memiliki keterbatasan karena validasi berfokus pada klaim yang dibawa token, bukan pada pembuktian bahwa pengirim \textit{request} benar-benar merupakan pemilik \textit{address} yang tercantum pada klaim. Jika token akses diperoleh oleh pihak lain, misalnya melalui kebocoran atau pencurian token, pihak tersebut dapat mengirimkan \textit{request} dengan klaim \textit{address} yang masih valid. Kondisi ini membuat token dapat digunakan oleh aktor yang tidak berwenang selama klaim di dalam token masih diterima oleh sistem. Oleh karena itu, diperlukan mekanisme tambahan untuk mengikat penggunaan token dengan bukti identitas dari aktor yang melakukan akses.

Selain verifikasi terhadap pengguna token, mekanisme pencabutan hak akses juga perlu diperhatikan. Pada DecMed, pasien dapat mencabut hak akses sebelum masa berlaku akses berakhir. Proses ini dilakukan melalui \textit{patient client} dengan mengirimkan \textit{request} revokasi ke PRE. Setelah diterima, status revokasi dicatat pada Redis dan penyimpanan \textit{on-chain} IOTA agar akses yang telah dicabut tidak lagi diterima oleh sistem.

Apabila mekanisme delegasi akses ditambahkan, sebagaimana dibutuhkan pada skenario layanan kesehatan kolaboratif di fasilitas pelayanan kesehatan pada Subbab \ref{subsubsec:analisis-skenario-layanan-klinis}, mekanisme revokasi tidak cukup hanya menangani token awal. Token hasil delegasi juga perlu diperhitungkan, termasuk hubungan antara token induk dan token turunannya. Dengan demikian, pencabutan akses dapat diterapkan tidak hanya terhadap akses yang diberikan langsung oleh pasien, tetapi juga terhadap akses yang muncul melalui proses delegasi.
  \subsection{Identifikasi Masalah}
\label{subsec:identifikasi-masalah}

Berdasarkan analisis terhadap arsitektur DecMed pada Subbab \ref{subsec:analisis-sistem-decmed}, Tugas Akhir ini mengidentifikasi sejumlah masalah yang perlu diselesaikan agar DecMed dapat mendukung kontrol akses granular dan delegasi \textit{offline}. Permasalahan utama terdiri atas tiga aspek. Pertama, token otorisasi DecMed belum mampu menyatakan batasan akses secara granular, seperti jenis data, aksi akses, masa berlaku, dan ruang lingkup delegasi. Kedua, data RME pada penyimpanan \textit{off-chain} belum direpresentasikan sebagai objek akses yang tersegmentasi sehingga pembatasan akses sulit ditegakkan pada unit data spesifik. Ketiga, alur pemberian akses DecMed masih bergantung pada keterlibatan langsung pasien sehingga belum mendukung delegasi antaraktor dalam proses layanan kesehatan yang kolaboratif.

Selain tiga masalah utama tersebut, terdapat beberapa masalah turunan yang perlu diperhatikan agar rancangan solusi dapat ditegakkan secara aman. Delegasi antaraktor harus dibatasi agar hak akses turunan tidak melebihi hak akses induknya. Perbedaan data klinis yang dapat diakses oleh masing-masing aktor juga perlu diperhitungkan karena \textit{role} umum seperti \texttt{MedicalPersonnel} belum cukup untuk membedakan kewenangan dokter, perawat, petugas laboratorium, dan apoteker. Selain itu, karena delegasi dapat membentuk rantai akses, mekanisme pencabutan dan audit pada DecMed perlu diperluas agar dapat mencakup akses turunan. Penggunaan token akses juga perlu dikaitkan dengan pembuktian identitas aktor aktif pada setiap permintaan agar token tidak hanya diperlakukan sebagai \textit{bearer token} yang dapat digunakan oleh pihak yang tidak semestinya.

Agar keterhubungan antara analisis masalah, ancaman keamanan, kebutuhan sistem, rancangan, implementasi, dan pengujian dapat ditelusuri, permasalahan tersebut dirinci menggunakan ID \texttt{M-TA} pada Tabel \ref{tab:identifikasi-masalah-decmed}.

\begin{table}[!ht]
	\centering
	\caption{Identifikasi Masalah DecMed}
	\label{tab:identifikasi-masalah-decmed}
	\small
	\begin{tabular}{|p{0.12\textwidth}|p{0.60\textwidth}|p{0.15\textwidth}|}
		\hline
		\textbf{ID}      & \textbf{Masalah}                                                                                         & \textbf{Dasar Analisis}                                                    \\ \hline
		\texttt{M-TA-01} & Token otorisasi DecMed belum mampu mengekspresikan batasan akses secara granular.                        & \ref{subsubsec:keterbatasan-token-kapabilitas-decmed}, T02, T05            \\ \hline
		\texttt{M-TA-02} & Data RME belum direpresentasikan sebagai objek akses yang tersegmentasi.                                 & \ref{subsubsec:keterbatasan-segmentasi-data-decmed}, T05                   \\ \hline
		\texttt{M-TA-03} & Alur pemberian akses belum mendukung delegasi antaraktor dalam layanan kolaboratif.                      & \ref{subsubsec:analisis-skenario-layanan-klinis}                           \\ \hline
		\texttt{M-TA-04} & Delegasi antaraktor memerlukan mekanisme atenuasi.                                                       & \ref{subsubsec:analisis-skenario-layanan-klinis},     T05, T08             \\ \hline
		\texttt{M-TA-05} & \textit{Role} \texttt{MedicalPersonnel} terlalu umum untuk membedakan kewenangan antar aktor.            & \ref{subsubsec:analisis-skenario-layanan-klinis}, T02, T05                 \\ \hline
		\texttt{M-TA-06} & Mekanisme revokasi dan audit akses perlu diperluas agar mencakup akses turunan pada rantai delegasi.     & \ref{subsubsec:analisis-mekanisme-penggunaan-dan-revokasi-akses}, T04, T08 \\ \hline
		\texttt{M-TA-07} & Token akses belum secara eksplisit membuktikan kepemilikan identitas aktor aktif pada setiap permintaan. & \ref{subsubsec:analisis-mekanisme-penggunaan-dan-revokasi-akses}, T01      \\ \hline
	\end{tabular}
\end{table}

Berdasarkan Tabel \ref{tab:identifikasi-masalah-decmed}, Tugas Akhir ini tidak hanya memandang keterbatasan DecMed sebagai persoalan token otorisasi, tetapi sebagai rangkaian masalah yang saling bergantung. \texttt{M-TA-01}, \texttt{M-TA-02}, dan \texttt{M-TA-03} merupakan masalah inti yang masing-masing berkaitan dengan kebutuhan kontrol akses granular, segmentasi objek data RME, dan delegasi antaraktor. Sementara itu, \texttt{M-TA-04}, \texttt{M-TA-05}, \texttt{M-TA-06}, dan \texttt{M-TA-07} merupakan masalah turunan atau implikasi yang perlu diselesaikan agar ketiga masalah inti tersebut dapat ditangani secara aman. \texttt{M-TA-04} menekankan perlunya pembatasan hak akses turunan, \texttt{M-TA-05} menekankan perlunya pembedaan kewenangan berdasarkan aktor, \texttt{M-TA-06} menekankan perlunya perluasan pencabutan dan audit terhadap rantai delegasi, sedangkan \texttt{M-TA-07} menekankan perlunya pembuktian identitas aktor aktif pada setiap permintaan akses. ID masalah tersebut kemudian digunakan sebagai dasar pemetaan kebutuhan, rancangan, implementasi, dan pengujian pada bagian selanjutnya.

\subsection{Kebutuhan Sistem}
\label{subsec:kebutuhan-sistem}

Berdasarkan masalah yang telah dipetakan pada Tabel \ref{tab:identifikasi-masalah-decmed}, kebutuhan sistem pada Tugas Akhir ini dirumuskan sebagai perluasan terhadap kebutuhan DecMed sebelumnya. Kebutuhan ini tidak dimaksudkan untuk menggantikan fungsi dasar DecMed, melainkan untuk memastikan bahwa sistem dapat mendukung kontrol akses granular, segmentasi data RME, delegasi berantai, revokasi akses turunan, audit delegasi, dan pembuktian identitas aktor aktif.

Kebutuhan sistem dibagi menjadi kebutuhan fungsional dan kebutuhan nonfungsional. Kebutuhan fungsional menjelaskan kemampuan yang harus disediakan oleh sistem dari sudut pandang aktor dan proses bisnis layanan kesehatan. Sementara itu, kebutuhan nonfungsional menjelaskan batasan kualitas dan keamanan yang harus dipenuhi dalam penerapan mekanisme kontrol akses dan delegasi.

\subsubsection{Kebutuhan Fungsional Tambahan}
\label{subsubsec:kebutuhan-fungsional-tambahan}

Kebutuhan fungsional tambahan dirumuskan dari sudut pandang aktor, sama seperti pola pada kebutuhan fungsional DecMed pada Lampiran \ref{lampiran:kebutuhan-fungsional-decmed}. Perumusan berbasis aktor digunakan untuk memperjelas pihak yang dapat menerima akses, jenis data yang dapat diakses, operasi yang dapat dilakukan, serta batasan akses yang harus ditegakkan. Batasan tersebut mencakup pasien, episode RME, dataset, fungsi klinis, aksi akses, waktu berlaku, penerima delegasi, dan kedalaman delegasi. Rangkuman kebutuhan fungsional tambahan DecMed ditunjukkan pada Tabel \ref{tab:kebutuhan-fungsional-tambahan-decmed}.

\begin{table}[!ht]
	\centering
	\caption{Kebutuhan Fungsional Tambahan DecMed}
	\label{tab:kebutuhan-fungsional-tambahan-decmed}
	\begin{tabular}{|p{2cm}|p{10.5cm}|}
		\hline
		\textbf{ID} & \textbf{Deskripsi}                                                                                                                    \\
		\hline
KF-TA-01    & Pasien dapat memberikan akses awal kepada petugas rekam medis sesuai konteks layanan kesehatan.                                       \\ \hline
KF-TA-02    & Pasien dapat mencabut akses awal dan akses turunannya yang telah diberikan.                                                           \\ \hline
		KF-TA-03    & Pasien dapat melihat riwayat akses dan delegasi terhadap RME miliknya.                                                                \\ \hline
		KF-TA-04    & Petugas rekam medis dapat mendelegasikan akses terbatas kepada perawat dan dokter sesuai cakupan akses.                               \\ \hline
		KF-TA-05    & Dokter dapat membaca dan menulis segmen klinis sesuai cakupan akses dan konteks pelayanan.                                            \\ \hline
		KF-TA-06    & Dokter dapat mendelegasikan akses terbatas kepada tenaga laboratorium dan apoteker sesuai cakupan akses.                              \\ \hline
		KF-TA-07    & Perawat dapat membaca data yang relevan dan menulis segmen klinis pemeriksaan awal sesuai cakupan akses.                              \\ \hline
		KF-TA-08    & Petugas laboratorium dapat membaca data yang relevan dan menulis segmen klinis terkait pemeriksaan laboratorium sesuai cakupan akses. \\ \hline
		KF-TA-09    & Apoteker dapat membaca data yang relevan dan menulis segmen klinis terkait pelayanan obat sesuai cakupan akses.                       \\ \hline
		KF-TA-10    & Sistem dapat memverifikasi bukti kepemilikan identitas aktor aktif pada setiap permintaan akses terhadap segmen RME.                  \\ \hline
	\end{tabular}
\end{table}
\FloatBarrier

Kebutuhan fungsional dengan ID KF-TA-01 merupakan penyesuaian dari F06 dan F07 pada Lampiran \ref{lampiran:kebutuhan-fungsional-decmed}. Penyesuaian ini dilakukan karena dalam skenario layanan kesehatan kolaboratif, akses awal terhadap RME hanya diberikan langsung oleh pasien kepada petugas rekam medis. Kemudian hak akses tersebut dapat didelegasikan secara terbatas kepada aktor tenaga medis yang terlibat dalam proses pelayanan, seperti dokter, perawat, petugas laboratorium, dan apoteker.

KF-TA-11 juga merupakan penyesuaian dari F09 pada Lampiran \ref{lampiran:kebutuhan-fungsional-decmed}. Penyesuaian ini terkait dengan mekanisme revokasi akses pada skenario layanan kesehatan kolaboratif. Dalam skenario tersebut, pasien memerlukan kemampuan untuk mencabut akses awal dan akses turunannya yang telah didelegasikan.

\subsubsection{Kebutuhan Nonfungsional}
\label{subsubsec:kebutuhan-nonfungsional}

Selain kebutuhan fungsional, sistem juga memiliki kebutuhan nonfungsional untuk memastikan bahwa mekanisme kontrol akses granular dan delegasi berantai dapat diterapkan secara aman. Kebutuhan nonfungsional pada Tugas Akhir ini difokuskan pada aspek keamanan, keterlacakan, dan kompatibilitas dengan arsitektur DecMed.

\begin{table}[!ht]
	\centering
	\caption{Kebutuhan Nonfungsional Tambahan DecMed}
	\label{tab:kebutuhan-nonfungsional-tambahan-decmed}
	\small
	\begin{tabular}{|p{2cm}|p{10.5cm}|}
		\hline
		\textbf{ID} & \textbf{Deskripsi}                                                                                                                                                              \\
		\hline
		KNF-TA-01   & Sistem harus menjaga kerahasiaan data RME dengan memastikan bahwa penerima akses hanya dapat memperoleh \textit{plaintext} dari segmen yang sesuai dengan cakupan otorisasinya. \\ \hline
		KNF-TA-02   & Sistem menjamin integritas token otorisasi dan batasan akses sehingga perubahan terhadap token atau \textit{caveat} menyebabkan proses verifikasi gagal.                        \\ \hline
		KNF-TA-03   & Sistem menjamin penolakan permintaan akses jila token, rantai delegasi, atau bukti identitas aktor aktif tidak valid.                                                           \\ \hline
		KNF-TA-04   & Sistem menyediakan keterlacakan akses dan delegasi dengan mencatat informasi delegasi.                                                                                          \\ \hline
		KNF-TA-05   & Sistem tidak mengizinkan akses untuk hak akses yang dicabut.                                                                                                                    \\ \hline
		KNF-TA-06   & Sistem harus memverifikasi keterikatan bukti kepemilikan identitas terhadap aktor aktif, token akses, dan konteks permintaan akses.                                             \\ \hline
	\end{tabular}
\end{table}
\FloatBarrier

\subsubsection{Pemetaan Masalah terhadap Kebutuhan Sistem}
\label{subsubsec:pemetaan-masalah-kebutuhan}

Setelah kebutuhan fungsional dan nonfungsional dirumuskan, dilakukan pemetaan antara masalah yang telah diidentifikasi dan kebutuhan sistem yang diusulkan. Pemetaan ini digunakan untuk memastikan bahwa setiap masalah memiliki kebutuhan yang menanganinya.

\begin{table}[!ht]
	\centering
	\caption{Pemetaan Masalah terhadap Kebutuhan Sistem}
	\label{tab:pemetaan-masalah-kebutuhan}
	\small
	\begin{tabular}{|p{0.15\textwidth}|p{0.35\textwidth}|p{0.38\textwidth}|}
		\hline
		\textbf{ID Masalah} & \textbf{Kebutuhan Fungsional}                                              & \textbf{Kebutuhan Nonfungsional}                           \\ \hline
		\texttt{M-TA-01}    & \texttt{KF-TA-01}, \texttt{KF-TA-04}--\texttt{KF-TA-10}                    & \texttt{KNF-TA-01}, \texttt{KNF-TA-02}, \texttt{KNF-TA-03} \\ \hline
		\texttt{M-TA-02}    & \texttt{KF-TA-05}, \texttt{KF-TA-07}, \texttt{KF-TA-08}, \texttt{KF-TA-09} & \texttt{KNF-TA-01}                                         \\ \hline
		\texttt{M-TA-03}    & \texttt{KF-TA-01}, \texttt{KF-TA-04}, \texttt{KF-TA-06}                    & \texttt{KNF-TA-04}                                         \\ \hline
		\texttt{M-TA-04}    & \texttt{KF-TA-04}, \texttt{KF-TA-06}                                       & \texttt{KNF-TA-02}, \texttt{KNF-TA-03}                     \\ \hline
		\texttt{M-TA-05}    & \texttt{KF-TA-05}, \texttt{KF-TA-07}, \texttt{KF-TA-08}, \texttt{KF-TA-09} & \texttt{KNF-TA-01}, \texttt{KNF-TA-03}                     \\ \hline
		\texttt{M-TA-06}    & \texttt{KF-TA-02}, \texttt{KF-TA-03}                                       & \texttt{KNF-TA-04}, \texttt{KNF-TA-05}                     \\ \hline
		\texttt{M-TA-07}    & \texttt{KF-TA-10}                                                          & \texttt{KNF-TA-03}, \texttt{KNF-TA-06}                     \\ \hline
	\end{tabular}
\end{table}

Kebutuhan fungsional dan nonfungsional tambahan tersebut menjadi dasar bagi perancangan solusi pada tahap berikutnya.
 
 
\section{Analisis Solusi}
\label{sec:analisis-solusi}

DecMed sudah memiliki fondasi yang sesuai untuk sistem RME terdesentralisasi melalui IOTA, IPFS, dan PRE. Oleh karena itu, solusi yang dianalisis pada bagian ini diarahkan untuk memperkuat lapisan kontrol akses agar menjadi lebih granular dan memfasilitasi kemampuan delegasi \textit{offline}.

Berdasarkan kebutuhan sistem pada Subbab \ref{subsec:kebutuhan-sistem}, solusi yang dianalisis pada Tugas Akhir ini dikelompokkan ke dalam enam aspek utama, yaitu token otorisasi, segmentasi objek data, delegasi akses \textit{offline}, serta penegakan otoritas. Pengelompokan ini diperlukan karena kontrol akses granular dan kemampuan delegasi \textit{offline} tidak dapat ditegakkan hanya melalui satu komponen. Token otorisasi diperlukan untuk menyatakan batasan akses, segmentasi data diperlukan agar batasan tersebut memiliki objek penegakan yang jelas, sedangkan \textit{Server} PRE diperlukan sebagai titik penegakan sebelum data terenkripsi diambil dan diproses \textit{re-encryption}. Dengan demikian, analisis solusi pada bagian ini berfungsi sebagai dasar pemilihan pendekatan sebelum rancangan teknis dijelaskan pada Subbab \ref{sec:rancangan-solusi}.

\subsection{Analisis Delegasi Akses dan Atenuasi \textit{Offline}}
\label{subsec:analisis-delegasi-atenuasi-offline}

Delegasi akses diperlukan berdasarkan hasil analisis pada Subbab \ref{subsubsec:analisis-skenario-layanan-klinis}. Dalam skenario DecMed sebelumnya, pasien memberikan akses kepada aktor tertentu secara langsung. Pola tersebut kurang sesuai untuk layanan klinis kolaboratif karena kebutuhan akses dapat berpindah antar aktor selama episode berlayanan berlangsung. Jika setiap perpindahan akses harus selalu diterbitkan langsung oleh pasien, maka proses pelayanan menjadi bergantung pada ketersediaan pasien.

Alternatif solusi pertama yang dipertimbangkan adalah melakukan delegasi akses melalui \textit{client} dan mencatat hasil delegasi pada penyimpanan \textit{on-chain}. Pada pendekatan ini, aktor yang memiliki token induk membentuk token turunan secara lokal dengan menambahkan batasan baru terhadap hak akses yang dimiliki. Token turunan tersebut kemudian dibagikan kepada penerima delegasi melalui penyimpanan \textit{on-chain}. Pendekatan ini dapat mengurangi keterlibatan pasien karena pemegang token dapat menurunkan hak akses yang sudah diterima sebelumnya.

Namun, pendekatan tersebut menjadikan \textit{client} sebagai komponen yang menentukan hasil delegasi. Kondisi ini kurang ideal karena \textit{client} dapat dimanipulasi untuk melakukan delegasi sesuka hati. Akibatnya, sistem berpotensi mencatat delegasi yang tidak konsisten dengan hak akses induk, meskipun token turunan tersebut kemudian dapat ditolak ketika digunakan untuk mengakses \textit{Server} PRE.

Alternatif berikutnya adalah melibatkan \textit{Server} PRE sebagai komponen yang memverifikasi permintaan delegasi sebelum token turunan dibentuk. Pada pendekatan ini, aktor delegator tetap menginisiasi delegasi melalui \textit{client}. \textit{Client} mengirimkan token induk yang dimiliki \textit{delegator}, lalu token dan hak akses \textit{delegator} diverifikasi oleh \textit{Server} PRE.

Jika permintaan delegasi dinyatakan sah, \textit{Server} PRE membentuk hak akses turunan dengan menambahkan dengan batasan baru yang mempersempit hak akses. PRE kemudian mengembalikan hak akses tersebut kepada \textit{client}. Setelah menerima hasil atenuasi dari PRE, \textit{client} kemudian dapat melakukan proses pencatatan lebih lanjut pada penyimpanan \textit{on-chain}. Dengan demikian, PRE hanya menjadi verifikator dan pembentuk hak akses turunan, sedangkan pencatatan delegasi tetap dilakukan oleh \textit{client}.

Berdasarkan pertimbangan tersebut, Tugas Akhir ini memilih mekanisme delegasi akses yang dimediasi oleh \textit{Server} PRE. Mekanisme ini memanfaatkan konsep atenuasi sehingga token turunan hanya dapat memiliki cakupan yang sama atau lebih sempit daripada token induk. Istilah \textit{offline} pada atenuasi delegasi dalam Tugas Akhir ini tidak dimaknai sebagai proses yang sepenuhnya dilakukan tanpa internet, melainkan sebagai proses delegasi yang tidak memerlukan keterlibatan pasien. Setelah pasien memberikan akses awal, aktor yang sah dapat mengajukan delegasi terbatas melalui \textit{client}, sedangkan PRE memverifikasi dan membentuk hak akses turunan berdasarkan hak akses induk yang masih berlaku.

Pemilihan mekanisme ini memberikan batas tanggung jawab yang lebih jelas. \textit{client} berperan sebagai inisiator delegasi dan pencatat delegasi pada penyimpanan \textit{on-chain}. \textit{Server} PRE berperan sebagai pembentuk hak akses turunan. Penyimpanan \textit{on-chain} mencatat riwayat delegasi dan bukti akses. Dengan pembagian tersebut, proses pelayanan tidak perlu selalu melibatkan pasien, tetapi pembentukan hak akses turunan tetap melewati validasi PRE sebelum dicatat dan digunakan.
 \subsection{Analisis Pemilihan Token Otorisasi}
\label{subsec:analisis-pemilihan-token-otorisasi}

Berdasarkan masalah yang telah diidentifikasi pada Subbab \ref{subsec:identifikasi-masalah},
mekanisme otorisasi pada DecMed perlu mendukung kontrol akses granular dan memiliki mekanisme atuneasi. Selain itu, pada Subbab \ref{subsec:analisis-delegasi-atenuasi-offline} dibahas bahwa mekanisme delegasi antaraktor perlu mendukung atenuasi \textit{offline}.
Oleh karena itu, diperlukan analisis perbandingan terhadap beberapa pendekatan token otorisasi untuk menentukan teknologi token otorisasi yang paling sesuai dengan kebutuhan sistem.

Alternatif token otorisasi yang dibandingkan pada analisis ini adalah JWT, PASETO, Macaroons, dan Biscuit. JWT dipilih sebagai pembanding karena merupakan bentuk token yang digunakan pada rancangan DecMed awal dan umum digunakan untuk merepresentasikan klaim otorisasi. PASETO dipilih karena memiliki tujuan penggunaan yang serupa dengan JWT, tetapi dengan kriptografi yang \textit{opinionated}. Macaroons dianalisis sebagai kandidat utama karena menyediakan mekanisme \textit{caveat} yang memungkinkan atenuasi \textit{offline}. Biscuit dianalisis sebagai pembanding karena mendukung verifikasi terdesentralisasi dan atenuasi \textit{offline}, tetapi menggunakan pendekatan yang lebih ekspresif melalui aturan berbasis \textit{logic language}.

Kriteria evaluasi antar token diturunkan dari kebutuhan sebagai berikut:
\begin{enumerate}
	\item \textbf{Representasi batasan akses granular}

	      Menilai apakah token dapat membawa batasan akses yang dibutuhkan agar token benar-benar dapat mengekspresikan pembatasan akses sesuai kebutuhan.
	\item \textbf{Atenuasi \textit{offline}}

	      Menilai apakah pemegang token dapat menambahkan batasan baru tanpa harus melibatkan pasien dan menerbitkan token baru. Sesuai dengan solusi dari Subbab \ref{subsec:analisis-delegasi-atenuasi-offline}.
	\item \textbf{Jaminan token turunan tidak memperluas hak akses}

	      Menganalisis jenis kunci kriptografi yang digunakan dan seberapa kuat sistem delegasi token mencegah token turunan digunakan untuk mengakses lebih banyak atau lebih luas daripada token induknya.
	\item \textbf{Model verifikasi dan pengelolaan kunci}

	      Mengkaji bagaimana masing-masing token diverifikasi dan bagaimana pengelolaan kuncinya.
	\item \textbf{Kesesuaian ekspresivitas dengan aturan akses DecMed}

	      DecMed membutuhkan token yang dapat mengekspresikan batasan akses secara langsung sebagai atribut pada token itu sendiri.
	\item \textbf{Kompleksitas integrasi ke PRE}

	      membahas tingkat kesulitan integrasi ke dalam sistem PRE sebagai penegak otoritas akses.
\end{enumerate}

Tabel \ref{tab:perbandingan-token-otorisasi} menunjukkan perbandingan antar token otorisasi berdasarkan kriteria evaluasi.

\begin{table}[!ht]
	\centering
	\caption{Perbandingan Pendekatan Token Otorisasi}
	\label{tab:perbandingan-token-otorisasi}
	\scriptsize
	\begin{tabular}{|p{0.22\textwidth}|p{0.15\textwidth}|p{0.15\textwidth}|p{0.15\textwidth}|p{0.15\textwidth}|}
		\hline
		\textbf{Kriteria} & \textbf{JWT}                                                                  & \textbf{PASETO} & \textbf{Macaroons} & \textbf{Biscuit} \\ \hline
		Representasi batasan akses granular
		                  & Melalui klaim
		                  & Melalui klaim
		                  & Melalui \textit{caveats}
		                  & Melalui \textit{fact}, \textit{rule}, dan \textit{check}                                                                                \\ \hline

		Atenuasi \textit{offline}
		                  & Tidak didukung secara bawaan
		                  & Tidak didukung secara bawaan
		                  & Didukung melalui penambahan \textit{caveat}
		                  & Didukung melalui penambahan blok pembatasan                                                                                             \\ \hline

		Jaminan token turunan tidak memperluas hak akses
		                  & Bergantung pada penerbitan token baru atau logika \textit{server}
		                  & Bergantung pada penerbitan token baru atau logika \textit{server}
		                  & Ditegakkan melalui verifikasi seluruh \textit{caveat}
		                  & Ditegakkan melalui blok dan pemeriksaan kriptografis                                                                                    \\ \hline

		Model verifikasi dan pengelolaan kunci
		                  & Dapat menggunakan kunci simetris atau asimetris
		                  & Dapat menggunakan mode lokal atau publik
		                  & Berbasis kunci simetris
		                  & Berbasis kunci asimetris                                                                                                                \\ \hline

		Kesesuaian ekspresivitas dengan kebutuhan DecMed
		                  & Cukup untuk klaim statis, kurang sesuai untuk delegasi lokal
		                  & Cukup untuk klaim statis, kurang sesuai untuk delegasi lokal
		                  & Sesuai untuk batasan berbasis atribut yang langsung
		                  & Lebih ekspresif, tetapi melebihi kebutuhan aturan akses yang relatif langsung                                                           \\ \hline

		Kompleksitas integrasi
		                  & Rendah
		                  & Rendah
		                  & Sedang
		                  & Tinggi                                                                                                                                  \\ \hline
	\end{tabular}

\end{table}

JWT dan PASETO dapat digunakan untuk membawa informasi otorisasi dalam bentuk klaim, seperti identitas subjek, cakupan akses, dan waktu kedaluwarsa. Namun, keduanya tidak menyediakan mekanisme bawaan untuk menambahkan batasan baru pada token yang sudah diterbitkan. Akibatnya, delegasi hak akses yang lebih sempit tetap membutuhkan penerbitan token baru oleh pihak penerbit atau mekanisme tambahan di sisi \textit{server}. Karakteristik tersebut kurang sesuai dengan kebutuhan delegasi \textit{offline} pada Tugas Akhir ini.

Macaroons dan Biscuit lebih dekat dengan kebutuhan karena keduanya mendukung atenuasi \textit{offline}. Perbedaan utama keduanya terletak pada model pembatasan dan model verifikasi. Macaroons menggunakan \textit{caveats} untuk menambahkan batasan pada token, sedangkan Biscuit menggunakan blok berisi \textit{fact}, \textit{rule}, dan \textit{check}. Model Biscuit memberikan ekspresivitas yang lebih tinggi, tetapi kebutuhan pembatasan akses pada DecMed masih dapat direpresentasikan melalui klaim pada token, seperti identitas pasien, identitas aktor, cakupan data, aksi, dan masa berlaku. Pembatasan tersebut hanya memanfaatkan \textit{fact} saja, karena pemeriksaan otoritas dari klaim pada token sudah cukup ditangani oleh \textit{server} PRE.Oleh karena itu, penggunaan Biscuit dinilai menambah kompleksitas integrasi yang tidak diperlukan untuk ruang lingkup Tugas Akhir ini.

Dari sisi model verifikasi, Macaroons lebih sesuai dengan rancangan DecMed karena proses verifikasi token tetap dipusatkan pada \textit{Server} PRE sebagai titik penegakan otorisasi. Dengan model tersebut, kunci verifikasi tidak perlu disebarkan ke banyak layanan. Biscuit lebih sesuai untuk arsitektur yang membutuhkan verifikasi token oleh banyak komponen berbeda menggunakan kunci publik. Pada rancangan DecMed, kebutuhan tersebut tidak menjadi fokus karena seluruh akses data RME tetap melewati PRE.

Berdasarkan perbandingan tersebut, Macaroons dipilih sebagai mekanisme token otorisasi pada DecMed karena memberikan keseimbangan antara representasi batasan akses granular, dukungan delegasi terbatas melalui atenuasi \textit{offline}, dan kompleksitas integrasi yang sesuai dengan ruang lingkup Tugas Akhir. Melalui mekanisme \textit{caveats}, token dapat membawa batasan akses dan pemegang token dapat menambahkan batasan baru secara lokal untuk mempersempit hak akses sebelum token didelegasikan kepada aktor lain. \subsection{Analisis Segmentasi Data RME}
\label{subsec:analisis-segmentasi-data-rme}

Masalah granularitas tidak hanya berada pada token, tetapi juga pada struktur data yang dilindungi. Sesuai analisis pada Subbab \ref{subsubsec:keterbatasan-segmentasi-data-decmed}, jika data RME disimpan sebagai satu \textit{encrypted blob}, maka pemberian akses terhadap blob tersebut akan membuka seluruh isi data meskipun aktor hanya membutuhkan sebagian informasi. Oleh karena itu, implementasi pemrosesan data pada DecMed perlu disesuaikan untuk memenuhi kebutuhan akses granular.

Alternatif pertama adalah melakukan pemfilteran setelah data klinis diambil dari penyimpanan \textit{off-chain}. Pada pendekatan ini, data tetap disimpan sebagai satu \textit{encrypted blob}, kemudian sistem memilih bagian data yang boleh ditampilkan berdasarkan hak akses pengguna. Pendekatan ini kurang sesuai dengan karakteristik DecMed. IPFS tidak menyediakan mekanisme pencarian atau pemfilteran isi data secara langsung. Selain itu, DecMed menyimpan data dalam bentuk terenkripsi, sehingga isi data baru dapat diproses setelah didekripsi. Jika pemfilteran dilakukan di \textit{client}, maka keseluruhan \textit{plaintext} data klinis sudah terbuka terlebih dahulu sebelum batas akses diterapkan.

Alternatif yang dipilih adalah melakukan segmentasi data klinis sebelum data disimpan pada penyimpanan \textit{off-chain}. Pada pendekatan ini, data klinis dipisahkan menjadi beberapa segmen. Setiap segmen memiliki metadata yang disimpan pada penyimpanan \textit{on-chain}. Metadata tersebut digunakan oleh sistem untuk memilih segmen yang sesuai dengan cakupan akses sebelum data diambil dari IPFS. Dengan cara ini, proses pengambilan data menjadi lebih terarah karena sistem hanya mengambil segmen yang relevan dengan hak akses pengguna.

Penentuan kategori segmentasi tidak dilakukan secara arbitrer. Opsi pertama adalah mengikuti pembagian data pada penelitian terkait yang dibahas pada Subbab \ref{subsec:composable-nft-rme}. Namun, pembagian tersebut kurang sesuai dijadikan dasar segmentasi pada Tugas Akhir ini karena ditujukan untuk pembuktian konseptual, bukan untuk merepresentasikan struktur data pelayanan medis. Oleh karena itu, rancangan segmentasi mengacu pada KMK 1423/2022 tentang Pedoman Variabel dan Metadata pada Penyelenggaraan Rekam Medis Elektronik yang telah dibahas pada Subbab \ref{subsec:isi-rme}.

Berdasarkan pedoman tersebut, data RME dapat dibedakan berdasarkan \textit{dataset} pelayanan. Dalam cakupan Tugas Akhir ini, \textit{dataset category} yang digunakan meliputi:

\begin{itemize}[]
	\item Rawat Jalan
	\item Rawat Inap
	\item Laboratorium
	\item Apotek
\end{itemize}
Pembagian berdasarkan \textit{dataset} pelayanan dipilih karena setiap layanan memiliki jenis data dan kebutuhan akses yang berbeda. Sebagai contoh, petugas laboratorium membutuhkan akses tulis terhadap \textit{dataset} laboratorium, sedangkan apoteker membutuhkan akses tulis terhadap \textit{dataset} apotek. Dengan pemisahan ini, sistem dapat membedakan konteks layanan sebelum menentukan segmen data yang dapat dibaca atau ditulis oleh aktor.

Selain berdasarkan \textit{dataset}, data klinis juga disegmentasi berdasarkan kategori fungsi klinis. \textit{Function category} digunakan untuk membedakan fungsi data di dalam suatu \textit{dataset}, seperti anamnesis, pemeriksaan fisik, pemeriksaan penunjang, diagnosis, persetujuan tindakan, atau terapi. Metadata fungsi klinis digunakan untuk menentukan apakah suatu segmen data relevan dengan peran dan kebutuhan akses tenaga medis.

Kombinasi \textit{dataset category} dan \textit{function category} dipilih agar segmentasi tidak terlalu kasar seperti pemisahan per dokumen penuh, tetapi juga tidak terlalu rinci hingga tingkat \textit{field}. Segmentasi hingga tingkat \textit{field} akan membuat kontrol akses bergantung pada setiap elemen data yang tersimpan. Pendekatan tersebut tidak dipilih karena kebutuhan akses pada rancangan ini dinilai berdasarkan konteks layanan dan fungsi klinis, bukan berdasarkan setiap \textit{field} secara terpisah.

Sebagai contoh, Formulir Umum / Asesmen Awal Rawat Jalan dapat dipetakan ke fungsi klinis seperti anamnesis, pemeriksaan fisik, dan pemeriksaan psikologi. Sementara itu, Pemeriksaan Spesialistik dapat dipetakan ke fungsi klinis seperti riwayat penggunaan obat, rencana rawat, instruksi medik dan keperawatan, pemeriksaan penunjang, diagnosis, persetujuan tindakan, dan terapi. Pemetaan lengkap kategori fungsi klinis dapat dilihat pada Lampiran \ref{lampiran:metadata-rme}.

Melalui pendekatan ini, data \textit{off-chain} tidak lagi diperlakukan sebagai satu objek besar, tetapi sebagai kumpulan segmen klinis dengan metadata yang terdefinisi. Segmentasi tersebut menyediakan objek penegakan yang lebih jelas bagi token otorisasi karena batasan akses pada token dapat dicocokkan langsung dengan \textit{dataset category} dan \textit{function category} pada metadata \textit{on-chain}. \subsection{Analisis Penegakan Otoritas pada \textit{Server} PRE}
\label{subsec:analisis-penegakan-otoritas-pre}

Pada DecMed, PRE digunakan untuk mengubah data terenkripsi dari pasien ke penerima akses tanpa sekalipun membuka \textit{plaintext}. Pada implementasi DecMed sebelumnya, PRE berfokus melakukan verifikasi token dan menjalankan proses \textit{re-encryption}. \textit{Server} PRE dipilih sebagai titik penegakan otoritas karena komponen ini berada pada jalur kritis sebelum data terenkripsi diproses. Pengguna tidak memperoleh \textit{plaintext} langsung dari penyimpanan \textit{off-chain}, melainkan membutuhkan proses \textit{re-encryption} yang dikendalikan oleh PRE. Oleh karena itu, PRE dapat berperan sebagai titik penegakan kebijakan akses yang memverifikasi token dan menolak permintaan sebelum data diambil dari IPFS apabila otorisasi tidak valid.

\subsubsection{Analisis Verifikasi Aktor dan Operasi Tulis}
\label{subsec:analisis-verifikasi-aktor-operasi-tulis}

Pada DecMed sebelumnya, \textit{role} \texttt{MedicalPersonnel} digunakan untuk merepresentasikan tenaga medis secara umum, dan memberikan hak untuk mengakses data klinis pasien. Namun, \textit{role} tersebut belum cukup untuk membedakan dokter, perawat, petugas laboratorium, dan apoteker. Padahal, keempat aktor tersebut memiliki kewenangan tulis yang berbeda. Sebagai contoh, dokter dapat menulis diagnosis dan terapi sedangkan petugas laboratorium hanya dapat menulis hasil laboratorium.

Oleh karena itu, diperlukan verifikasi tambahan terhadap \textit{subrole} tenaga medis, terutama untuk operasi tulis. \textit{Server} PRE tetap perlu memastikan bahwa aktor aktif memiliki \textit{subrole} yang sesuai dengan \textit{dataset category} dan \textit{function category} yang akan ditulis. Verifikasi ini mencegah kondisi ketika aktor membawa token dengan izin tulis, tetapi secara kewenangan profesional tidak berhak memperbarui segmen data tertentu.

\subsubsection{Analisis Verifikasi Bukti Identitas Aktor Aktif}
\label{subsec:analisis-verifikasi-bukti-identitas-aktor-aktif}

Pada DecMed sebelumnya, proses otorisasi berfokus pada pemeriksaan klaim yang terdapat pada token. Namun, mekanisme tersebut belum memastikan bahwa pengirim \textit{request} merupakan aktor yang identitasnya tercantum pada token. Kondisi ini menimbulkan kebutuhan verifikasi tambahan, yaitu verifikasi bukti identitas aktor aktif sebagaimana dirumuskan pada KNF-TA-06.

Identitas pengguna pada DecMed direpresentasikan melalui \textit{wallet} yang dibuat saat pendaftaran akun. \textit{Wallet address} menjadi identifikator unik yang dapat membedakan setiap pengguna pada sistem. Oleh karena itu, solusi yang diusulkan adalah menggunakan \textit{wallet} sebagai dasar verifikasi identitas aktor aktif.

Verifikasi dilakukan dengan meminta aktor aktif membuktikan kepemilikan \textit{wallet} yang sesuai dengan klaim pada token. Dengan mekanisme ini, \textit{request} tidak hanya membawa token otorisasi, tetapi juga membawa bukti bahwa pengirim \textit{request} menguasai \textit{wallet} dari aktor yang sedang menggunakan hak akses tersebut. Hasil verifikasi ini kemudian digunakan oleh \textit{Server} PRE untuk menentukan apakah aktor aktif berwenang melanjutkan operasi akses terhadap data RME.

\subsubsection{Analisis Revokasi dan Audit Delegasi}
\label{subsec:analisis-revokasi-audit-delegasi}

Delegasi akses menyebabkan token tidak lagi berdiri sendiri, tetapi membentuk relasi antara token induk dan token turunan. Oleh karena itu, mekanisme revokasi pada DecMed perlu diperluas agar pencabutan token induk dapat berdampak pada token turunan yang berasal dari token tersebut. Tanpa pencabutan berantai, aktor yang memperoleh token turunan masih dapat menggunakan akses meskipun akses asalnya telah dicabut oleh pasien.

Revokasi berantai diperlukan untuk mempertahankan kendali pasien sebagai pemilik data. Ketika pasien mencabut akses awal, sistem harus dapat mengenali bahwa token-token turunan yang berada dalam rantai delegasi yang sama tidak lagi memiliki dasar kewenangan yang sah. Dengan demikian, pemeriksaan revokasi tidak cukup hanya dilakukan terhadap token yang dikirim pada saat permintaan akses, tetapi juga perlu mempertimbangkan hubungan token tersebut dengan token asalnya.

Selain revokasi, delegasi juga membutuhkan audit yang mencatat hubungan antaraktor. Audit tidak cukup hanya mencatat bahwa suatu segmen diakses, tetapi juga perlu mencatat sumber kewenangan akses tersebut, termasuk aktor pemberi delegasi, aktor penerima delegasi, waktu delegasi, cakupan akses, dan token yang digunakan. Informasi ini diperlukan agar pasien dan sistem dapat menelusuri bagaimana akses terhadap data RME terbentuk selama proses pelayanan. \subsection{Pemetaan Masalah terhadap Solusi}
\label{subsec:pemetaan-kebutuhan-solusi}
\label{subsec:pemetaan-masalah-solusi}

Setelah pendekatan solusi dianalisis, setiap solusi pada bagian ini diberi ID \texttt{S-TA} untuk menjaga ketertelusuran terhadap masalah yang telah diidentifikasi pada Tabel \ref{tab:identifikasi-masalah-decmed}. Daftar solusi utama ditunjukkan pada Tabel \ref{tab:daftar-solusi-ta}.

\begin{table}[!ht]
	\centering
	\caption{Daftar Solusi yang Diusulkan Berdasarkan Hasil Analisis}
	\label{tab:daftar-solusi-ta}
	\small
	\begin{tabular}{|p{0.12\textwidth}|p{0.5\textwidth}|p{0.25\textwidth}|}
		\hline
		\textbf{ID}      & \textbf{Nama Solusi}                                                                                                        & \textbf{Masalah Terkait}                                                                 \\ \hline
		\texttt{S-TA-01} & Token otorisasi berbasis Macaroons dengan \textit{caveats}                                                                  & \texttt{M-TA-01}, \texttt{M-TA-04}                                                       \\ \hline
		\texttt{S-TA-02} & Segmentasi data RME berdasarkan \texttt{dataset\_category} dan \texttt{function\_category}                                  & \texttt{M-TA-02}                                                                         \\ \hline
		\texttt{S-TA-03} & Delegasi akses melalui atenuasi \textit{offline} Macaroons                                                                  & \texttt{M-TA-03}, \texttt{M-TA-04}                                                       \\ \hline
		\texttt{S-TA-04} & Penegakan akses pada \textit{Server} PRE berbasis \textit{effective capability}, \textit{wallet proof}, dan metadata segmen & \texttt{M-TA-01}, \texttt{M-TA-02}, \texttt{M-TA-04}, \texttt{M-TA-05}, \texttt{M-TA-07} \\ \hline
		\texttt{S-TA-05} & Validasi subrole tenaga medis terhadap dataset untuk operasi tulis                                                          & \texttt{M-TA-05}                                                                         \\ \hline
		\texttt{S-TA-06} & Pencabutan akses berantai dan audit log akses dan delegasi                                                                  & \texttt{M-TA-06}                                                                         \\ \hline
	\end{tabular}
\end{table}



%
  
\section{Rancangan Solusi}
\label{sec:rancangan-solusi}

Berdasarkan hasil analisis solusi pada Subbab \ref{sec:analisis-solusi}, rancangan solusi difokuskan pada pengembangan mekanisme kontrol akses granular dan delegasi terbatas pada DecMed melalui integrasi Macaroons. Rancangan ini mempertahankan fondasi utama DecMed, yaitu IOTA sebagai penyimpanan \textit{on-chain}, IPFS sebagai penyimpanan \textit{off-chain}, dan PRE sebagai mekanisme distribusi akses terhadap data terenkripsi. Perancangan ini mencakup arsitektur sistem, token otorisasi berbasis Macaroons, segmentasi data \textit{off-chain}, penegakan otoritas pada \textit{Server} PRE, validasi subrole tenaga medis, pencabutan akses, audit, alur akses data, dan delegasi berantai dengan atenuasi \textit{offline}.

\subsection{Rancangan Arsitektur Solusi}
\label{subsec:rancangan-arsitektur-solusi}



Arsitektur solusi mempertahankan keseluruhan arsitektur kerangka kerja DecMed. Tiga \textit{client} utama DecMed, yaitu \textit{Patient Client}, \textit{Hospital Client}, dan \textit{Ministry Client} tetap ada namun disesuaikan sesuai kebutuhan Tugas Akhir ini. \textit{Patient Client} digunakan pasien untuk menerbitkan token akses awal. \textit{Hospital Client} digunakan oleh petugas rekam medis, dokter, perawat, apoteker dan petugas laboratorium untuk mengakses atau memperbarui data sesuai kewenangannya. \textit{Ministry Client} tetap berada pada ruang lingkup manajemen fasyankes sebagaimana rancangan DecMed awal.

Perubahan utama berada pada \textit{server} PRE, token otorisasi berbasis Macaroons, dan implementasi Move \textit{smart contract}. Pada rancangan ini, token akses diganti dengan Macaroon yang membawa daftar \textit{caveats}. Setiap \textit{request} akses data harus menyertakan Macaroon. \textit{Server} PRE kemudian memverifikasi hak akses, lalu menjalankan proses \textit{re-encryption} untuk operasi baca dan enkripsi data ketika operasi tulis. Implementasi Move \textit{smart contract} disesuaikan untuk menyimpan metadata yang diperlukan oleh mekanisme segmentasi dan kontrol akses.

Pada rancangan ini, \textit{server} PRE juga berperan sebagai komponen pembentuk token Macaroon turunan pada proses delegasi. \textit{Client delegator} mengirimkan permintaan atenuasi kepada \textit{server} PRE. \textit{Server} PRE memverifikasi \textit{request} delegasi. Apabila valid, \textit{server} PRE mengembalikan token turunan untuk diproses lebih lanjut oleh \textit{client delegator} untuk dicatat pada penyimpanan \textit{on-chain}.

Skema penyimpanan \textit{on-chain} dan \textit{off-chain} disesuaikan agar dapat mendukung kontrol akses granular sesuai hasil analisis solusi pada Subbab \ref{subsec:analisis-segmentasi-data-rme}. Penyimpanan data pada sistem yang diusulkan dibagi menjadi dua lapisan. Penyimpanan pada lapisan \textit{on-chain} disesuaikan dengan mengubah implementasi dari Move \textit{smart contract}. Sedangkan penyimpanan pada lapisan \textit{off-chain} disesuaikan dengan mengubah mekanisme penyimpanan data yang diproses pada \textit{hospital client}.

Gambaran arsitektur sistem secara keseluruhan dapat dilihat pada gambar \ref{fig:arsitektur-solusi}. Warna merah digunakan untuk merepresentasikan komponen yang diubah dari rancangan DecMed awal. Legenda dari diagram arsitektur ini dapat dilihat pada lampiran \ref{lampiran:legenda-diagram-arsitektur-solusi}.

\begin{figure}[!ht]
	\centering
	\includegraphics[width=1\textwidth]{resources/chapter-3/rancangan/decmed_archi_changed.png}
	\caption{Arsitektur Solusi (Dokumentasi Penulis)}
	\label{fig:arsitektur-solusi}
\end{figure} \subsection{Rancangan Segmentasi Data RME}
\label{subsec:rancangan-segmentasi-data-rme}



Sesuai dengan analisis solusi pada Subbab \ref{subsec:pemetaan-kebutuhan-solusi}, rancangan segmentasi data RME memecah data klinis menjadi beberapa segmen terenkripsi. Setiap segmen dibentuk berdasarkan kombinasi \textit{dataset category} dan \textit{function category}. \textit{Dataset category} merepresentasikan konteks layanan, seperti Rawat Jalan, Rawat Inap, Laboratorium, dan Apotek. \textit{Function category} merepresentasikan fungsi data di dalam \textit{dataset}, seperti anamnesis, diagnosis, terapi, pemeriksaan penunjang, atau \textit{dispensing}. Kombinasi kedua kategori tersebut menjadi unit data terkecil yang digunakan dalam proses kontrol akses.

Tabel \ref{tab:mapping-dataset-category} menunjukkan \textit{dataset category} yang digunakan pada rancangan ini.

\begin{table}[!ht]
	\centering
	\caption{Tabel Pemetaan \textit{Dataset Category}}
	\label{tab:mapping-dataset-category}
	\small
	\begin{tabular}{|p{0.23\textwidth}|p{0.27\textwidth}|}
		\hline
		\textbf{\textit{Dataset Category}} & \textbf{\textit{Enum Value}} \\ \hline
		Rawat Jalan                        & \texttt{RAWAT\_JALAN}        \\ \hline
		Rawat Inap                         & \texttt{RAWAT\_INAP}         \\ \hline
		Laboratorium                       & \texttt{LABORATORIUM}        \\ \hline
		Apotek                             & \texttt{APOTEK}              \\ \hline
	\end{tabular}
\end{table}

Setiap \textit{dataset category} kemudian dibagi lagi berdasarkan \textit{function category}. Kategori fungsi klinis mengacu pada hasil klasifikasi metadata RME yang disajikan pada Lampiran \ref{lampiran:metadata-rme}. Selain fungsi klinis, rancangan ini juga tetap menyediakan kategori untuk dokumentasi administratif. Dokumentasi administratif tidak menjadi fokus utama Tugas Akhir ini, sehingga dokumen administratif dikelompokkan ke dalam satu kategori umum, yaitu \texttt{ADMINISTRATIVE\_GENERAL}. Kategori ini tetap disediakan karena beberapa kebutuhan fungsional DecMed masih bergantung pada keberadaan data administratif, khususnya F01, F06, F07, F14, dan F21 pada Tabel di Lampiran \ref{tab:kebutuhan-fungsional-decmed}.

Tabel \ref{tab:mapping-function-category} menunjukkan \textit{function category} yang digunakan pada implementasi.

\begin{table}[!ht]
	\centering
	\caption{Tabel Pemetaan \textit{Function Category}}
	\label{tab:mapping-function-category}
	\small
	\begin{tabular}{|p{0.4\textwidth}|p{0.45\textwidth}|}
		\hline
		\textbf{\textit{Function Category}} & \textbf{\textit{Enum Value}}                \\ \hline

		Dokumentasi Administratif
		                                    & \texttt{ADMINISTRATIVE\_GENERAL}
		\\ \hline

		Anamnesis
		                                    & \texttt{ANAMNESIS}
		\\ \hline

		Pemeriksaan Fisik
		                                    & \texttt{PEMERIKSAAN\_FISIK}
		\\ \hline

		Pemeriksaan Psikologis, Sosial Ekonomi, Spiritual
		                                    & \texttt{PEMERIKSAAN\_PSIKOLOGIS}
		\\ \hline

		Riwayat Penggunaan Obat
		                                    & \texttt{RIWAYAT\_PENGGUNAAN\_OBAT}
		\\ \hline

		Rencana Rawat
		                                    & \texttt{RENCANA\_RAWAT}
		\\ \hline

		Perencanaan Pemulangan Pasien
		                                    & \texttt{PERENCANAAN\_PEMULANGAN}
		\\ \hline

		Instruksi Medik dan Keperawatan
		                                    & \texttt{INSTRUKSI\_MEDIK\_DAN\_KEPERAWATAN}
		\\ \hline

		Pemeriksaan Penunjang
		                                    & \texttt{PEMERIKSAAN\_PENUNJANG}
		\\ \hline

		Diagnosis
		                                    & \texttt{DIAGNOSIS}
		\\ \hline

		\textit{Informed Consent}
		                                    & \texttt{INFORMED\_CONSENT}
		\\ \hline

		Terapi
		                                    & \texttt{TERAPI}
		\\ \hline

		Laboratorium
		                                    & \texttt{LABORATORIUM}
		\\ \hline

		Peresepan
		                                    & \texttt{PERESEPAN}
		\\ \hline

		\textit{Dispensing}
		                                    & \texttt{DISPENSING}
		\\ \hline
	\end{tabular}

\end{table}

Pada rancangan ini, \textit{function category} tidak selalu dibentuk dari setiap subkelompok metadata secara langsung. Beberapa subkelompok metadata dapat digabung ke dalam satu \textit{function category} selama masih memiliki fungsi pelayanan dan kebutuhan akses yang sama. Sebagai contoh, subkelompok pada \textit{dataset} Laboratorium, seperti permintaan pemeriksaan, spesimen klinis, hasil pemeriksaan, serta interpretasi dan validasi hasil, dikelompokkan ke dalam \textit{function category} \texttt{LABORATORIUM} karena seluruhnya merepresentasikan satu rangkaian fungsi pemeriksaan laboratorium. Sementara itu, \textit{dataset} Apotek dibedakan menjadi \texttt{PERESEPAN} dan \texttt{DISPENSING} karena keduanya merepresentasikan dua fungsi yang berbeda, yaitu pencatatan resep dan proses penyiapan serta penyerahan obat.

Secara penyimpanan, satu entri metadata \textit{on-chain} menunjuk pada satu segmen data terenkripsi di penyimpanan \textit{off-chain}. Rincian skema metadata \textit{on-chain} disajikan pada Tabel \ref{tab:onchain-schema}, sedangkan skema objek data segmen \textit{off-chain} disajikan pada Tabel \ref{tab:offchain-schema}.

Rancangan ini membuat \textit{Server} PRE dapat menegakkan akses pada tingkat segmen. Ketika aktor meminta akses terhadap data RME, \textit{Server} PRE mencocokkan \textit{caveats} pada token Macaroon dengan metadata segmen yang diminta. Sebagai contoh, token yang hanya mengizinkan akses tulis terhadap \textit{dataset} Laboratorium tidak dapat digunakan untuk menulis segmen pada \textit{dataset} Rawat Jalan.
 \subsection{Rancangan Token Otorisasi Berbasis Macaroons}
\label{subsec:rancangan-token-macaroons}



Macaroons digunakan sebagai pengganti token JWT DecMed sesuai dengan hasil analisis solusi pada Subbab \ref{subsec:analisis-pemilihan-token-otorisasi}. Token Macaroon tetap mengikuti model CapBAC karena token berperan sebagai bukti kapabilitas yang dibawa oleh pemegang akses.

Pada rancangan ini, \textit{Server} PRE bertindak sebagai \textit{issuer} token Macaroon awal dan memegang \texttt{RootKey} yang digunakan untuk proses penandatanganan token. \texttt{RootKey} disimpan pada \textit{environment variable} \textit{Server} PRE dan tidak pernah diberikan kepada \textit{client}. Token awal diterbitkan oleh \textit{Server} PRE setelah terdapat persetujuan dari pasien sebagai pemilik data. Dengan demikian, sifat \textit{patient-centric} tetap dipertahankan karena penerbitan akses awal tetap bergantung pada persetujuan pasien, sedangkan otoritas kriptografis penerbitan token berada pada \textit{Server} PRE.Setelah token awal diterbitkan, pemegang token dapat melakukan atenuasi melalui \textit{client} miliknya dan mengirimkan \textit{request} untuk mengatenuasi Macaroons ke \textit{server} PRE.

Karena \textit{caveat} pada Macaroons bersifat \textit{append-only}, \textit{caveat} pada token induk tidak dapat dihapus atau diganti pada token turunan. Oleh karena itu, rancangan \textit{caveat} harus dibuat dalam bentuk batasan yang bersifat monoton. Setiap \textit{caveat} baru hanya dapat mempersempit hak akses dan tidak dapat memperluas hak akses. Dengan demikian, token turunan yang dibuat pada \textit{client} tetap aman karena seluruh batasan pada token induk tetap ikut terbawa dan diverifikasi oleh \textit{server} PRE.

Berdasarkan implementasi DecMed sebelumnya, token memiliki \textit{claim} \texttt{purpose} dan terdapat dua buah token untuk operasi baca atau tulis. Implementasi Move \textit{smart contract}, \textit{client}, dan PRE juga mencatat akses yang telah diberikan kepada tenaga medis untuk 2 \textit{purpose}, yaitu \texttt{Read} dan \texttt{Update}. Pada pengembangan ini, pola 2 buah token masih dipertahankan dengan membedakan nilai \textit{caveat} \texttt{purpose}, yaitu \texttt{Read} atau \texttt{Update}. Hal ini dilakukan sehingga rancangan token Macaroon tetap kompatibel dengan implementasi DecMed.

Pada rancangan ini, tidak ada \textit{caveat} yang digunakan sebagai penanda pemegang token saat ini. Hal ini disebabkan karena nilai pemegang token akan berubah ketika token didelegasikan, sedangkan Macaroon tidak mendukung penghapusan atau penggantian \textit{caveat} lama. Apabila token induk memiliki \texttt{holder\_address = 0xPETUGAS\_REKAM\_MEDIS} dan token turunan menambahkan \texttt{holder\_address = 0xDOKTER}, maka token tersebut akan memiliki dua \textit{caveat} identitas yang saling kontra. Kondisi tersebut dapat menyebabkan token tidak dapat digunakan karena satu permintaan akses tidak mungkin berasal dari dua alamat pemegang yang berbeda.

Sebagai gantinya, rancangan ini menggunakan pendekatan rantai delegasi dengan \textit{caveat} \texttt{root\_subject}, \texttt{delegated\_by}, dan \texttt{delegated\_to}. \texttt{root\_subject} menyatakan aktor penerima token awal dari pasien, sedangkan \texttt{delegated\_by} dan \texttt{delegated\_to} digunakan untuk merepresentasikan proses delegasi pada token turunan. Dengan pendekatan ini, sistem dapat menambahkan informasi delegasi baru tanpa perlu mengubah \textit{caveat} yang sudah ada sebelumnya. Untuk memastikan delegasi berjalan hanya pada satu fasyankes, \textit{caveat} \texttt{hospital\_cid} digunakan untuk mengikat token terhadap ID rumah sakit.

Selain informasi delegasi, token juga memuat batasan akses terhadap data RME tertentu. \textit{Caveat} \texttt{patient\_address} digunakan untuk mengikat token terhadap data RME pasien tertentu. \textit{Caveat} \texttt{related\_rme\_id} digunakan untuk membatasi token terhadap satu episode pelayanan. Kedua \textit{caveat} tersebut memastikan bahwa token tidak dapat digunakan untuk mengakses data pasien lain atau episode pelayanan lain yang tidak termasuk dalam cakupan token.

Namun \textit{caveat} \texttt{related\_rme\_id} hanya berlaku untuk token \texttt{Update} dan tidak berlaku untuk token \texttt{Read}. Hal ini karena pada skenario alur klinis pada pelayanan kesehatan di sebuah fasyankes, seorang dokter pada umumnya membutuhkan akses ke semua data RME pasien untuk menentukan tindakan medis yang tepat. Namun, untuk aksi \texttt{Update} pada data RME, dokter hanya perlu mengakses data RME yang berkaitan dengan episode pelayanan yang sedang berlangsung.

Pada rancangan ini, \texttt{related\_rme\_id} menggunakan format \texttt{RME-\char`\{sequence\_6\_digit\char`\}}, misalnya \texttt{RME-000001}. Format ini mengikuti implementasi DecMed sebelumnya dan dipilih karena sudah cukup untuk kebutuhan prototipe.

Mekanisme pembatasan kategori data dilakukan dengan metode \textit{whitelisting}. Token harus secara eksplisit menyatakan \textit{dataset category} dan \textit{function category} yang dapat diakses oleh pemegang token. Rancangan ini menggunakan \textit{caveat} \texttt{read\_dataset\_in}, \texttt{write\_dataset\_in}, \texttt{read\_function\_in}, dan \texttt{write\_function\_in}. Keempat \textit{caveat} tersebut disesuaikan dengan \textit{metadata} segmen RME yang disimpan pada penyimpanan \textit{on-chain}, yaitu \texttt{dataset\_category} dan \texttt{function\_category}. Dengan demikian, \textit{Server} PRE dapat memvalidasi apakah segmen yang diminta termasuk dalam cakupan akses yang diizinkan oleh token.

Pembatasan waktu akses dilakukan dengan \textit{caveat} \texttt{expires\_before}. \textit{Caveat} ini menyatakan batas waktu maksimum token dapat digunakan. Pada token turunan, nilai \texttt{expires\_before} harus sama dengan atau lebih awal daripada token induk. Pada proses verifikasi, \textit{caveat} \texttt{expires\_before} yang diambil adalah nilai paling awal dari seluruh \textit{caveat} \texttt{expires\_before}.

Pembatasan jumlah delegasi dilakukan dengan \textit{caveat} \texttt{max\_delegation\_depth}. Nilai ini menyatakan jumlah maksimum delegasi yang masih diperbolehkan dari token terkait. Token awal dapat memiliki nilai \texttt{max\_delegation\_depth = 2} yang menandakan bahwa token Macaroon tersebut dapat diturunkan 2 kali. Pada saat membuat token turunan, token turunan dapat ditambahkan \textit{caveat} \texttt{max\_delegation\_depth = 1}, yang berarti kemampuan menurunkan token lebih lanjut telah berkurang 1. \textit{Caveat} ini dapat divalidasi lebih lanjut lagi dengan memanfaatkan \textit{caveat} \texttt{delegated\_by} dan \texttt{delegated\_to} untuk berapa banyak turunan yang telah terjadi.

Untuk membantu verifikasi token yang sudah direvokasi, \textit{caveat} \texttt{parent\_token\_hash} digunakan untuk mengikat token turunan dengan token induk. Nilai \texttt{parent\_token\_hash} adalah \textit{hash} dari token induk. Rincian lebih lanjut tentang mekanisme revokasi token Macaroon dijelaskan pada Subbab \ref{subsubsec:revokasi-token-macaroons}.

Rancangan \textit{caveats} yang digunakan pada token Macaroon ditunjukkan pada Tabel \ref{tab:rancangan-caveats-macaroons-baru}.

\begin{table}[!ht]
	\centering
	\caption{Rancangan \textit{Caveats} pada Token Macaroon}
	\label{tab:rancangan-caveats-macaroons-baru}
	\small
	\begin{tabular}{|p{0.22\textwidth}|p{0.55\textwidth}|p{0.13\textwidth}|}
		\hline
		\textbf{\textit{Caveat}} & \textbf{Fungsi}                                                          & \textbf{Pada Token} \\ \hline

		\texttt{purpose}
		                         & Menyatakan tujuan penggunaan token (\texttt{Read} atau \texttt{Update}). & Keduanya
		\\ \hline

		\texttt{patient\_address}
		                         & Mengikat token terhadap data RME milik pasien tertentu.                  & Keduanya
		\\ \hline

		\texttt{related\_rme\_id}
		                         & Menentukan episode pelayanan atau data RME yang dapat diakses.           & Update
		\\ \hline

		\texttt{root\_subject}
		                         & Menyatakan aktor penerima token awal dari pasien.                        & Keduanya
		\\ \hline

		\texttt{delegated\_by}
		                         & Menyatakan aktor yang melakukan delegasi token.                          & Keduanya
		\\ \hline

		\texttt{delegated\_to}
		                         & Menyatakan aktor penerima token turunan.                                 & Keduanya
		\\ \hline

		\texttt{read\_dataset\_in}
		                         & Membatasi \textit{dataset category} yang dapat dibaca.                   & Read
		\\ \hline

		\texttt{write\_dataset\_\newline in}
		                         & Membatasi \textit{dataset category} yang dapat ditulis.                  & Update
		\\ \hline

		\texttt{read\_function\_\newline in}
		                         & Membatasi \textit{function category} yang dapat dibaca.                  & Read
		\\ \hline

		\texttt{write\_function\_\newline in}
		                         & Membatasi \textit{function category} yang dapat ditulis.                 & Update
		\\ \hline

		\texttt{expires\_before}
		                         & Membatasi masa berlaku token.                                            & Keduanya
		\\ \hline

		\texttt{max\_delegation\_\newline depth}
		                         & Membatasi jumlah delegasi turunan yang masih diperbolehkan.              & Keduanya
		\\ \hline

		\texttt{hospital\_cid}
		                         & Mengikat token dengan rumah sakit tertentu.                              & Keduanya
		\\ \hline

		\texttt{parent\_token\_\newline hash}
		                         & Mengikat token turunan dengan token induk untuk proses revokasi.         & Keduanya
		\\ \hline
	\end{tabular}
\end{table}

Sebagai contoh, pasien dapat memberikan token awal kepada petugas rekam medis dengan cakupan akses yang cukup untuk mengoordinasikan pelayanan pada satu episode RME tertentu. Petugas rekam medis diberikan token awal dengan cakupan yang luas karena berperan sebagai koordinator awal pelayanan. Meskipun cakupannya luas, namun aksi yang dilakukan oleh petugas rekam medis dapat dibatasi dengan penegakan otoritas berbasis \textit{role} yang akan dijelaskan pada Subbab \ref{subsubsec:verifikasi-role-tenaga-medis-untuk-akses-dataset}. Contoh \textit{caveat} token awal untuk petugas rekam medis disajikan pada Kode \ref{lst:contoh-token-baca-awal-rekam-medis} dan Kode \ref{lst:contoh-token-pembaruan-awal-rekam-medis}.

Pada saat token turunan digunakan oleh tenaga medis, \textit{Server} PRE akan memverifikasi seluruh \textit{caveat} yang terdapat pada token, baik \textit{caveat} dari token awal maupun \textit{caveat} tambahan dari token turunan. Mekanisme verifikasi token pada \textit{Server} PRE dijelaskan pada Subbab \ref{subsubsec:verifikasi-caveats-macaroons-di-pre}.

Rancangan ini juga mendukung prinsip \textit{least privilege}. Dokter dapat memiliki cakupan akses yang relatif luas untuk menjalankan tugas klinisnya, tetapi token turunan yang diberikan kepada petugas laboratorium atau apoteker dapat dibatasi hanya pada segmen yang relevan. Sebagai contoh, petugas laboratorium hanya dapat membaca segmen \texttt{PEMERIKSAAN\_PENUNJANG} pada \textit{dataset} \texttt{RAWAT\_JALAN} dan menulis segmen \texttt{LABORATORIUM} pada \textit{dataset} \texttt{LABORATORIUM}. Petugas tersebut tidak dapat membaca diagnosis, anamnesis, resep, atau segmen klinis lain yang tidak relevan dengan tugasnya. \subsection{Rancangan Penegakan Otoritas}
\label{subsec:rancangan-penegakan-otoritas}



Penegakan otoritas pada rancangan ini merujuk pada proses memastikan \textit{request} akses terhadap data RME hanya dilayani apabila aktor yang mengajukan permintaan memiliki kewenangan yang sah.

Bagian ini membahas tiga aspek utama penegakan otoritas pada \textit{Server} PRE. Pertama, Subbab \ref{subsubsec:verifikasi-caveats-macaroons-di-pre} menjelaskan verifikasi \textit{caveats} Macaroon untuk memastikan token masih valid. Kedua, Subbab \ref{subsubsec:verifikasi-role-tenaga-medis-untuk-akses-dataset} membahas validasi tambahan terhadap \textit{subrole} \texttt{MedicalPersonnel} agar akses tulis hanya dapat dilakukan pada \textit{dataset category} yang sesuai dengan kewenangannya. Ketiga, Subbab \ref{subsubsec:pencabutan-hak-akses-macaroon-delegasi} menjelaskan rancangan mekanisme revokasi token Macaroon, termasuk revokasi agar token turunannya tidak dapat digunakan.

Sebagai gambaran, Gambar \ref{fig:seqdiag-overview-penegakan-otoritas} menunjukkan alur penegakan otoritas pada \textit{Server} PRE. Secara garis besar, alur akses dimulai dengan aktor melalui \textit{client} untuk mengirimkan permintaan akses ke \textit{Server} PRE. \textit{Server} PRE kemudian memverifikasi token Macaroon, memverifikasi status revokasi dari token, dan menentukan aktor aktif. Jika semua verifikasi berhasil, \textit{Server} PRE melanjutkan proses ke 3 kemungkinan skenario, yaitu mengambil daftar RME yang dapat diakses, melakukan pembacaan data RME yang melibatkan proses \textit{re-encryption}, dan melakukan penulisan data RME.

\begin{figure}[!ht]
	\centering
	\includegraphics[width=\textwidth]{resources/chapter-3/rancangan/penegakan otoritas/seqdiag-overview.png}
	\caption{Alur penegakan otoritas pada \textit{Server} PRE}
	\label{fig:seqdiag-overview-penegakan-otoritas}
\end{figure}

\subsubsection{Verifikasi \textit{Caveats} Macaroons}
\label{subsubsec:verifikasi-caveats-macaroons-di-pre}

\textit{Server} PRE bertindak sebagai titik penegakan kebijakan akses. \textit{Server} PRE memverifikasi token Macaroon dan seluruh \textit{caveats}. Jika terdapat \textit{caveat} yang tidak dikenali atau diluar rancangan pada Tabel \ref{tab:rancangan-caveats-macaroons-baru}, permintaan akses ditolak.

Rancangan langkah verifikasi token Macaroon untuk permintaan akses baca data RME adalah sebagai berikut:

\begin{enumerate}[leftmargin=*]
	\item \textit{Server} PRE memverifikasi validitas token Macaroon melalui mekanisme \textit{Chained-HMAC} menggunakan \texttt{RootKey}.
	\item \textit{Server} PRE memverifikasi bahwa \texttt{patient\_address} dan \texttt{related\_rme\_id} pada token sesuai dengan data RME yang diminta.
	\item \textit{Server} PRE menghitung aktor aktif. Jika tidak terdapat \texttt{delegated\_to}, maka aktor aktif adalah \texttt{root\_subject}. Jika terdapat satu atau lebih \texttt{delegated\_to}, maka aktor aktif adalah nilai \texttt{delegated\_to} terakhir.
	\item \textit{Server} PRE memverifikasi tanda tangan \textit{wallet} dari aktor aktif pada permintaan akses.
	\item \textit{Server} PRE menghitung kedalaman delegasi berdasarkan jumlah \texttt{delegated\_to} pada token dan memastikan nilainya tidak melebihi \texttt{max\_delegation\_depth}.
	\item \textit{Server} PRE memvalidasi seluruh batas waktu pada \texttt{expires\_before} dan menggunakan batas waktu yang paling awal sebagai masa berlaku efektif token.
	\item \textit{Server} PRE menghitung cakupan akses efektif berdasarkan irisan seluruh \texttt{read\_dataset\_in}, \texttt{write\_dataset\_in}, \texttt{read\_function\_in}, dan \texttt{write\_function\_in}.
	\item \textit{Server} PRE mengambil metadata segmen dari IOTA dan mencocokkan \texttt{dataset\_category} serta \texttt{function\_category} segmen dengan cakupan akses efektif token.
	\item Jika seluruh verifikasi berhasil, \textit{Server} PRE melanjutkan proses pengambilan data terenkripsi dari IPFS dan melakukan \textit{re-encryption}. Jika salah satu pemeriksaan gagal, permintaan akses ditolak.
\end{enumerate}

\subsubsection{Verifikasi \textit{Role} dan \textit{Subrole} Tenaga Medis untuk Akses \textit{Dataset}}
\label{subsubsec:verifikasi-role-tenaga-medis-untuk-akses-dataset}

Dari hasil analisis pada Subbab \ref{subsec:analisis-verifikasi-aktor-operasi-tulis}, diperlukan verifikasi tambahan terhadap \textit{subrole} \texttt{MedicalPersonnel} agar akses tulis hanya dapat dilakukan pada \textit{dataset category} yang sesuai dengan kewenangannya. Pada implementasi Tugas Akhir ini, terdapat 4 \textit{subrole} turunan dari \textit{role} \texttt{MedicalPersonnel}, yaitu dokter, perawat, petugas laboratorium, dan apoteker. Sebagaimana dijelaskan pada Subbab \ref{subsec:analisis-segmentasi-data-rme}, setiap \textit{subrole} tersebut memiliki kebutuhan akses data yang berbeda-beda untuk melakukan akses jenis \textit{update}. Pemetaan kebutuhan akses data untuk setiap \textit{subrole} dapat dilihat pada Tabel \ref{tab:pemetaan-hak-akses-update-subrole}.

\begin{table}[!ht]
	\centering
	\caption{Pemetaan Hak Akses Update Berdasarkan subrole Tenaga Medis}
	\label{tab:pemetaan-hak-akses-update-subrole}
	\begin{tabular}{|p{0.15\textwidth}|p{0.35\textwidth}|p{0.35\textwidth}|}
		\hline
		\textbf{\textit{subrole}} & \textbf{\textit{Dataset Category}} & \textbf{\textit{Function Category}}                  \\
		\hline
		Perawat                   & Rawat Jalan/Rawat Inap             & Anamnesis, Pemeriksaan Fisik, Pemeriksaan Psikologis \\
		\hline
		Dokter                    & Rawat Jalan/Rawat Inap             & Seluruh \textit{function category}                   \\
		\hline
		Lab                       & Laboratorium                       & Laboratorium                                         \\
		\hline
		Apoteker                  & Apotek                             & Peresepan dan Dispensing                             \\
		\hline
	\end{tabular}
\end{table}

Berdasarkan pemetaan tersebut, terdapat verifikasi tambahan untuk \textit{role} \texttt{MedicalPersonnel}, yaitu melakukan validasi apakah \textit{subrole} pengguna yang saat ini mencoba melakukan aksi \textit{write} memang sesuai dengan jenis \textit{dataset} yang diakses. Alur tambahan tersebut ditunjukkan pada Gambar \ref{fig:seqdiag-validasi-subrole-dataset}. Pada alur ini, \textit{Server} PRE terlebih dahulu mendapatkan informasi \textit{subrole} tenaga medis melalui \textit{smart contract}. Informasi tersebut kemudian digunakan untuk memvalidasi pasangan \textit{role} dan \textit{subrole} terhadap kategori \textit{dataset} yang coba ditulis. Apabila \textit{subrole} aktor tidak memiliki kewenangan untuk menulis pada kategori \textit{dataset} tersebut, maka \textit{request} ditolak. Bagian berlatar merah pada Gambar \ref{fig:seqdiag-validasi-subrole-dataset} menunjukkan alur yang berubah atau ditambahkan pada rancangan Tugas Akhir ini dibandingkan Implementasi DecMed sebelumnya. Perubahan tersebut mencakup pengambilan informasi \textit{subrole} dari \textit{smart contract} serta validasi kesesuaian \textit{role}, \textit{subrole}, dan kategori \textit{dataset} sebelum akses tulis diberikan.

\begin{figure}[!ht]
	\centering
	\includegraphics[width=0.95\textwidth,height=0.80\textheight,keepaspectratio]{resources/chapter-3/rancangan/seqdiag-validasi-subrole X dataset.png}
	\caption{\textit{Sequence Diagram} Rancangan Validasi \textit{subrole} Tenaga Medis terhadap \textit{Dataset Category}}
	\label{fig:seqdiag-validasi-subrole-dataset}
\end{figure}

Validasi menggunakan data \textit{role} dan \textit{subrole} yang terdapat pada penyimpanan\textit{on-chain} karena pada Macaroon tidak terdapat informasi \textit{role} dari pengguna aktif saat ini, sesuai dengan rancangan \textit{caveats} pada Tabel \ref{tab:rancangan-caveats-macaroons-baru}. Sementara itu, \textit{registry personnel} pada penyimpanan \textit{on-chain} merepresentasikan status kelembagaan aktor, misalnya apakah aktor masih terdaftar sebagai \textit{personnel} rumah sakit dengan \textit{role} atau \textit{subrole} tertentu. Dengan demikian, \textit{request} hanya dilanjutkan apabila aktor memiliki token yang valid sekaligus masih memiliki otoritas kelembagaan yang sesuai.

\subsubsection{Verifikasi \textit{Wallet} Aktor Aktif}
\label{subsubsec:verifikasi-wallet-aktor-aktif}

Sesuai analisis pada Subbab \ref{subsec:verifikasi-bukti-kepemilikan-wallet}, diperlukan verifikasi tambahan terhadap \textit{wallet} aktor aktif untuk memastikan \textit{request} dilakukan oleh aktor aktif yang sah.

Rancangan proses verifikasi dilakukan dengan mekanisme \textit{presigned request}. Ketika \textit{client} mengirimkan \textit{request} ke \textit{Server} PRE, \textit{client} juga menambahkan header yang berisi konteks \textit{request} yang telah di-\textit{sign} menggunakan kunci privat \textit{wallet} pengguna. Konteks \textit{request} memuat informasi seperti \texttt{token\_id}, \texttt{patient\_address}, \texttt{related\_rme\_id}, \texttt{operation}, \texttt{segment\_id}, \texttt{dataset\_category}, \texttt{function\_category}, dan \texttt{timestamp}. Keberadaan \textit{timestamp} pada konteks yang ditandatangani juga membuat \textit{signature} hanya berlaku untuk rentang waktu tertentu. Apabila \textit{signature} tidak valid maka \textit{request} ditolak.

\subsubsection{Pencabutan Hak Akses Macaroon Delegasi}
\label{subsubsec:pencabutan-hak-akses-macaroon-delegasi}

Sesuai dengan analisis pada Subbab \ref{subsec:analisis-revokasi-audit-delegasi}, diperlukan mekanisme revokasi dan audit delegasi untuk memastikan bahwa pencabutan akses dapat berdampak pada seluruh token turunan yang berasal dari token tersebut. Revokasi akses pada implementasi Tugas Akhir ini masih menjadikan pasien sebagai \textit{resource owner} yang berwenang untuk mencabut akses terhadap RME.

Revokasi akses pada implementasi ini dilakukan melalui kombinasi pencatatan status di IOTA \textit{on-chain} dan Redis PRE, dengan \textit{server} PRE bertindak sebagai titik penegakan kebijakan akses. Lapisan \textit{on-chain} digunakan untuk mencatat perubahan status akses dan audit trail revokasi, sedangkan Redis PRE digunakan untuk menyimpan \textit{hash} dari token yang telah dicabut. Alur pencabutan akses oleh pasien ditunjukkan pada Gambar \ref{fig:seqdiag-revoke-access}. Pada alur tersebut, pasien memilih akses yang akan dicabut melalui \textit{Patient Client}. \textit{Patient Client} kemudian mengeksekusi fungsi \texttt{revoke\_access()} pada \textit{smart contract} Move untuk menandai akses sebagai dicabut, menghapus hak akses pada penyimpanan \textit{on-chain}, melakukan pencabutan berantai terhadap penerima delegasi turunan, serta mencatat kejadian revokasi ke dalam \textit{audit log}. Setelah transaksi \textit{on-chain} berhasil, \textit{Patient Client} mengirimkan \textit{payload} revokasi ke \textit{Server} PRE agar menyimpan \textit{hash} dari token yang telah dicabut pada Redis PRE dan menghapus \textit{re-encryption key} terkait dengan token Macaroon yang direvokasi.

\begin{figure}[!ht]
	\centering
	\includegraphics[width=0.95\textwidth,height=0.80\textheight,keepaspectratio]{resources/chapter-3/rancangan/seqdiag-revoke_access.png}
	\caption{\textit{Sequence Diagram} Rancangan Alur Pencabutan Hak Akses Macaroon Delegasi}
	\label{fig:seqdiag-revoke-access}
\end{figure}

Bagian berlatar merah pada Gambar \ref{fig:seqdiag-revoke-access} menunjukkan alur baru yang ditambahkan pada Tugas Akhir ini. Alur baru tersebut mencakup pencabutan berantai terhadap seluruh \textit{downstream delegatees}, pengiriman \textit{payload} revokasi ke \textit{Server} PRE, pembentukan \textit{revocation key}, penyimpanan status revokasi pada Redis dengan TTL sesuai masa berlaku token, serta penghapusan kunci \textit{re-encryption}. Dengan penambahan alur ini, revokasi tidak hanya tercatat pada IOTA, tetapi juga langsung ditegakkan oleh \textit{Server} PRE ketika token Macaroon digunakan kembali.

Pada saat verifikasi, \textit{Server} PRE memeriksa \textit{hash} token yang sedang digunakan dan setiap \textit{caveat} \texttt{parent\_token\_hash} ke \textit{revocation key} pada Redis.  Dengan demikian, apabila pasien mencabut akses pada delegator awal, maka token Macaroon turunan yang dibuat dari akses tersebut juga akan dianggap tidak valid. Selain itu, pemerikskaan status hak akses juga dilakukan pada penyimpanan \textit{on-chain} yang mencatat status hak akses. \subsection{Rancangan Alur Akses dan Delegasi Macaroon}
\label{subsec:rancangan-alur-akses-dan-delegasi-macaroon}



Sesuai dengan analisis pada Subbab \ref{subsec:analisis-delegasi-atenuasi-offline} dan skenario pada Subbab \ref{subsubsec:analisis-skenario-layanan-klinis}, rancangan alur akses dan delegasi pada Tugas Akhir ini dirancang untuk menjawab kebutuhan delegasi akses antar aktor layanan kesehatan tanpa mengharuskan pasien menerbitkan ulang token setiap kali terjadi perpindahan tugas dalam proses pelayanan.

Alur akses dan delegasi pada rancangan ini mengikuti alur pelayanan kesehatan yang umum terjadi pada proses rawat jalan dan rawat inap. Alur ini dibahas pada Subbab \ref{subsubsec:analisis-skenario-layanan-klinis}. Pasien terlebih dahulu diterima oleh petugas rekam medis. Setelah itu, petugas rekam medis memberikan akses kepada perawat untuk melakukan pemeriksaan awal. Petugas rekam medis juga memberikan akses kepada dokter terkait agar dokter dapat memeriksa pasien dan menentukan tindakan medis berikutnya. Apabila dokter membutuhkan pemeriksaan penunjang, dokter dapat mendelegasikan akses kepada petugas laboratorium. Setelah hasil pemeriksaan tersedia dan dokter menentukan terapi atau pengobatan, dokter dapat mendelegasikan akses kepada apoteker untuk memproses resep dan penyerahan obat. Rancangan alur tersebut ditunjukkan pada Gambar \ref{fig:alur-akses-delegasi}.

\begin{figure}[!ht]
	\centering
	\includegraphics[width=0.95\textwidth,height=0.80\textheight,keepaspectratio]{resources/chapter-3/analisis/flowdiagram-alur-akses.png}
	\caption{\textit{Flow Diagram} Alur Akses dan Delegasi pada Proses Pelayanan Pasien}
	\label{fig:alur-akses-delegasi}
\end{figure}

Delegasi akses pada rancangan ini terdiri atas tiga tahap utama. Tahap pertama adalah penerbitan token awal kepada petugas rekam medis sebagai aktor yang menerima pasien dan mengoordinasikan akses awal dalam proses pelayanan. Tahap kedua adalah delegasi lokal dari petugas rekam medis kepada perawat dan dokter. Tahap ketiga adalah delegasi lanjutan dari dokter kepada petugas laboratorium dan apoteker, sesuai kebutuhan tindakan medis.

Penjelasan kemudian dilanjutkan dengan bagian berikutnya, yaitu pola teknis proses pemberian token Macaroon atau proses delegasi oleh \textit{delegator} hingga akhirnya digunakan oleh \textit{delegatee} untuk mengakses data RME. Pola teknis ini dibagi menjadi dua bagian, yaitu proses pemberian token Macaroon delegasi hingga tercatat pada \textit{on-chain} dan proses pengambilan dan penggunaan token Macaroon delegasi.

\subsubsection{Penerbitan Token Awal oleh Pasien kepada Petugas Rekam Medis}
\label{subsubsec:penerbitan-token-awal-oleh-pasien-kepada-petugas-rekam-medis}

Penerbitan token awal merupakan proses pemberian otorisasi pertama pada saat pasien memulai layanan kesehatan. Pada tahap ini, pasien tetap menjadi pihak yang mengendalikan pemberian akses awal terhadap data RME miliknya. Pasien memberikan persetujuan akses melalui \textit{Patient Client}, kemudian \textit{Server} PRE menerbitkan token awal berdasarkan persetujuan tersebut. Pasien mengetahui identitas dari petugas rekam medis melalui QR yang dipindai, kemudian membuktikan  kepemilikan \textit{wallet} kepada PRE dengan \textit{nonce} sebelum token awal diterbitkan.

Pada rancangan ini, petugas rekam medis berperan sebagai aktor awal yang menerima pasien dan menginisiasi proses pelayanan. Token ini tidak dimaksudkan agar petugas rekam medis dapat mengakses seluruh data klinis pasien secara bebas, melainkan sebagai token induk yang dapat digunakan untuk menurunkan akses kepada aktor layanan kesehatan yang relevan. Terdapat 2 token Macaroon yang diterbitkan pada tahapan ini, yaitu token \texttt{Read} dan token \texttt{Update}. Penggunaan langsung token tersebut tetap dibatasi oleh pemeriksaan otoritas yang dibahas pada Subbab \ref{subsec:rancangan-penegakan-otoritas}, sehingga cakupan token induk tidak otomatis berarti seluruh data dapat dibaca bebas oleh petugas rekam medis. Contoh \textit{caveat} yang ditambahkan pada token awal untuk petugas rekam medis dapat dilihat pada Lampiran \ref{lst:contoh-token-baca-awal-rekam-medis} dan \ref{lst:contoh-token-pembaruan-awal-rekam-medis}.


Untuk memastikan RME tetap dapat diproses dengan mekanisme Umbral PRE, dibutuhkan sebuah \textit{data PRE seed}. \textit{Data PRE seed} ini digunakan untuk membangkitkan pasangan kunci PRE data yang dibutuhkan dalam proses enkripsi, dekripsi, dan \textit{re-encryption}. \textit{Data PRE seed} tersebut dibuat pada \textit{patient client} ketika akses awal diberikan, lalu disimpan dalam metadata akses di \textit{on-chain} dalam kondisi terenkripsi menggunakan kunci publik PRE penerima akses. Pada proses delegasi, \textit{data PRE seed} yang sudah ada didekripsi oleh pemberi delegasi, kemudian dienkripsi ulang menggunakan kunci publik PRE penerima delegasi dan disimpan kembali bersama metadata delegasi di \textit{on-chain}.

Alur penerbitan token awal kepada petugas rekam medis diawali dengan pasien bertemu dengan petugas rekam medis yang menerima permintaan pelayanan serta memperlihatkan QR untuk discan oleh pasien. Pasien memberikan persetujuan akses awal melalui \textit{Patient Client} setelah memilih konteks pelayanan (rawat jalan atau rawat inap), batasan waktu akses, dan jenis akses yang diberikan. \textit{Patient Client} mengirimkan permintaan penerbitan token kepada \textit{Server} PRE. \textit{Server} PRE memverifikasi identitas pasien, identitas petugas rekam medis, serta batasan yang diberikan. Jika lolos proses verifikasi, \textit{Server} PRE menerbitkan token Macaroon awal dengan \textit{caveats} yang mengikat token pada pasien, ID RME terkait, petugas rekam medis sebagai penerima awal, cakupan akses, masa berlaku, serta batas delegasi. \textit{Server} PRE mencatat penerbitan token awal ke dalam penyimpanan \textit{on-chain}. Token awal didistribusikan kepada petugas rekam medis melalui mekanisme yang dijelaskan pada Subbab \ref{subsubsec:proses-pemberian-dan-penggunaan-macaroon-delegasi} dengan memanfaatkan penyimpanan \textit{on-chain}.

\subsubsection{Delegasi dari Petugas Rekam Medis kepada Perawat dan Dokter}
\label{subsec:delegasi-dari-petugas-rekam-medis-kepada-perawat-dan-dokter}

Setelah menerima token awal, petugas rekam medis dapat mendelegasikan akses kepada perawat dan dokter terkait sesuai kebutuhan pelayanan. Delegasi ini dilakukan dengan mengirimkan permintaan atenuasi kepada \textit{server} PRE. \textit{Server} PRE memverifikasi \textit{request} delegasi. Apabila valid, \textit{server} PRE mengembalikan token turunan untuk diproses lebih lanjut oleh \textit{client delegator} untuk dicatat pada penyimpanan \textit{on-chain}.

Delegasi kepada perawat diberikan agar perawat dapat membaca RME pasien untuk melakukan penanganan, dan menulis hasil pemeriksaan awal. Sementara itu, delegasi kepada dokter diberikan agar dokter dapat membaca data seluruh rekam medis pasien, menuliskan RME untuk episode layanan saat ini, serta melakukan delegasi lanjutan apabila diperlukan. Pemetaan delegasi dari petugas rekam medis ditunjukkan pada Tabel \ref{tab:pemetaan-delegasi-rekam-medis}. Contoh \textit{caveats} yang ditambahkan pada token turunan untuk perawat dan dokter dapat dilihat pada Lampiran \ref{lst:contoh-token-delegasi-dokter} dan \ref{lst:contoh-token-delegasi-perawat}.

\begin{table}[!ht]
	\centering
	\caption{Pemetaan Delegasi Token dari Petugas Rekam Medis}
	\label{tab:pemetaan-delegasi-rekam-medis}
	\small
	\begin{tabular}{|p{0.22\textwidth}|p{0.28\textwidth}|p{0.4\textwidth}|}
		\hline
		\textbf{Penerima Delegasi} & \textbf{Tujuan Delegasi}                                                                                                                                                              & \textbf{Cakupan Akses} \\ \hline

		Perawat
		                           & Melakukan anamnesis dan pemeriksaan awal pasien.
		                           & Menulis segmen anamnesis, pemeriksaan fisik, dan pemeriksaan psikologis pada \textit{dataset} rawat jalan.                                                                                                     \\ \hline

		Dokter
		                           & Memeriksa pasien, membaca hasil pemeriksaan awal, menentukan tindakan medis, diagnosis, terapi, dan kebutuhan pemeriksaan penunjang.
		                           & Membaca dan menulis segmen klinis yang relevan pada \textit{dataset} rawat jalan atau rawat inap, serta memiliki kemampuan delegasi lanjutan untuk apoteker dan petugas laboratorium.                          \\ \hline
	\end{tabular}
\end{table}

\subsubsection{Delegasi dari Dokter kepada Laboratorium dan Apoteker}
\label{subsec:delegasi-dari-dokter-kepada-laboratorium-dan-apoteker}

Setelah dokter menerima token turunan dari petugas rekam medis, dokter dapat melakukan delegasi lanjutan apabila dibutuhkan dalam proses pelayanan. Delegasi lanjutan dilakukan kepada petugas laboratorium apabila dokter membutuhkan pemeriksaan penunjang. Setelah dokter menentukan terapi, dokter juga dapat mendelegasikan akses kepada apoteker untuk memproses resep dan penyerahan obat.  Contoh \textit{caveats} yang ditambahkan pada token turunan untuk laboratorium dan apoteker dapat dilihat pada Lampiran \ref{lst:contoh-token-delegasi-lab} dan \ref{lst:contoh-token-delegasi-apoteker}. Pemetaan delegasi lanjutan dari dokter ditunjukkan pada Tabel \ref{tab:pemetaan-delegasi-dokter-lanjutan}.

\begin{table}[!ht]
	\centering
	\caption{Pemetaan Delegasi Lanjutan dari Dokter}
	\label{tab:pemetaan-delegasi-dokter-lanjutan}
	\small
	\begin{tabular}{|p{0.22\textwidth}|p{0.28\textwidth}|p{0.4\textwidth}|}
		\hline
		\textbf{Penerima Delegasi} & \textbf{Tujuan Delegasi}                                                                        & \textbf{Cakupan Akses} \\ \hline

		Petugas Laboratorium
		                           & Melakukan pemeriksaan penunjang berdasarkan permintaan dokter dan mengunggah hasil pemeriksaan.
		                           & Membaca segmen pemeriksaan penunjang dan menulis pada \textit{dataset} laboratorium.                                     \\ \hline

		Apoteker
		                           & Memproses resep, menyiapkan obat, dan mencatat penyerahan obat kepada pasien.
		                           & Membaca segmen terapi dan menulis pada \textit{dataset} apotek.                                                          \\ \hline
	\end{tabular}
\end{table}

\subsubsection{Proses Pemberian dan Penggunaan Token Macaroon Delegasi}
\label{subsubsec:proses-pemberian-dan-penggunaan-macaroon-delegasi}

Secara umum, proses pemberian dan penggunaan token Macaroon delegasi dibagi menjadi dua bagian. Bagian pertama adalah proses pembentukan dan pencatatan delegasi, \textit{client} yang meminta atenuasi token ke \textit{server} PRE. Bagian ini mencakup pembentukan token Macaroon turunan, enkripsi metadata delegasi, serta pencatatan metadata akses terenkripsi pada penyimpanan \textit{on-chain}. Alur proses tersebut ditunjukkan pada Gambar \ref{fig:seq-pemberian-macaroon-delegasi}.

\begin{figure}[htbp]
	\centering
	\includegraphics[width=0.92\textwidth]{resources/chapter-3/rancangan/seqdiag-delegasi-tahap1-3.png}
	\caption{Sequence diagram pembentukan dan pencatatan token Macaroon delegasi}
	\label{fig:seq-pemberian-macaroon-delegasi}
\end{figure}

Berdasarkan Gambar  \ref{fig:seq-pemberian-macaroon-delegasi}, proses delegasi diawali ketika \textit{delegator} memilih cakupan akses yang akan diberikan kepada \textit{delegatee}. Setelah cakupan akses ditentukan, \textit{client delegator} mengambil token Macaroon yang dimiliki oleh \textit{delegator} dan melakukan \textit{request} atenuasi token ke \textit{server} PRE.

\textit{Server} PRE memverifikasi \textit{request} delegasi. Validasi dilakukan untuk memeriksa apakah cakupan akses turunan lebih luas, verifikasi \textit{signature} dari \textit{delegator}, status akses \textit{delegator} pada \textit{on-chain}, serta batasan cakupan yang diminta. Apabila valid, \textit{server} PRE mengembalikan token turunan untuk diproses lebih lanjut oleh \textit{client delegator} untuk dicatat pada penyimpanan \textit{on-chain}.

Setelah token Macaroon turunan terbentuk, token tersebut tidak disimpan secara langsung dalam bentuk \textit{plaintext} pada penyimpanan \textit{on-chain}. Aplikasi \textit{client} milik \textit{delegator} terlebih dahulu mengambil kunci publik milik \textit{delegatee} dari IOTA, termasuk kunci publik PRE milik \textit{delegatee}. Kunci publik PRE kemudian digunakan untuk mengenkripsi objek akses yang berisi token Macaroon turunan dan metadata delegasi. Dengan pendekatan ini, metadata akses yang dicatat pada \textit{on-chain} hanya tersedia dalam bentuk terenkripsi yang hanya bisa dibuka oleh \textit{delegatee}.

Metadata akses terenkripsi kemudian dikirimkan ke disimpan pada IOTA menggunakan \textit{smart contract}. Pada tahap ini, \textit{smart contract} mencatat metadata terenkripsi pada data akses milik \textit{delegatee} dan juga memvalidasi prasyarat delegasi yang dapat diperiksa dari penyimpanan \textit{on-chain}. Prasyarat tersebut meliputi status aktif akun, kesesuaian fasyankes antara \textit{delegator} dan \textit{delegatee}, serta keberadaan dan masa berlaku akses \textit{delegator} terhadap pasien terkait. Apabila prasyarat tersebut valid, metadata akses terenkripsi disimpan pada daftar akses milik \textit{delegatee}. Apabila prasyarat tidak terpenuhi, transaksi delegasi ditolak.

Bagian kedua dari proses ini adalah pengambilan dan penggunaan token Macaroon delegasi oleh \textit{delegatee}. Bagian ini terjadi setelah metadata akses terenkripsi berhasil dicatat pada blockchain. Alur pengambilan dan penggunaan token Macaroon delegasi ditunjukkan pada Gambar~\ref{fig:seq-penggunaan-macaroon-delegasi}.

\begin{figure}[htbp]
	\centering
	\includegraphics[width=\textwidth]{resources/chapter-3/rancangan/seqdiag-delegasi-tahap4-5.png}
	\caption{Sequence diagram pengambilan dan penggunaan token Macaroon delegasi}
	\label{fig:seq-penggunaan-macaroon-delegasi}
\end{figure}

Berdasarkan Gambar \ref{fig:seq-penggunaan-macaroon-delegasi}, ketika \textit{delegatee} membuka aplikasi, \textit{client} mengambil metadata akses terenkripsi yang tersimpan pada blockchain. Metadata tersebut kemudian didekripsi secara lokal menggunakan kunci privat PRE milik \textit{delegatee}. Apabila proses dekripsi berhasil, aplikasi memperoleh token Macaroon turunan beserta metadata pendukungnya. Token Macaroon turunan tersebut kemudian dapat digunakan sebagai token otorisasi ketika \textit{delegatee} mengirim permintaan akses baca atau tulis ke \textit{Server} PRE. Selain membentuk token turunan, aplikasi \textit{client} juga membentuk bukti delegasi yang di-\textit{sign} menggunakan \textit{wallet} milik \textit{delegator}. Bukti ini digunakan agar \textit{Server} PRE dapat memastikan bahwa token turunan memang dibuat oleh pihak yang berwenang dalam rantai delegasi.

Pada tahap penggunaan token, \textit{Server} PRE melakukan verifikasi terhadap token Macaroon yang dikirimkan oleh \textit{delegatee}. Verifikasi ini mencakup pemeriksaan validitas token melalui mekanisme \textit{chained-HMAC}, pemeriksaan status revokasi token, verifikasi bukti kepemilikan \textit{wallet} serta bukti delegasi, dan pemeriksaan seluruh \textit{caveats} yang melekat pada token. Selain itu, \textit{Server} PRE juga dapat mengambil informasi aktor aktif dari \textit{on-chain}, seperti \textit{role} dan status akun, untuk memastikan permintaan akses berasal dari pihak yang sesuai. \textit{Caveats} tersebut diverifikasi terhadap konteks permintaan akses, seperti identitas pasien, RME terkait, jenis hak akses, dataset, fungsi klinis, waktu kedaluwarsa, dan identitas \textit{delegatee}. Apabila seluruh \textit{caveats} terpenuhi, permintaan akses diterima sesuai batas delegasi yang diberikan. Sebaliknya, apabila terdapat \textit{caveats} yang tidak terpenuhi, \textit{Server} PRE menolak permintaan akses tersebut.

Dengan mekanisme ini, IOTA digunakan sebagai media pencatatan dan distribusi metadata akses yang dapat diaudit, sedangkan kerahasiaan token tetap dilindungi melalui enkripsi. Token Macaroon tidak tersedia dalam bentuk terbuka pada \textit{state on-chain}, sehingga pihak selain \textit{delegatee} tidak dapat langsung memanfaatkan metadata akses yang tercatat pada blockchain. Selain itu, karena pembatasan hak akses dibawa langsung oleh \textit{caveats} pada Macaroon, \textit{server} PRE tetap dapat menegakkan batas delegasi pada saat token digunakan.

Untuk mengurangi risiko \textit{delegator} melakukan delegasi akses berlebihan, jumlah delegasi yang dapat dilakukan oleh seorang pemegang hak akses dapat dibatasi dengan mencatat jumlah delegasi yang telah dilakukan oleh \textit{delegator} pada penyimpanan \textit{on-chain}. Jumlah delegasi ini dapat diperiksa oleh \textit{server} PRE saat memverifikasi \textit{request} delegasi. Apabila jumlah delegasi melebihi batas yang ditentukan, \textit{server} PRE menolak \textit{request} delegasi. Jumlah delegasi ini dapat dibatasi dengan dengan aturan bahwa seorang \textit{delegator} hanya dapat melakukan delegasi sebanyak 1 kali untuk setiap \textit{role}. \subsection{Pemetaan Rancangan terhadap Solusi}
\label{subsec:pemetaan-rancangan-solusi}

Setelah rancangan solusi dijelaskan, setiap rancangan dipetakan kembali terhadap solusi yang telah dianalisis pada Subbab \ref{subsec:pemetaan-kebutuhan-solusi}. Pemetaan ini digunakan untuk menunjukkan bahwa setiap rancangan \texttt{R-TA} memiliki keterkaitan langsung dengan solusi \texttt{S-TA} yang diusulkan.

\begin{table}[!ht]
	\centering
	\caption{Pemetaan Rancangan terhadap Solusi}
	\label{tab:daftar-rancangan-solusi}
	\small
	\begin{tabular}{|p{0.12\textwidth}|p{0.54\textwidth}|p{0.24\textwidth}|}
		\hline
		\textbf{ID}      & \textbf{Nama Rancangan}                                                                                                                                 & \textbf{Merealisasikan Solusi}                                         \\ \hline
		\texttt{R-TA-01} & Rancangan arsitektur solusi DecMed dengan Macaroons, segmentasi, dan PRE sebagai \textit{policy enforcement point}.                                     & \texttt{S-TA-01}, \texttt{S-TA-02}, \texttt{S-TA-03}, \texttt{S-TA-04} \\ \hline
		\texttt{R-TA-02} & Rancangan segmentasi data RME berdasarkan \textit{dataset} dan \textit{function category}.                                                              & \texttt{S-TA-02}                                                       \\ \hline
		\texttt{R-TA-03} & Rancangan \textit{caveats} Macaroons.                                                                                                                   & \texttt{S-TA-01}                                                       \\ \hline
		\texttt{R-TA-04} & Rancangan rantai delegasi Macaroons dengan \texttt{root\_subject}, \texttt{delegated\_by}, \texttt{delegated\_to}, dan \texttt{max\_delegation\_depth}. & \texttt{S-TA-03}, \texttt{S-TA-04}                                     \\ \hline
		\texttt{R-TA-05} & Rancangan verifikasi \textit{caveats}, \textit{effective capability}, dan bukti kepemilikan \textit{wallet} pada \textit{Server} PRE.                   & \texttt{S-TA-04}                                                       \\ \hline
		\texttt{R-TA-06} & Rancangan validasi \textit{subrole} tenaga medis terhadap \textit{dataset category}.                                                                    & \texttt{S-TA-05}                                                       \\ \hline
		\texttt{R-TA-07} & Rancangan revokasi delegasi token Macaroons.                                                                                                            & \texttt{S-TA-06}                                                       \\ \hline
		\texttt{R-TA-08} & Rancangan alur akses dan delegasi pada skenario layanan kesehatan kolaboratif.                                                                          & \texttt{S-TA-03}, \texttt{S-TA-04}, \texttt{S-TA-05}                   \\ \hline
	\end{tabular}
\end{table}
   \chapter{Implementasi dan Evaluasi}
\label{bab:implementasi-dan-evaluasi}

\section{Lingkungan Implementasi}
\label{sec:lingkungan-implementasi}

Implementasi sistem dilakukan pada lingkungan prototipe yang merepresentasikan komponen utama DecMed, yaitu \textit{patient client}, \textit{hospital client}, \textit{ministry client}, \textit{Server} PRE, penyimpanan \textit{off-chain} IPFS, Redis, IOTA \textit{Gas Station}, dan \textit{smart contract} IOTA Move pada jaringan IOTA Localnet. Lingkungan ini dipilih agar rancangan pada Bab \ref{chap:analisis-dan-rancangan} dapat diuji tanpa bergantung pada integrasi langsung dengan sistem informasi rumah sakit yang sesungguhnya.

Pada implementasi DecMed sebelumnya, semua komponen dijalankan pada lingkungan lokal. Semua \textit{client}, \textit{Server} PRE, Redis PRE, kluster IPFS, IOTA \textit{Gas Station}, dan jaringan IOTA jenis Localnet dijalankan langsung di komputer lokal. Perangkat lokal yang digunakan tidak cukup kuat untuk menjalankan semua komponen secara bersamaan. Oleh karena itu, dilakukan pembagian komponen implementasi tersebut ke dalam lingkungan lokal dan dua buah \textit{Virtual Private Server} (VPS).

Komponen \textit{client} dan \textit{Server} PRE dijalankan pada lingkungan lokal. Lingkungan lokal digunakan untuk memudahkan pengembangan fitur, pengujian alur akses, serta pengamatan respons sistem. Komponen pendukung yang membutuhkan ketersediaan jaringan, yaitu kluster IPFS dan IOTA \textit{Gas Station}, dijalankan pada VPS. \textit{Smart contract} IOTA Move di-\textit{deploy} pada jaringan IOTA Localnet yang juga berjalan di VPS. IOTA Localnet digunakan karena sudah memadai untuk menyimulasikan interaksi dengan jaringan IOTA. VPS digunakan untuk menjalankan komponen pendukung karena komponen tersebut tidak banyak berubah selama implementasi, tetapi tetap dibutuhkan selama pengembangan. Pemanfaatan VPS juga mengurangi beban perangkat lokal karena perangkat lokal tidak perlu menjalankan seluruh komponen pendukung.

Berikut adalah spesifikasi dari perangkat lokal yang digunakan untuk implementasi:
\begin{itemize}
	\item Sistem Operasi: Ubuntu 24.04.4, berjalan di lingkungan WSL2
	\item CPU: Intel Core i5-11400H (3 core fisik dan 6 \textit{logical thread} dialokasikan untuk WSL2)
	\item Memori RAM: 16 GB (8 GB teralokasi untuk WSL2)
	\item \textit{Tools}:
	      \begin{itemize}
		      \item Rust (\texttt{rustc 1.93.0}, \texttt{cargo 1.93.0})
		      \item Node.js \texttt{v24.13.0}
		      \item Docker \texttt{29.5.3}
	      \end{itemize}
\end{itemize}

Sementara itu, spesifikasi kedua VPS yang digunakan adalah sebagai berikut:
\begin{itemize}
	\item Sistem Operasi: Ubuntu 24.04.4 LTS
	\item Jumlah vCPU: 8 vCPU
	\item Memori RAM: 8 GB
	\item \textit{Tools}:
	      \begin{itemize}
		      \item Docker: Docker \texttt{29.4.2}
	      \end{itemize}
\end{itemize} \section{Implementasi Sistem}
\label{sec:implementasi-sistem}

Bagian ini menjelaskan implementasi dari rancangan solusi yang telah dijabarkan pada Bab \ref{chap:analisis-dan-rancangan}. Implementasi difokuskan pada integrasi Macaroons sebagai token otorisasi, segmentasi data RME, mekanisme atenuasi dan delegasi akses, penegakan otoritas pada \textit{Server} PRE, perubahan \textit{smart contract} IOTA Move, serta penyesuaian antarmuka \textit{client}.

Daftar komponen implementasi dan keterkaitannya dengan rancangan solusi ditunjukkan pada Tabel \ref{tab:daftar-implementasi-sistem}. ID \texttt{I-TA} digunakan sebagai rujukan pada implementasi terkait. Rancangan ID \texttt{R-TA-01} merupakan rancangan arsitektur sistem yang implementasinya merupakan keseluruhan implementasi sistem DecMed, sehingga tidak muncul pada tabel.

\begin{table}[!ht]
	\centering
	\caption{Daftar Implementasi Sistem}
	\label{tab:daftar-implementasi-sistem}
	\small
	\begin{tabular}{|p{0.12\textwidth}|p{0.48\textwidth}|p{0.30\textwidth}|}
		\hline
		\textbf{ID}      & \textbf{Komponen}                                                                                          & \textbf{Merealisasikan Rancangan}                    \\ \hline
		\texttt{I-TA-01} & \textit{Crate} \texttt{decmed-rme-segment}.                                                                & \texttt{R-TA-02}                                     \\ \hline
		\texttt{I-TA-02} & \textit{Crate} \texttt{decmed-macaroon-auth} dan \textit{Crate} \texttt{macaroon}.                         & \texttt{R-TA-03}, \texttt{R-TA-04}, \texttt{R-TA-05} \\ \hline
		\texttt{I-TA-03} & Implementasi penegakan otoritas pada \textit{Server} PRE.                                                  & \texttt{R-TA-05}, \texttt{R-TA-06}, \texttt{R-TA-07} \\ \hline
		\texttt{I-TA-04} & Implementasi \textit{smart contract} IOTA Move untuk metadata segmen, akses, dan audit.                    & \texttt{R-TA-02}, \texttt{R-TA-04}, \texttt{R-TA-07} \\ \hline
		\texttt{I-TA-05} & Implementasi \textit{Patient Client} untuk pemberian akses awal, pencabutan, dan audit.                    & \texttt{R-TA-07}, \texttt{R-TA-08}                   \\ \hline
		\texttt{I-TA-06} & Implementasi \textit{Hospital Client} untuk akses baca dan tulis segmen RME dan delegasi \textit{offline}. & \texttt{R-TA-02}, \texttt{R-TA-04}, \texttt{R-TA-08} \\ \hline
	\end{tabular}
\end{table}

\subsection{Batasan Implementasi}
\label{subsec:batasan-implementasi}

Implementasi pada Tugas Akhir ini memiliki beberapa batasan sebagai berikut:

\begin{enumerate}
	\item Implementasi dilakukan dalam lingkungan komputer lokal dan VPS.
	      \begin{itemize}
		      \item Semua \textit{client} dan \textit{Server} PRE dijalankan langsung di komputer lokal.
		      \item VPS menjalankan kluster IPFS dan IOTA \textit{Gas Station}.
		      \item \textit{Smart contract} yang digunakan di-\textit{deploy} pada IOTA Localnet yang berjalan di VPS.
	      \end{itemize}
	\item Implementasi mempertahankan arsitektur kerangka kerja DecMed sebagaimana yang dijelaskan pada Subbab \ref{subsec:rancangan-arsitektur-solusi}.
	\item Kluster IPFS yang digunakan dalam implementasi hanya menggunakan konfigurasi \textit{default} sesuai implementasi DecMed sebelumnya. Konfigurasi tersebut terdiri atas tiga buah \textit{peers}.
	\item Segmentasi data sesuai dengan yang dibahas pada \ref{subsec:rancangan-segmentasi-data-rme} hanya dilakukan hingga tingkat kombinasi \textit{dataset category} dan \textit{function category}, kemudian disimpan sebagai \textit{text}.
	\item Nilai \texttt{max\_delegation\_depth} pada token Macaroon hanya diatur sampai tingkat 2. Hal ini untuk memastikan bahwa token turunan hanya dapat diturunkan sebanyak 2 kali sesuai dengan skenario pada Subbab \ref{subsubsec:analisis-skenario-layanan-klinis} dan rancangan pada Subbab \ref{subsec:rancangan-alur-akses-dan-delegasi-macaroon}.
\end{enumerate} \subsection{Implementasi Arsitektur Sistem}
\label{subsec:implementasi-arsitektur-sistem}



Implementasi arsitektur sistem mengikuti rancangan pada Subbab \ref{subsec:rancangan-arsitektur-solusi}. Komponen dasar DecMed dipertahankan, yaitu IOTA untuk metadata \textit{on-chain}, IPFS untuk penyimpanan \textit{off-chain} segmen data yang terenkripsi, dan \textit{Server} PRE sebagai penegak otoritas dan proses \textit{re-encryption} data. Arsitektur DecMed secara keseluruhan dapat dilihat pada gambar \ref{fig:arsitektur-solusi}.

Selain perubahan sesuai rancangan yang sudah dijelaskan pada subbab \ref{sec:rancangan-solusi}, implementasi juga memperbaiki pola pengorganisasian kode. Pada DecMed awal, beberapa bagian yang memiliki fungsi serupa ditulis berulang pada beberapa komponen. Pola tersebut membuat perubahan kode menjadi lebih sulit di-\textit{maintain} karena perubahan pada satu bagian perlu direplikasi pada komponen lain. Pada implementasi ini, bagian yang digunakan bersama dipisahkan menjadi \textit{crate} Rust agar dapat digunakan ulang oleh beberapa komponen.

\textit{Crate} yang diimplementasikan dan digunakan adalah \texttt{decmed-rme-segment} dan \texttt{decmed-macaroon-auth}.
\texttt{decmed-rme-segment} berisi definisi struktur segmentasi RME, definisi \textit{dataset category} dan \textit{function category}, serta fungsi pendukungnya. Sementara itu, \texttt{decmed-macaroon-auth} berisi implementasi token otorisasi berbasis Macaroons di DecMed, termasuk pembuatan token, atenuasi token, dan proses verifikasi. \textit{crate} \texttt{decmed-macaroon-auth} memiliki dependensi utama yaitu \textit{crate} \texttt{macaroon} yang merupakan \textit{crate open source} implementasi Macaroons dengan pembaruan \textit{dependency}. \texttt{decmed-rme-segment} digunakan oleh \textit{Patient Client}, \textit{Hospital Client}, dan \textit{Server} PRE, sedangkan \texttt{decmed-macaroon-auth} digunakan oleh \textit{Hospital Client} dan \textit{Server} PRE. \textit{Ministry Client} tidak memiliki dependensi langsung terhadap kedua \textit{crate} tersebut karena perannya tetap terbatas pada kebutuhan registrasi dan aktivasi rumah sakit. Hubungan antar komponen implementasi dan dependensi \textit{crate} bersama ditunjukkan pada \textit{package diagram} pada Gambar \ref{fig:package-diagram-implementasi} dengan legenda dari \textit{package diagram} tersebut dapat dilihat pada Lampiran \ref{lampiran:legenda-package-diagram}.

\begin{figure}[!ht]
	\centering
	\includegraphics[width=0.9\textwidth]{resources/chapter-4/Package Diagram.png}
	\caption{Package Diagram Implementasi}
	\label{fig:package-diagram-implementasi}
\end{figure}

Dengan pemisahan tersebut, implementasi segmentasi data dan otorisasi tidak tersebar sebagai kode duplikatif pada setiap komponen.

 \subsection{Implementasi Segmentasi Data RME}
\label{subsec:implementasi-segmentasi-data-rme}



Segmentasi data RME diimplementasikan dengan memecah satu data RME menjadi beberapa objek data terenkripsi berdasarkan kombinasi \textit{dataset category} dan \textit{function category}. Kombinasi kategori tersebut mengikuti rancangan pada Subbab \ref{subsec:rancangan-segmentasi-data-rme}. Setiap segmen data RME disimpan sebagai satu entri metadata pada penyimpanan \textit{on-chain} dan satu objek data pada penyimpanan \textit{off-chain}.

Segmentasi data RME diimplementasikan pada \textit{crate} \texttt{decmed-rme-segment}. \textit{Crate} ini mendefinisikan \texttt{DatasetCategory} dan \texttt{FunctionCategory} yang digunakan untuk kontrol akses granular data RME yang diakses. Implementasi \texttt{DatasetCategory} ditunjukkan pada Kode \ref{lst:implementasi-dataset-category-segmentasi-rme}, implementasi \texttt{FunctionCategory} ditunjukkan pada Kode \ref{lst:implementasi-function-category-segmentasi-rme}.

\begin{lstlisting}[language=Rust, caption=Implementasi \texttt{DatasetCategory} pada \textit{crate} \texttt{decmed-rme-segment}, label={lst:implementasi-dataset-category-segmentasi-rme}]
	pub enum DatasetCategory {
		RAWAT_JALAN,
		RAWAT_INAP,
		LABORATORIUM,
		APOTEK,
	}
\end{lstlisting}

\begin{lstlisting}[language=Rust, caption=Implementasi \texttt{FunctionCategory} pada \textit{crate} \texttt{decmed-rme-segment}, label={lst:implementasi-function-category-segmentasi-rme}]
	pub enum FunctionCategory {
		// ADMINISTRATIVE
		ADMINISTRATIVE_GENERAL,
	
		// RAWAT_JALAN & RAWAT_INAP
		ANAMNESIS,
		PEMERIKSAAN_FISIK,
		PEMERIKSAAN_PSIKOLOGIS,
		RIWAYAT_PENGGUNAAN_OBAT,
		RENCANA_RAWAT,
		INSTRUKSI_MEDIK_DAN_KEPERAWATAN,
		PEMERIKSAAN_PENUNJANG,
		DIAGNOSIS,
		INFORMED_CONSENT,
		TERAPI,
	
		// RAWAT_INAP
		PERENCANAAN_PEMULANGAN,
	
		// LABORATORIUM
		LABORATORIUM,
	
		// APOTEK
		PERESEPAN,
		DISPENSING,
	}
\end{lstlisting}

Perbedaan antara \textit{dataset category} rawat jalan dan rawat inap terdapat pada \textit{function category} \texttt{PERENCANAAN\_PEMULANGAN}. Segmen data ini hanya ada pada \textit{dataset category} \texttt{RAWAT\_INAP}, sedangkan pada \textit{dataset category} \texttt{RAWAT\_JALAN} tidak ada. Selain itu, \textit{function category} \texttt{ADMINISTRATIVE\_GENERAL} ada pada setiap \textit{dataset category}.

\subsubsection{Implementasi Segmentasi Metadata RME pada Penyimpanan \textit{On-Chain}}
\label{subsec:implementasi-segmentasi-metadata-rme-onchain}

Definisi struktur metadata segmentasi ditulis pada \textit{crate} \texttt{decmed-rme-segment}. \textit{Crate} ini mendefinisikan \texttt{RmeSegmentMetadata} yang merupakan representasi metadata segmen sebelum diserialisasi menjadi JSON base64 dan disimpan pada \textit{field} \texttt{metadata} milik \texttt{PatientMedicalMetadata} yang sudah ditunjukkan pada Kode \ref{lst:decmed-patient-medical-metadata}. Implementasi struktur \texttt{RmeSegmentMetadata} terdapat pada Kode \ref{lst:implementasi-segmentasi-metadata-rme}.
Struktur tersebut sesuai dengan rancangan skema metadata segmen RME pada Lampiran Tabel \ref{tab:onchain-schema}.

\begin{lstlisting}[language=Rust, caption=Implementasi struktur RmeSegmentMetadata pada crate decmed-rme-segment, label={lst:implementasi-segmentasi-metadata-rme}]
pub struct RmeSegmentMetadata {
	pub segment_id: String,
	pub related_rme_id: String,
	pub patient_address: String,
	pub author_address: String,
	pub hospital_cid: String,
	pub dataset_category: DatasetCategory,
	pub function_category: FunctionCategory,
	pub ipfs_cid: String,
	pub integrity_hash: String,
	pub capsule: String,
	pub enc_key_and_nonce: String,
	pub created_at: String,
}
\end{lstlisting}

Sebelum metadata segmen disimpan pada IOTA, struktur metadata divalidasi terlebih dahulu. Validasi tersebut mencakup pemeriksaan setiap \textit{field} wajib terisi, validitas kombinasi \texttt{dataset\_category} dan \texttt{function\_category}, serta kesesuaian \texttt{integrity\_hash} terhadap data terenkripsi. \texttt{dataset\_category} dan \texttt{function\_category} memiliki nilai sesuai dengan Kode \ref{lst:implementasi-dataset-category-segmentasi-rme} dan \ref{lst:implementasi-function-category-segmentasi-rme}. \texttt{patient\_address} dan \texttt{author\_address} memuat \textit{address wallet} dari pasien dan aktor yang menulis isi segmen. \texttt{hospital\_cid} merupakan \textit{identifier} fasyankes, menunjukkan bahwa segmen data tersebut dihasilkan oleh tenaga medis dari fasyankes tersebut. \texttt{ipfs\_cid} merupakan \textit{pointer} ke data segmen terenkripsi yang disimpan pada penyimpanan \textit{off-chain}. \texttt{integrity\_hash} merupakan hash SHA-256 dari data segmen terenkripsi yang digunakan untuk memverifikasi integritas data. \texttt{capsule} merupakan kapsul PRE yang diperlukan untuk proses \textit{re-encryption}. \texttt{enc\_key\_and\_nonce} merupakan kunci simetris dan \textit{nonce} yang sudah terenkripsi untuk membuka segmen data. \texttt{created\_at} memuat informasi tanggal dan waktu ketika segmen tercatat dengan format \texttt{YYYY-MM-DD HH:MM:SS} serta informasi \textit{timezone}.

\subsubsection{Implementasi Segmentasi Data RME pada Penyimpanan \textit{Off-Chain}}
\label{subsec:implementasi-segmentasi-data-rme-offchain}

Definisi struktur segmentasi data penyimpanan \textit{off-chain} didefinisikan pada \textit{crate} \texttt{decmed-rme-segment}. \textit{Crate} ini mendefinisikan \texttt{RmeSegmentData}, struktur data segmen RME yang disimpan pada penyimpanan \textit{off-chain}. Implementasi struktur \texttt{RmeSegmentData} terdapat pada Kode \ref{lst:implementasi-segmentasi-data-rme}. Struktur tersebut sesuai dengan rancangan skema data segmen RME pada Lampiran Tabel \ref{tab:offchain-schema}.

\begin{lstlisting}[language=Rust, caption=Implementasi struktur \texttt{RmeSegmentData} pada \textit{crate} \texttt{decmed-rme-segment}, label={lst:implementasi-segmentasi-data-rme}]
pub struct RmeSegmentData {
    pub segment_id: String,
    pub related_rme_id: String,
    pub dataset_category: DatasetCategory,
    pub function_category: FunctionCategory,
    pub patient_address: String,
    pub service_date: String,
    pub author_address: String,
    pub payload: Value,
    pub payload_hash: String,
}
\end{lstlisting}

\texttt{segment\_id} merupakan ID unik untuk satu segmen data RME. \texttt{dataset\_category} dan \texttt{function\_category} digunakan untuk kontrol akses granular data RME. atribut tersebut memiliki nilai sesuai dengan Kode \ref{lst:implementasi-dataset-category-segmentasi-rme} dan \ref{lst:implementasi-function-category-segmentasi-rme}. \texttt{patient\_address} dan \texttt{author\_address} memuat \textit{address wallet} dari pasien dan aktor yang menulis isi segmen.\texttt{service\_date} merupakan tanggal dan waktu ketika segmen tercatat dengan format \texttt{YYYY-MM-DD HH:MM:SS} beserta informasi \textit{timezone}. \texttt{payload} berisi konten utama segmen dengan tipe data \texttt{Value}. Pada implementasi ini, \texttt{payload} berisi data dalam bentuk \texttt{text}. \texttt{payload\_hash} merupakan hash integritas dari konten \texttt{payload}.



Pada saat data RME disimpan, \textit{Hospital Client} membentuk objek untuk setiap segmen. Setiap segmen berisi data klinis atau administratif sesuai kategorinya, kemudian dienkripsi sebelum diunggah ke IPFS. Hasil unggahan IPFS berupa CID disimpan bersama metadata segmen pada IOTA sebagai \texttt{ipfs\_cid}. Metadata tersebut disimpan dalam bentuk serialisasi JSON dan dikodekan dalam base64 pada \textit{field} \texttt{metadata} milik \texttt{PatientMedicalMetadata}. Keputusan ini menjaga kompatibilitas dengan struktur penyimpanan DecMed yang masih berbentuk daftar metadata \textit{flat}, tetapi tetap memungkinkan \textit{Server} PRE membaca kategori segmentasi dari isi metadata.







%
 \subsection{Implementasi Token Otorisasi Berbasis Macaroons}
\label{subsec:implementasi-token-macaroons}

\subsubsection{Komponen Implementasi Macaroons}
\label{subsubsec:komponen-implementasi-macaroons}

Seperti yang sudah dijelaskan pada subbab \ref{subsec:implementasi-arsitektur-sistem}, terdapat dua \textit{crate} yang berkaitan dengan implementasi Macaroons pada implementasi Tugas Akhir ini, yaitu \texttt{macaroon} dan \texttt{decmed-macaroon-auth}.

\textit{Crate} \texttt{macaroon} merupakan hasil penyesuaian dari implementasi Macaroons berbasis Rust yang tersedia secara \textit{open source} pada URL \url{https://github.com/macaroon-rs/macaroon}. Penyesuaian dilakukan karena \textit{crate} asli memerlukan beberapa penyesuaian agar lebih sesuai dengan kebutuhan implementasi sistem DecMed, terutama dari sisi kompatibilitas \textit{toolchain} dan pengelolaan \textit{dependency}. Perubahan utama pada \textit{crate} \texttt{macaroon} adalah penggantian \textit{dependency} kriptografi berbasis \texttt{sodiumoxide} dengan \textit{crate} pure Rust, yaitu \texttt{hmac}, \texttt{sha2}, \texttt{chacha20poly1305}, dan \texttt{rand}.

Setelah perubahan tersebut dilakukan, \textit{unit testing} pada \textit{crate} dijalankan untuk memastikan bahwa fungsi utama Macaroons tetap berjalan dengan benar. Pengujian tersebut mencakup pembuatan Macaroon, penambahan \textit{caveat}, serialisasi, deserialisasi, serta verifikasi token. Dengan demikian, \textit{crate} \texttt{macaroon} digunakan sebagai fondasi implementasi Macaroons pada sistem DecMed.

\texttt{decmed-macaroon-auth} merupakan \textit{wrapper} yang dibangun untuk mempermudah penggunaan \textit{crate} \texttt{macaroon} pada implementasi sistem DecMed. \textit{Wrapper} ini menyediakan fungsi-fungsi utilitas untuk kebutuhan otorisasi DecMed, seperti proses penerbitan token Macaroon, melakukan atenuasi dengan penambahan \textit{caveat} baru, definisi struktur \textit{caveats} Macaroon, verifikasi Macaroon, dan fungsi pendukung lainnya. Rincian modul dan implementasi yang diimplementasikan pada \textit{crate} \texttt{decmed-macaroon-auth} ditunjukkan pada Lampiran \ref{lampiran:implementasi-crate-decmed-macaroon-auth}.

\subsubsection{Struktur \textit{Caveats} Macaroons}
\label{subsubsec:struktur-caveats-macaroons}

Struktur \textit{caveats} pada token Macaroon pada implementasi ini ditunjukkan pada Tabel \ref{tab:implementasi-caveats-macaroons}. \textit{caveats} terdefinisi pada \textit{crate} \texttt{decmed-macaroon-auth}, sehingga \textit{hospital client} membentuk token dengan menambahkan \textit{caveats} sesuai dengan daftar \textit{caveats} yang terdefinisi pada \textit{crate}, dan PRE akan menolak \textit{request} jika terdapat \textit{caveats} yang tidak termasuk ke dalam daftar \textit{caveats}.

\begin{table}[!ht]
	\centering
	\caption{\textit{Caveats} yang Diimplementasikan pada Token Macaroon}
	\label{tab:implementasi-caveats-macaroons}
	\small
	\begin{tabular}{|p{0.34\textwidth}|p{0.56\textwidth}|}
		\hline
		\textbf{\textit{Caveat}}        & \textbf{Tipe Data}
		\\ \hline
		\texttt{patient\_address}       & \texttt{String} (alamat IOTA format heksadesimal)
		\\ \hline
		\texttt{related\_rme\_id}       & \texttt{String} (ID episode RME)
		\\ \hline
		\texttt{root\_subject}          & \texttt{String} (alamat IOTA format heksadesimal)
		\\ \hline
		\texttt{delegated\_by}          & \texttt{String} (alamat IOTA format heksadesimal)
		\\ \hline
		\texttt{delegated\_to}          & \texttt{String} (alamat IOTA format heksadesimal)
		\\ \hline
		\texttt{read\_dataset\_in}      & \texttt{HashSet<DatasetCategory>}
		\\ \hline
		\texttt{write\_dataset\_in}     & \texttt{HashSet<DatasetCategory>}
		\\ \hline
		\texttt{read\_function\_in}     & \texttt{HashSet<FunctionCategory>}
		\\ \hline
		\texttt{write\_function\_in}    & \texttt{HashSet<FunctionCategory>}
		\\ \hline
		\texttt{expires\_before}        & \texttt{DateTime<Utc>} (\texttt{YYYY-MM-DD HH:MM:SS})
		\\ \hline
		\texttt{max\_delegation\_depth} & \texttt{u32}
		\\ \hline
		\texttt{hospital\_cid}          & \texttt{String} (CID IPFS)
		\\ \hline
		\texttt{parent\_token\_hash}    & \texttt{String} (hash SHA-256 dari token induk)
		\\ \hline
		\texttt{purpose}                & \texttt{String} (nilai: \texttt{Read} atau \texttt{Update})
		\\ \hline
	\end{tabular}
\end{table}

\subsubsection{Penerbitan Token Macaroon}
\label{subsubsec:penerbitan-awal-token-macaroons}

Token otorisasi diimplementasikan menggunakan Macaroons sesuai rancangan pada Subbab \ref{subsec:rancangan-token-macaroons}. Pada implementasi ini, penerbitan token dilakukan oleh \textit{Server} PRE setelah permintaan akses dari pasien dinyatakan valid.

Pada proses penerbitan akses awal, \textit{Server} PRE menerbitkan dua jenis token untuk penerima akses, yaitu token \texttt{Read} dan token \texttt{Update}. Token \texttt{Read} memuat \textit{caveat} utama berupa \texttt{patient\_address}, \texttt{root\_subject}, \texttt{read\_dataset\_in}, \texttt{read\_function\_in}, \texttt{expires\_before}, \texttt{max\_delegation\_depth}, \texttt{hospital\_cid}, dan \texttt{purpose = Read}. Token \texttt{Update} memuat \textit{caveat} berupa \texttt{patient\_address}, \texttt{related\_rme\_id}, \texttt{root\_subject}, \texttt{write\_dataset\_in}, \texttt{write\_function\_in}, \texttt{expires\_before}, \texttt{max\_delegation\_depth}, \texttt{hospital\_cid}, dan \texttt{purpose = Update}.

Ketika melakukan penerbitan untuk petugas rekam medis, fungsi yang dipanggil adalah \texttt{issue\_admin\_personnel\_token} yang termasuk implementasi pada \textit{crate} \texttt{decmed-macaroon-auth}. Alur lengkap pemberian akses awal dari pasien kepada petugas rekam medis, termasuk pemindaian QR code, verifikasi identitas, pembentukan material PRE, penyimpanan ke Redis, dan pencatatan grant akses ke \textit{smart contract}, dijelaskan pada subbab \ref{subsec:implementasi-atenuasi-offline-delegasi}.

\subsubsection{Atenuasi Token Macaroon}
\label{subsubsec:atenuasi-token-macaroons}

Pada implementasi ini, pemegang token utama adalah para tenaga kesehatan, dengan role \texttt{AdministrativePersonnel} atau \texttt{MedicalPersonnel}. Pemegang token melakukan atenuasi token melalui \textit{hospital client} dengan melakukan \textit{request} atenuasi token ke \textit{server} PRE. Implementasi dari fungsi atenuasi ini terdapat pada \textit{crate} \texttt{decmed-macaroon-auth} dengan nama fungsi \texttt{attenuate\_macaroon}. Fungsi ini bekerja dengan tahapan melakukan deserialisasi token Macaroon, lalu menambahkan \textit{caveat} baru sesuai dengan \textit{caveats} yang ingin ditambahkan, kemudian melakukan serialisasi kembali terhadap token Macaroon yang telah diatenuasi sebelum di-\textit{return} sebagai \textit{response}. Perhitungan \textit{signature} baru dari Macaroon pada proses \textit{chained-HMAC} dilakukan pada setiap kali \textit{caveat} baru ditambahkan.

\subsubsection{Revokasi Token Macaroon}
\label{subsubsec:revokasi-token-macaroons}

Revokasi token pada implementasi ini sesuai dengan rancangan pada Subbab \ref{subsec:rancangan-penegakan-otoritas}. Implementasi revokasi akses melibatkan pencatatan di IOTA \textit{on-chain}, Redis PRE, dan \textit{server} PRE.

Setiap token turunan menyimpan keterkaitan terhadap token induknya melalui \textit{caveat} \texttt{parent\_token\_hash}. Informasi ini digunakan untuk mendukung pencabutan berantai, sehingga pencabutan terhadap token pada tingkat tertentu dapat berdampak pada token-token turunan yang bergantung kepadanya.

Pemeriksaan pencabutan dilakukan oleh \textit{Server} PRE sebelum akses terhadap data terenkripsi diberikan. Saat \textit{request} diterima oleh PRE, maka PRE akan memeriksa \textit{hash} dari token Macaroon yang digunakan, dan mengecek apakah ada di daftar \textit{revoked keys} yang tersimpan pada Redis PRE dengan pola \verb|revoked:token:<hash>|. Hal ini juga dilakukan dengan semua \textit{hash} token yang ada pada \textit{caveat} \texttt{parent\_token\_hash}. Jika ada salah satu \textit{hash} yang ditemukan, maka \textit{request} akan ditolak. Selain itu, proses ini juga akan menandai hak akses yang tercatat pada \textit{on-chain} sebagai \textit{revoked}. \subsection{Delegasi Akses dan Atenuasi \textit{Offline}}
\label{subsec:implementasi-atenuasi-offline-delegasi}



Mekanisme ini merupakan realisasi dari rancangan pada Subbab \ref{subsec:rancangan-alur-akses-dan-delegasi-macaroon}. Implementasi dibagi menjadi dua bagian. Bagian pertama adalah proses pemberian akses awal dari pasien kepada petugas rekam medis sebagai penerima akses awal. Bagian kedua adalah proses delegasi akses antar tenaga kesehatan.

\subsubsection{Delegasi Akses dari Pasien kepada Petugas Rekam Medis}
\label{subsubsec:delegasi-akses-dari-pasien-kepada-petugas-rekam-medis}

Implementasi delegasi akses dari pasien kepada petugas rekam medis merupakan realisasi dari rancangan pada Subbab \ref{subsec:delegasi-dari-petugas-rekam-medis-kepada-perawat-dan-dokter} dan \ref{subsec:delegasi-dari-dokter-kepada-laboratorium-dan-apoteker}.

Proses diawali ketika pasien melakukan \textit{scan} QR code milik \texttt{AdministrativePersonnel}. QR code tersebut memuat sebuah \textit{string} dengan format \verb|{hospitalCid}@{personnelIota}@{personnelPrePK}|. Informasi tersebut terdiri dari \texttt{hospital\_cid} sebagai \textit{identifier} fasyankes, \texttt{personnel\_iota\_address} sebagai alamat IOTA milik \texttt{AdministrativePersonnel}, dan \texttt{personnel\_pre\_public\_key} sebagai PRE \textit{public key} milik \texttt{AdministrativePersonnel} dalam format base64. PRE \textit{public key} tersebut digunakan untuk mengenkripsi \texttt{data\_pre\_secret\_key\_seed} yang nantinya hanya dapat dibuka oleh petugas penerima akses.

Setelah QR code berhasil di-\textit{scan}, \textit{Patient Client} mengambil informasi publik \texttt{AdministrativePersonnel} melalui \textit{smart contract} \verb|get_hospital_personnel_info|. Langkah ini dilakukan untuk memastikan identitas petugas yang akan menerima akses. Pasien kemudian memilih jenis layanan, yaitu \texttt{RAWAT\_JALAN} atau \texttt{RAWAT\_INAP}, serta memilih mode akses yang akan diberikan. Implementasi mendukung tiga mode akses, yaitu \texttt{read}, \texttt{update}, dan \texttt{read\_update}. Mode \texttt{read} hanya menerbitkan token baca, mode \texttt{update} hanya menerbitkan token pembaruan, sedangkan mode \texttt{read\_update} menerbitkan kedua token tersebut. Pada implementasi, batas waktu akses bawaan untuk \texttt{RAWAT\_JALAN} adalah 1 hari, sedangkan batas waktu akses bawaan untuk \texttt{RAWAT\_INAP} adalah 3 hari. Namun, implementasi juga mendukung nilai \texttt{expires\_before} khusus sehingga batas waktu akses dapat disesuaikan dengan kebutuhan.

Selanjutnya, \textit{Patient Client} meminta \textit{nonce} kepada \textit{Server} PRE melalui \textit{endpoint} \texttt{/nonce}. Sebelum \textit{nonce} diterbitkan, \textit{Server} PRE memeriksa terlebih dahulu apakah alamat pasien telah terdaftar pada \textit{smart contract} \verb|is_patient_registered|. \textit{Nonce} yang diterima kemudian di-\textit{sign} menggunakan kunci IOTA pasien. Tanda tangan ini digunakan sebagai bukti bahwa permintaan penerbitan akses benar-benar berasal dari pasien yang sah.

Bersamaan dengan proses tersebut, \textit{Patient Client} membuat \texttt{data\_pre\_secret\_key\_seed} baru dan menurunkan \texttt{data\_pre\_public\_key} dari \textit{seed} tersebut. \textit{Seed} kemudian dienkripsi menggunakan \texttt{personnel\_pre\_public\_key}. Selain itu, \textit{Patient Client} juga membuat \textit{key fragment} \texttt{k\_frag} dari PRE \textit{secret key} pasien menuju \texttt{data\_pre\_public\_key}. Material PRE ini disimpan pada Redis PRE agar \textit{Server} PRE dapat melakukan \textit{re-encryption} ketika petugas mengakses data pasien. TTL untuk material PRE dihitung berdasarkan waktu kedaluwarsa aktual dari token yang diterbitkan.

Permintaan penerbitan akses kemudian dikirim ke \textit{endpoint} \texttt{/keys}. \textit{Payload} yang dikirim mencakup alamat IOTA pasien, alamat IOTA petugas, \texttt{hospital\_cid}, \texttt{encounter\_dataset}, mode akses, \textit{signature} pasien, \texttt{k\_frag}, \texttt{data\_pre\_public\_key}, \texttt{enc\_data\_pre\_secret\_key\_seed}, \texttt{data\_pre\_secret\_key\_seed\_capsule}, \texttt{patient\_pre\_public\_key}, dan \texttt{signer\_pre\_public\_key}. \textit{Server} PRE kemudian memverifikasi \textit{signature} pasien terhadap \textit{nonce} dan memeriksa apakah \textit{role} petugas penerima akses adalah \texttt{AdministrativePersonnel}.

Jika verifikasi berhasil, \textit{Server} PRE membentuk \texttt{MacaroonKey} dari \texttt{RootKey} dan menerbitkan token Macaroon sesuai mode akses yang diminta. Jika mode yang dipilih adalah \texttt{read}, \textit{Server} PRE hanya menerbitkan token \texttt{Read}. Jika mode yang dipilih adalah \texttt{update}, \textit{Server} PRE hanya menerbitkan token \texttt{Update}. Jika mode yang dipilih adalah \texttt{read\_update}, \textit{Server} PRE menerbitkan token \texttt{Read} dan token \texttt{Update}.

Token \texttt{Read} memiliki \textit{caveat} \texttt{patient\_address}, \texttt{root\_subject}, \texttt{read\_dataset\_in}, \texttt{read\_function\_in}, \texttt{expires\_before}, \texttt{max\_delegation\_depth}, \texttt{hospital\_cid}, dan \texttt{purpose = Read}. Nilai \texttt{read\_dataset\_in} dan \texttt{read\_function\_in} pada token ini mencakup seluruh \textit{dataset category} dan seluruh \textit{function category}. Dengan demikian, token \texttt{Read} memberikan kemampuan baca terhadap data RME pasien.

Token \texttt{Update} memiliki \textit{caveat} \texttt{patient\_address}, \texttt{related\_rme\_id}, \texttt{root\_subject}, \texttt{write\_dataset\_in}, \texttt{write\_function\_in}, \texttt{expires\_before}, \texttt{max\_delegation\_depth}, \texttt{hospital\_cid}, dan \texttt{purpose = Update}. Berbeda dengan token \texttt{Read}, cakupan \texttt{write\_dataset\_in} pada token \texttt{Update} tidak mencakup seluruh \textit{dataset category}. Nilainya dibatasi pada \texttt{encounter\_dataset} yang dipilih pasien, ditambah \texttt{LABORATORIUM} dan \texttt{APOTEK}. Nilai \texttt{write\_function\_in} diturunkan dari \textit{function category} yang termasuk ke dalam cakupan \textit{dataset category} tersebut.

\texttt{related\_rme\_id} hanya dibuat apabila akses yang diterbitkan mencakup kemampuan \texttt{Update}. Pada kondisi tersebut, \textit{Server} PRE membuat \texttt{related\_rme\_id} baru sebagai identitas episode rekam medis yang sedang dibuat atau diperbarui. ID tersebut menggunakan format \verb|RME-{sequence_6_digit}|, misalnya \verb|RME-000001|. Perhitungan \textit{sequence} dilakukan pada \textit{Server} PRE dengan memanfaatkan Redis untuk menyimpan nilai \textit{sequence}. Apabila pasien hanya memberikan akses \texttt{read}, \textit{Server} PRE tidak membuat \texttt{related\_rme\_id}.

Setelah token diterbitkan, \textit{Server} PRE menyimpan material akses PRE ke Redis PRE. Entri PRE menggunakan \textit{key} \verb|{personnel_iota_address}@{patient_iota_address}|. \textit{Server} PRE kemudian menghitung \textit{hash} dari token yang diterbitkan. PRE lalu mengembalikan \texttt{access\_token\_read}, \texttt{access\_token\_update}, \textit{hash} masing-masing token, serta \texttt{related\_rme\_id} kepada \textit{Patient Client}.

Pada tahap akhir, \textit{Patient Client} mengenkripsi token dan metadata akses menggunakan \texttt{personnel\_pre\_public\_key}. Metadata akses kemudian dicatat ke IOTA melalui \textit{smart contract} \texttt{create\_access}. Dengan demikian, token Macaroon awal telah dirilis kepada \texttt{AdministrativePersonnel} secara terenkripsi, \textit{hash} token telah tersedia untuk keperluan audit dan revokasi, dan material PRE telah tersimpan di Redis PRE untuk proses akses data berikutnya.

\subsubsection{Delegasi Akses antar Tenaga Kesehatan}
\label{subsubsec:delegasi-akses-antar-tenaga-kesehatan}

Implementasi delegasi akses antar tenaga kesehatan merupakan realisasi dari rancangan pada Subbab \ref{subsec:delegasi-dari-petugas-rekam-medis-kepada-perawat-dan-dokter} dan \ref{subsec:delegasi-dari-dokter-kepada-laboratorium-dan-apoteker}. Tenaga medis memiliki \textit{role} \texttt{MedicalPersonnel} dengan \textit{subrole} masing-masing sesuai dengan tabel \ref{tab:pemetaan-hak-akses-update-subrole}.

Proses diawali ketika \textit{delegator} memilih \textit{delegatee} yang akan menerima akses. Informasi kandidat \textit{delegatee} diperoleh melalui \textit{smart contract} \texttt{get\_delegatee\_candidates}. Data yang digunakan dalam proses delegasi mencakup alamat IOTA \textit{delegatee} dan PRE \textit{public key} milik \textit{delegatee}. PRE \textit{public key} tersebut digunakan untuk mengenkripsi metadata delegasi dan material akses agar hanya dapat dibuka oleh \textit{delegatee}.

Sebelum membentuk token turunan, \textit{Hospital Client} memuat material kriptografis milik \textit{delegator}, yaitu PRE \textit{secret key}, alamat IOTA, dan pasangan kunci IOTA. \textit{Delegator} kemudian mendekripsi \texttt{data\_pre\_secret\_key\_seed} yang terenkripsi, lalu mengenkripsi ulang \textit{seed} tersebut menggunakan PRE \textit{public key} milik \textit{delegatee}. Proses ini memungkinkan \textit{delegatee} memperoleh material yang diperlukan untuk membuka data yang didelegasikan.

\textit{Hospital client} kemudian mengirim \textit{request} atenuasi token ke \textit{server} PRE. request memuat token induk, requested \textit{scope}, \textit{expiry}, alamat \textit{delegator/delegatee/patient}, \textit{parent delegation signature} bila token induk merupakan token hasil delegasi, serta \texttt{delegation\_request\_signature}. Signature terakhir merupakan tanda tangan IOTA \textit{delegator} atas \texttt{DelegationRequestProofContext} dan digunakan PRE untuk memastikan bahwa permintaan atenuasi benar-benar disetujui oleh \textit{delegator}.

Pertama PRE memverifikasi autentikasi token induk melalui mekanisme \textit{chained-HMAC}, mengecek apakah token ada di daftar yang sudah direvokasi, dan dilanjutkan dengan pemeriksaan \texttt{EffectiveCapability} \texttt{EffectiveCapability} dari token induk divalidasi terhadap \texttt{requested scope}. Validasi ini memastikan bahwa waktu kedaluwarsa token turunan tidak melebihi token induk, cakupan \textit{dataset} dan \textit{function} merupakan subset dari token induk, serta nilai \texttt{max\_delegation\_depth} tidak bertambah. Definisi \texttt{EffectiveCapability} dapat dilihat pada \ref{lst:definisi-effective-capability}.

Pemeriksaan kemudian dilanjutkan dengan pemeriksaan metadata akses pada IOTA melalui \texttt{get\_delegation\_access\_snapshot}. Pemeriksaan ini memeriksa \texttt{expiry} dari hak akses yang dimiliki oleh \textit{delegator}, fasyankes sama, \textit{delegation depth} saat ini, dan memeriksa apakah terdapat delegasi aktif terhadap sebuah \textit{subrole} yang telah didelegasikan oleh \textit{delegator}. Jika sudah ada, maka \textit{request} ditolak.

Apabila verifikasi berhasil, \textit{server} PRE melakukan atenuasi pada token induk dengan menambahkan \textit{caveat} melalui fungsi \texttt{attenuate\_macaroon} dari \textit{crate} \texttt{decmed-macaroon-auth}. Fungsi ini menerima token Macaroon induk dan parameter atenuasi berupa alamat \textit{delegator}, alamat \textit{delegatee}, cakupan \textit{dataset category}, cakupan \textit{function category}, waktu kedaluwarsa, dan batas kedalaman delegasi. \textit{Server} PRE kemudian mengembalikan token turunan untuk diproses lebih lanjut oleh \textit{client delegator} untuk dicatat pada penyimpanan \textit{on-chain}. \textit{Response} ke \textit{client delegator} memuat token turunan, \textit{hash} token turunan, dan konteks bukti delegasi yang akan menjadi bukti delegasi.

\textit{Caveat} yang mungkin ditambahkan pada token turunan adalah \texttt{parent\_token\_hash}, \texttt{delegated\_by}, \texttt{delegated\_to}, \texttt{read\_dataset\_in}, \texttt{write\_dataset\_in}, \texttt{read\_function\_in}, \texttt{write\_function\_in}, \texttt{expires\_before}, dan \texttt{max\_delegation\_depth}. Jika mengikuti \textit{preset} pada \textit{hospital client}, maka \textit{caveats} yang ditambahkan dapat dilihat pada Lampiran \ref{lampiran:contoh-caveat-macaroon} sesuai dengan alur delegasinya.

Setelah token turunan diterima oleh \textit{client delegator}, \textit{client} membuat bukti delegasi menggunakan \texttt{DelegationProofContext}. Struktur \texttt{DelegationProofContext} ditunjukkan pada Kode \ref{lst:definisi-delegation-proof-context}. Nilai \texttt{scope\_fingerprint} merupakan \textit{hash} dengan struktur \texttt{r:\{ds\_read\}|w:\{ds\_write\}|rf:\{fn\_read\}|wf:\{fn\_write\}}, dengan \verb|ds_read|, \verb|ds_write|, \verb|fn_read|, dan \verb|fn_write| merupakan \textit{dataset category} dan \textit{function category} yang didelegasikan.  \texttt{DelegationProofContext} kemudian diserialisasi dan ditandatangani menggunakan kunci IOTA milik \textit{delegator}.

\begin{lstlisting}[language=Rust, caption={Definisi \texttt{DelegationProofContext}}, label={lst:definisi-delegation-proof-context}]
pub struct DelegationProofContext {
    pub token_hash: String,
    pub delegated_by: String,
    pub delegated_to: String,
    pub related_rme_id: String,
    pub expires_before: String,
    pub scope_fingerprint: String,
}
\end{lstlisting}

Token turunan, \textit{hash} token, tanda tangan delegasi, metadata akses, dan \texttt{data\_pre\_secret\_key\_seed} yang telah dienkripsi ulang kemudian dikemas menjadi metadata delegasi. Metadata tersebut dienkripsi menggunakan PRE \textit{public key} milik \textit{delegatee}. Dengan demikian, token dan material akses hanya dapat dibuka oleh \textit{delegatee}.

Tahap berikutnya adalah pencatatan delegasi ke IOTA melalui \textit{smart contract} fungsi \texttt{create\_delegated\_access}. Fungsi ini menerima metadata delegasi, data audit delegasi, alamat pasien, dan alamat \textit{delegatee}. Pada sisi \textit{smart contract}, dilakukan validasi bahwa \textit{delegatee} terdaftar dan aktif, keduanya berada pada fasyankes yang sama, serta \textit{delegator} masih memiliki akses yang valid terhadap pasien terkait. Jika validasi berhasil, akses baru untuk \textit{delegatee} dimasukkan ke tabel akses \texttt{Read} dan \texttt{Update}.

Selain menyimpan akses untuk \textit{delegatee}, fungsi \texttt{create\_delegated\_access} juga mencatat audit delegasi pada data pasien. Entri audit tersebut memuat alamat \textit{delegator}, \textit{root subject} dari rantai delegasi, alamat \textit{delegatee}, jenis akses, \texttt{related\_rme\_id}, kedalaman delegasi, \textit{hash} token turunan, \textit{hash} token induk, waktu kedaluwarsa. Dengan demikian, perpindahan akses antar tenaga kesehatan tetap dapat ditelusuri tanpa membuka isi token Macaroon dan metadata akses yang terenkripsi.

Khusus untuk delegasi dari \texttt{AdministrativePersonnel} kepada \texttt{MedicalPersonnel}, delegasi dari \texttt{AdministrativePersonnel} pertama akan menambahkan \textit{function category} \texttt{ADMINISTRATIVE\_GENERAL} pada setiap \textit{dataset category}. Hal ini mengasumsikan bahwa \texttt{AdministrativePersonnel} sudah memeriksa data \texttt{ADMINISTRATIVE\_GENERAL} pasien, sehingga siap dicatat pada setiap \textit{dataset category}. Hal ini dilakukan agar sesuai dengan Lampiran \ref{lampiran:metadata-rme} yang memuat dokumentasi administratif pasien pada masin-masing \textit{dataset category}.

Ketika \textit{delegatee} menggunakan akses yang telah didelegasikan, \textit{request} melewati \textit{middleware} otorisasi pada \textit{Server} PRE. \textit{Server} PRE memverifikasi sesuai dengan implementasi yang dijelaskan pada subbab \ref{subsec:implementasi-penegakan-otoritas-pre}. Untuk memperoleh material PRE, \textit{Server} PRE mencari entri Redis berdasarkan subjek aktif, aktor pada rantai delegasi, hingga \texttt{root\_subject}. Mekanisme ini memungkinkan token turunan menggunakan material PRE dari akses awal, sementara pembatasan akses tetap ditegakkan dengan memeriksa Macaroon.

Dengan mekanisme tersebut, delegasi antar tenaga kesehatan dapat dilakukan oleh \textit{delegator} melalui atenuasi Macaroon di \textit{server} PRE. Hak akses \textit{delegatee} tetap tidak dapat melebihi hak akses \textit{delegator}, proses delegasi tercatat pada \textit{smart contract}, dan penggunaan akses berikutnya tetap diverifikasi oleh \textit{Server} PRE. \subsection{Implementasi Penegakan Otoritas}
\label{subsec:implementasi-penegakan-otoritas-pre}

Penegakan otoritas pada \textit{Server} PRE diimplementasikan sebagai tahap validasi berlapis sebelum operasi terhadap data RME dijalankan. Secara umum, mekanisme ini memastikan bahwa setiap \textit{request} yang membutuhkan otorisasi membawa Macaroon yang valid, belum kedaluwarsa, belum dicabut, sesuai dengan cakupan akses yang diminta, dan digunakan oleh aktor aktif yang sah.

Alur implementasi penegakan otoritas pada \textit{Server} PRE disajikan sebagai referensi pada Lampiran \ref{lampiran:alur-implementasi-penegakan-otoritas-pre}. Bagian awal alur pada Gambar \ref{fig:alur-penegakan-otoritas-pre-a} mencakup autentikasi token, perhitungan \textit{effective capability}, pembentukan \textit{delegation chain}, dan pemeriksaan revokasi. Kelanjutan alur pada Gambar \ref{fig:alur-penegakan-otoritas-pre-b} mencakup validasi \textit{role} dan \textit{subrole}, validasi konteks segmen, verifikasi bukti kepemilikan \textit{wallet}, serta eksekusi operasi PRE, IPFS, dan IOTA. Status implementasi penegakan otoritas pada \textit{Server} PRE disajikan pada Tabel \ref{tab:status-implementasi-penegakan-otoritas-pre}.

\begin{table}[!ht]
	\centering
	\caption{Status Implementasi Penegakan Otoritas pada \textit{Server} PRE}
	\label{tab:status-implementasi-penegakan-otoritas-pre}
	\small
	\begin{tabular}{|p{0.62\textwidth}|p{0.28\textwidth}|}
		\hline
		\textbf{Bagian}                                                                           & \textbf{Status}  \\ \hline
		Validasi \textit{Bearer Macaroon}                                                         & Terimplementasi. \\ \hline
		Validasi \textit{chained-HMAC} pada token Macaroon                                        & Terimplementasi. \\ \hline
		\textit{Parsing} \textit{caveats} pada token Macaroon                                     & Terimplementasi. \\ \hline
		Perhitungan \texttt{EffectiveCapability}                                                  & Terimplementasi. \\ \hline
		Penentuan aktor aktif melalui \texttt{DelegationChain}                                    & Terimplementasi. \\ \hline
		Pemeriksaan revokasi token melalui Redis                                                  & Terimplementasi. \\ \hline
		Pemeriksaan \textit{role} dan \textit{subrole} aktor melalui \textit{Move Smart Contract} & Terimplementasi. \\ \hline
		Validasi aksi baca dan tulis terhadap cakupan \textit{dataset} dan \textit{function}      & Terimplementasi. \\ \hline
		Verifikasi \textit{wallet signature} terhadap aktor aktif dan konteks akses segmen        & Terimplementasi. \\ \hline
\end{tabular}
\end{table}

\subsubsection{Autentikasi Token Macaroon}
\label{subsec:autentikasi-token-macaroon}

Tahap awal validasi dilakukan pada \textit{middleware} PRE. \textit{Middleware} membaca token Macaroon dari \textit{header} \texttt{Authorization} dengan skema \texttt{Bearer}, melakukan deserialisasi token, mem-parsing \textit{caveats} DecMed, serta memverifikasi keaslian token menggunakan mekanisme \textit{chained-HMAC} dengan \texttt{RootKey} PRE. Apabila token tidak dapat dideserialisasi, bukan merupakan token DecMed, atau nilai HMAC tidak valid, \textit{request} langsung ditolak.

Setelah token dinyatakan autentik, \textit{Server} PRE membentuk \texttt{EffectiveCapability} dari seluruh \textit{caveats} yang terdapat pada token. Struktur ini merepresentasikan cakupan akses efektif yang berlaku setelah proses pembatasan dan penggabungan \textit{caveats}. Untuk \textit{caveat} berbentuk \textit{set}, seperti \texttt{read\_dataset\_in}, \texttt{write\_dataset\_in}, \texttt{read\_function\_in}, dan \texttt{write\_function\_in}, cakupan efektif diperoleh melalui irisan antar-\textit{caveats}. Untuk \texttt{expires\_before}, waktu kedaluwarsa efektif diambil dari nilai paling awal. Struktur \texttt{EffectiveCapability} ditunjukkan pada Kode \ref{lst:definisi-effective-capability}.

\begin{lstlisting}[language=Rust, caption={Definisi \texttt{EffectiveCapability}}, label={lst:definisi-effective-capability}]
pub struct EffectiveCapability {
    pub read_datasets: HashSet<DatasetCategory>,
    pub write_datasets: HashSet<DatasetCategory>,
    pub read_functions: HashSet<FunctionCategory>,
    pub write_functions: HashSet<FunctionCategory>,
    pub expires_before: Option<DateTime<Utc>>,
    pub root_max_delegation_depth: Option<u32>,
    pub remaining_max_delegation_depth: Option<u32>,
    pub patient_address: Option<String>,
    pub related_rme_id: Option<String>,
    pub hospital_cid: Option<String>,
}
\end{lstlisting}

Selain menghitung \texttt{EffectiveCapability}, \textit{Server} PRE membentuk \texttt{DelegationChain} untuk menentukan aktor aktif dari token. Jika token belum didelegasikan, aktor aktif adalah \texttt{root\_subject}. Jika token telah didelegasikan, aktor aktif adalah nilai \texttt{delegated\_to} pada langkah delegasi terakhir. Rantai delegasi juga divalidasi agar setiap langkah delegasi tersambung secara konsisten, yaitu penerima delegasi pada satu langkah harus menjadi pemberi delegasi pada langkah berikutnya.
Struktur \texttt{DelegationChain} ditunjukkan pada Kode \ref{lst:definisi-delegation-chain}.

\begin{lstlisting}[language=Rust, caption={Definisi \texttt{DelegationStep} dan \texttt{DelegationChain}}, label={lst:definisi-delegation-chain}]
pub struct DelegationStep {
    pub delegated_by: String,
    pub delegated_to: String,
}

pub struct DelegationChain {
    pub root_subject: String,
    pub steps: Vec<DelegationStep>,
    pub active_subject: String,
}
\end{lstlisting}

Setelah \texttt{EffectiveCapability} dan \texttt{DelegationChain} terbentuk, \textit{Server} PRE memeriksa waktu kedaluwarsa token. Token ditolak apabila nilai \texttt{expires\_before} pada \texttt{EffectiveCapability} telah melewati waktu saat \textit{request} diproses. \textit{Server} PRE juga memeriksa batas kedalaman delegasi sudah melewati \texttt{max\_delegation\_depth} dengan menghitung jumlah \texttt{steps} pada \texttt{DelegationChain} untuk memastikan jumlah delegasi yang digunakan tidak melebihi batas yang ditentukan Apabila ada yang tidak terpenuhi, \textit{request} langsung ditolak.

\subsubsection{Pemeriksaan Status Revokasi Token}
\label{subsec:pemeriksaan-status-revokasi-token}

Pemeriksaan revokasi token dilakukan melalui Redis. \textit{Server} PRE menghitung \textit{hash} dari token yang sedang digunakan, kemudian membentuk kunci revokasi dengan format \texttt{revoked:token:<token\_hash>}. Selain token yang sedang digunakan, \textit{Server} PRE juga memeriksa \textit{hash} token induk yang tersimpan dalam \textit{caveat} \texttt{parent\_token\_hash}. Dengan demikian, pencabutan token induk juga dapat menyebabkan token hasil delegasi atau atenuasi ditolak.

\subsubsection{Verifikasi \textit{Role} dan \textit{Subrole} melalui \textit{Move Smart Contract}}
\label{subsec:validasi-role-subrole-move-smart-contract}

Setelah token dinyatakan valid dan tidak dicabut, \textit{Server} PRE memeriksa identitas kelembagaan aktor aktif melalui \textit{smart contract} \texttt{get\_hospital\_personnel\_auth\_info} menggunakan alamat IOTA aktor aktif. Fungsi ini mengembalikan \textit{role} dan \textit{subrole} personel yang tersimpan pada penyimpanan \textit{on-chain}. Eksistensi aktor aktif dalam sistem juga tervalidasi melalui pemeriksaan ini.

\subsubsection{Verifikasi Konteks Segmen}
\label{subsec:verifikasi-konteks-segmen}

Validasi konteks segmen dilakukan pada \textit{handler} operasi baca dan tulis. Pada operasi baca, \textit{handler} mengambil metadata segmen dari IOTA, kemudian mendekode metadata JSON yang tersimpan dalam format base64. Pada operasi tulis, \textit{handler} membentuk metadata \textit{preview} dari \textit{payload} segmen sebelum data disimpan. Metadata ini digunakan untuk membangun \texttt{SegmentAccessContext} yang berisi \texttt{segment\_id}, \texttt{patient\_address}, \texttt{related\_rme\_id}, \texttt{dataset\_category}, dan \texttt{function\_category}.

Pemeriksaan akses segmen dilakukan dengan membandingkan konteks segmen terhadap \texttt{EffectiveCapability}. Pada operasi baca, \textit{Server} PRE memastikan bahwa \texttt{patient\_address} sesuai dengan token, serta \texttt{dataset\_category} dan \texttt{function\_category} berada dalam cakupan \texttt{read\_dataset\_in} dan \texttt{read\_function\_in}. Pada operasi tulis, \textit{Server} PRE memastikan bahwa \texttt{patient\_address} sesuai, \texttt{related\_rme\_id} sesuai, serta \texttt{dataset\_category} dan \texttt{function\_category} berada dalam cakupan \texttt{write\_dataset\_in} dan \texttt{write\_function\_in}.

Untuk operasi tulis segmen, \textit{Server} PRE juga melakukan pemeriksaan tambahan berdasarkan \textit{role} dan \textit{subrole} aktor aktif. Aktor dengan \textit{role} \texttt{MedicalPersonnel} hanya dapat menulis kategori dataset yang sesuai dengan \textit{subrole}-nya. Dokter dan perawat dapat menulis dataset \texttt{RAWAT\_JALAN} dan \texttt{RAWAT\_INAP}, petugas laboratorium dapat menulis dataset \texttt{LABORATORIUM}, sedangkan apoteker dapat menulis dataset \texttt{APOTEK}. Aktor dengan \textit{role} \texttt{AdministrativePersonnel} hanya dapat membuat segmen dengan \textit{function category} \texttt{ADMINISTRATIVE\_GENERAL}. Selain itu, pada operasi tulis, \textit{Server} PRE memeriksa kesesuaian \texttt{hospital\_cid} pada token dengan \texttt{hospital\_cid} pada segmen yang akan disimpan.

\subsubsection{Verifikasi Bukti Kepemilikan \textit{Wallet}}
\label{subsec:verifikasi-bukti-kepemilikan-wallet}

Setelah konteks akses dinyatakan sesuai, \textit{Server} PRE memverifikasi bukti kepemilikan \textit{wallet} dari aktor aktif. Bukti ini dikirimkan melalui \textit{header} \texttt{x-decmed-wallet-signature}, sedangkan waktu pembuktian dikirimkan melalui \textit{header} \texttt{x-decmed-wallet-timestamp}. Tanda tangan dibuat terhadap representasi \texttt{WalletProofContext}. Struktur \texttt{WalletProofContext} ditunjukkan pada Kode \ref{lst:struktur-walletproofcontext}. \textit{Server} PRE kemudian memverifikasi \textit{wallet signature} tersebut terhadap alamat IOTA aktor aktif yang diperoleh dari \texttt{DelegationChain}. Timestamp pada konteks pembuktian juga divalidasi agar berada dalam rentang waktu yang dapat diterima. Pada implementasi ini, waktu yang diterima adalah rentang waktu 5 menit.

\begin{lstlisting}[language=Rust, caption={Struktur \texttt{WalletProofContext} sebagai konteks verifikasi \textit{wallet signature}}, label={lst:struktur-walletproofcontext}]
pub struct WalletProofContext {
	pub token_id: String,       // ID Macaroon
	pub patient_address: String,// Alamat pasien
	pub related_rme_id: String, // ID episode RME
	pub operation: AccessMode,  // Read / Update
	pub segment_id: String,     // ID segmen yang diakses
	pub dataset_category: DatasetCategory,
	pub function_category: FunctionCategory,
	pub timestamp: String,		// YYYY-MM-DD HH:MM:SS
}
\end{lstlisting}

Pada jalur operasi baca, apabila \textit{wallet signature} belum disertakan, \textit{Server} PRE mengembalikan respons \texttt{401 UNAUTHORIZED} yang berisi \texttt{proof\_context}. \textit{Client} kemudian dapat menandatangani konteks tersebut menggunakan kunci privat \textit{wallet} aktor aktif dan mengirimkan ulang \textit{request} dengan menyertakan \texttt{x-decmed-wallet-signature} serta \texttt{x-decmed-wallet-timestamp}.

\subsubsection{Eksekusi Operasi \textit{Re-encryption}, IPFS, dan IOTA}
\label{subsec:eksekusi-operasi-re-encryption-ipfs-iota}

Setelah seluruh pemeriksaan otorisasi selesai, \textit{Server} PRE melanjutkan permintaan ke operasi utama sesuai jenis aksi yang diminta. Pada tahap ini, \textit{Server} PRE berperan sebagai penghubung antara data terenkripsi yang tersimpan di IPFS, metadata segmen yang tersimpan di IOTA, dan mekanisme \textit{proxy re-encryption}.

Pada operasi baca, \textit{Server} PRE terlebih dahulu mengambil metadata RME dari IOTA berdasarkan pasien dan indeks data yang diminta. Metadata tersebut digunakan untuk memperoleh informasi penyimpanan data, seperti \texttt{ipfs\_cid}, \texttt{capsule} dan \texttt{enc\_key\_and\_nonce}. Setelah metadata diperoleh, \textit{Server} PRE mengambil data medis terenkripsi dari IPFS menggunakan \texttt{ipfs\_cid} pada metadata. Selanjutnya, \textit{Server} PRE memuat material PRE yang diperlukan, yaitu \texttt{kfrag}, kunci publik pasien, kunci publik penanda tangan, dan kunci publik data RME.

Proses \textit{re-encryption} dilakukan dengan Umbral PRE terhadap \textit{capsule} yang terkait dengan data medis. Hasil proses ini adalah \texttt{cfrag}, yaitu fragmen kapsul yang memungkinkan \textit{client} penerima membuka data terenkripsi.

\textit{Response} yang dikembalikan ke \textit{client} untuk operasi baca berisi data medis terenkripsi, \texttt{cfrag} hasil \textit{re-encryption}, informasi kunci publik yang diperlukan oleh \textit{client}, \textit{capsule}, \textit{encrypted key and nonce}, dan waktu pembuatan data.

Pada operasi tulis, \textit{client} mengirimkan segmen RME yang sudah terenkripsi. \textit{Server} PRE menyimpan \textit{payload} terenkripsi ke IPFS. Setelah penyimpanan berhasil, IPFS mengembalikan CID sebagai penunjuk terhadap data terenkripsi tersebut.

CID dari IPFS kemudian dimasukkan ke dalam metadata segmen final. Metadata ini memuat informasi sesuai dengan skema pada Lampiran \ref{lampiran:skema-penyimpanan-rme}. Setelah metadata final terbentuk, \textit{Server} PRE menyimpannya ke IOTA melalui \textit{smart contract} \texttt{create\_medical\_record\_segment}.

\textit{Response} yang dikembalikan ke \textit{client} untuk operasi tulis berisi ringkasan metadata segmen yang baru dibuat, yaitu \texttt{segment\_id}, \texttt{related\_rme\_id}, \texttt{dataset\_category}, \texttt{function\_category}, \texttt{ipfs\_cid}, \texttt{integrity\_hash}, dan \texttt{created\_at}. \textit{Response} ini menjadi konfirmasi bahwa data terenkripsi telah tersimpan di IPFS dan metadata referensinya telah dicatat di IOTA.

Pada operasi \textit{list}, \textit{Server} PRE mengambil kumpulan metadata dari IOTA secara bertahap menggunakan mekanisme \textit{pagination}. Setiap metadata yang berhasil di-\textit{decode} RME kemudian diubah menjadi \textit{list item} yang berisi informasi ringkas, yaitu \texttt{index}, \texttt{segment\_id}, \texttt{related\_rme\_id}, alamat pasien, \texttt{dataset category}, \texttt{function category}, \texttt{ipfs\_cid}, \texttt{created\_at} dan \texttt{author\_address}. Daftar metadata tersebut kemudian diproses agar hanya menampilkan metadata aktif sesuai dengan \texttt{EffectiveCapability}. Setelah itu, hasilnya diurutkan dan dipotong sesuai parameter \texttt{cursor} dan \texttt{limit} dari \textit{request}. \textit{Response} yang dikembalikan berisi daftar item metadata aktif dan \texttt{next\_cursor} apabila masih terdapat data lanjutan yang dapat diambil oleh \textit{client}.

Dengan alur ini, IPFS digunakan untuk menyimpan isi data RME terenkripsi, sedangkan IOTA digunakan untuk menyimpan metadata referensial yang diperlukan untuk pencarian, validasi konteks, dan navigasi data. \textit{Server} PRE menghubungkan keduanya dengan mekanisme Umbral PRE pada operasi baca, pencatatan metadata pada operasi tulis, dan penyusunan daftar metadata aktif pada operasi \textit{list}. \subsection{Implementasi \textit{Smart Contract} IOTA Move}
\label{subsec:implementasi-smart-contract-iota-move}



\textit{Smart contract} IOTA Move diimplementasikan untuk menyimpan metadata yang diperlukan oleh mekanisme segmentasi dan kontrol akses. Data RME tetap disimpan secara \textit{off-chain} pada IPFS, sedangkan IOTA menyimpan informasi yang dibutuhkan untuk menemukan, memverifikasi, dan menghubungkan segmen data yang diakses dengan pemiliknya.

Pada implementasi Tugas Akhir ini, struktur penyimpanan metadata klinis di Move masih bersifat flat melalui \texttt{medical\_metadata: TableVec<PatientMedicalMetadata>}. Struktur \texttt{PatientMedicalMetadata} menyimpan \texttt{index} dan \texttt{metadata}. Informasi segmentasi seperti \texttt{dataset\_category}, \texttt{function\_category}, \texttt{related\_rme\_id}, CID IPFS, \texttt{integrity\_hash}, dan \texttt{capsule} disimpan di dalam \texttt{metadata} sebagai JSON base64. Dengan cara ini, \textit{Server} PRE dapat mengambil metadata segmen dari IOTA sebelum memutuskan apakah token Macaroon berhak mengakses segmen tersebut.

Fungsi pada \textit{smart contract} disesuaikan untuk mendukung pembuatan metadata RME, penambahan segmen baru, pembaruan referensi segmen, pembacaan metadata, penyimpanan metadata akses personel, serta pencatatan audit delegasi. Operasi yang mengubah metadata hanya dapat dilakukan oleh aktor yang memiliki akses tulis yang sah pada lapisan aplikasi dan \textit{smart contract}. Rincian \textit{module} dan \textit{method} pada implementasi \textit{smart contract} IOTA Move disajikan pada Lampiran \ref{lampiran:implementasi-smart-contract-sistem-decmed}.

 \subsection{Implementasi \textit{Patient Client}}
\label{subsec:implementasi-patient-client}

\textit{Patient Client} disesuaikan untuk mendukung skenario penerbitan token awal oleh pasien, pembacaan RME, peninjauan audit akses, dan pencabutan akses. \textit{Patient Client} tidak hanya berfungsi sebagai tampilan pengguna, tetapi juga sebagai titik pembentukan konteks permintaan pasien yang dikirimkan ke \textit{Server} PRE. Tangkapan layar implementasi disajikan pada Lampiran \ref{lampiran:implementasi-patient-client}.

Terdapat beberapa bagian dari \textit{Patient Client} yang sudah ada dari implementasi DecMed sebelumnya, seperti halaman register, halaman login, halaman \textit{Home}, halaman Profil, Halaman log akses, halaman \textit{scan} QR, dan halaman detail RME. Bagian-bagian tersebut diadaptasi dan disesuaikan untuk mendukung implementasi Tugas Akhir ini. \textit{Patient Client} mengalami penyesuaian dan tambahan halaman baru pada implementasi Tugas Akhir ini untuk beberapa halaman berikut,

\subsubsection{Halaman \textit{Home} (\textit{List} RME)}
\label{subsubsec:halaman-home-daftar-rme}

Halaman \textit{home} merupakan halaman utama yang pertama kali dilihat pasien setelah berhasil masuk ke dalam aplikasi. Halaman ini menampilkan seluruh RME pasien dalam bentuk daftar \textit{card} yang tersusun secara vertikal. Setiap \textit{card} merepresentasikan satu RME unik dan menampilkan tanggal pembuatan RME dan ID RME. Daftar RME diurutkan berdasarkan tanggal terbaru, sehingga RME yang paling baru selalu muncul di bagian atas. Setiap \textit{card} dapat diklik, di mana klik tersebut akan mengarahkan pasien menuju halaman detail RME yang bersangkutan.

Tampilan dari halaman \textit{home} dapat dilihat pada
Gambar \ref{fig:halaman-home-daftar-rme}.

\subsubsection{Halaman Detail RME}
\label{subsubsec:halaman-detail-rme}

Halaman detail RME ditampilkan setelah pasien memilih salah satu \textit{card} dari daftar RME pada halaman \textit{home}. Halaman ini menampilkan detail lengkap dari sebuah RME per dataset. Bagian atas halaman menampilkan ID dari RME. Di bawahnya, terdapat \textit{tab} daftar \textit{dataset category} RME tersebut. \textit{Tab} yang terpilih menentukan entri data RME yang ditampilkan dalam bentuk \textit{card} yang memuat nama \textit{function category}. Pasien juga dapat kembali ke halaman \textit{home} melalui tautan navigasi balik yang tersedia di bagian atas halaman.

Tampilan dari halaman detail RME dapat dilihat pada Gambar \ref{fig:halaman-detail-rme}.

\subsubsection{Halaman Audit Delegasi}
\label{subsubsec:halaman-audit-delegasi}

Halaman audit delegasi menampilkan informasi mengenai seluruh rantai delegasi akses yang terbentuk dari token Macaroon yang telah diterbitkan oleh pasien. Setiap entri pada halaman ini merepresentasikan satu rantai delegasi dan menampilkan identitas personel rumah sakit yang menjadi penerima akses awal, jenis akses yang diberikan, ID RME yang terkait, serta tanggal kedaluwarsa dari akses awal tersebut. Status keseluruhan rantai delegasi juga ditampilkan dalam bentuk label berwarna, yakni \textit{Active} berwarna hijau, \textit{Expired} berwarna kuning, atau \textit{Revoked} berwarna merah. Di dalam setiap entri rantai delegasi, terdapat subdaftar yang merinci setiap delegasi, yaitu pasangan delegasi dari satu personel ke personel lainnya, beserta informasi kedalaman delegasi, tanggal kedaluwarsa, dan status aktif atau dicabutnya delegasi tersebut.

Tampilan dari halaman detail audit delegasi RME dapat dilihat pada Gambar \ref{fig:halaman-detail-audit-delegasi}.

\subsubsection{Halaman Detail \textit{Scan} QR}
\label{subsubsec:halaman-detail-scan-qr}

Halaman \textit{scan} QR merupakan halaman yang digunakan pasien untuk memulai proses pemberian akses kepada personel rumah sakit. Halaman ini menyediakan antarmuka unggah berkas gambar QR yang diperoleh dari rumah sakit, di mana pasien dapat memilih berkas gambar dalam format PNG atau JPEG. Setelah berkas QR diunggah dan diproses, sistem akan memvalidasi isi QR dan mengidentifikasi identitas personel rumah sakit yang terkandung di dalamnya. Apabila QR valid, sebuah dialog konfirmasi akan muncul yang menampilkan nama personel dan nama rumah sakit tujuan pemberian akses. Pada dialog konfirmasi ini, pasien dapat mengonfigurasi cakupan akses yang akan diberikan, meliputi jenis layanan medis berupa rawat jalan atau
rawat inap, mode akses berupa \textit{Read}, \textit{WRITE}, atau keduanya, serta tanggal dan waktu kedaluwarsa akses. Setelah mengonfirmasi pilihan tersebut, pasien akan diarahkan ke dialog masukan PIN enam digit sebagai verifikasi identitas sebelum token Macaroon awal diterbitkan dan akses resmi diberikan kepada personel rumah sakit yang dituju.

Tampilan dari halaman \textit{scan} QR dan halaman konfirmasi akses awal dapat dilihat pada Gambar \ref{fig:halaman-detail-scan-qr} dan Gambar \ref{fig:halaman-konfirmasi-akses-awal}.
 \subsection{Implementasi \textit{Hospital Client}}
\label{subsec:implementasi-hospital-client}

\textit{Hospital Client} disesuaikan untuk mendukung skenario penggunaan token oleh tenaga medis, pembuatan token turunan, dan akses terhadap segmen data RME. Tangkapan layar implementasi disajikan pada Lampiran \ref{lampiran:implementasi-hospital-client}.

\subsubsection{Halaman \textit{Dashboard} (\textit{List} RME)}
\label{subsubsec:halaman-dashboard-hospital-client}

Halaman \textit{dashboard} merupakan halaman utama yang ditampilkan setelah aktor rumah sakit berhasil masuk ke sistem. Konten halaman ini bersifat \textit{role-based}, sehingga tampilan yang disajikan berbeda-beda bergantung pada peran aktor yang sedang aktif. Untuk peran \texttt{MedicalPersonnel}, halaman ini menampilkan dua tab yang dapat dipilih secara bergantian, yaitu tab \textit{Read} dan tab \textit{Write}.

Tab \textit{Read} menampilkan daftar pasien yang token aksesnya memiliki kapabilitas baca, sedangkan tab \textit{Write} menampilkan daftar pasien yang token aksesnya memiliki kapabilitas tulis. Setiap entri dalam kedua tab tersebut menampilkan nama pasien dan dapat diklik untuk menavigasi ke halaman berikutnya. Untuk entri pada tab \textit{Read}, klik akan membawa aktor menuju halaman daftar RME pasien yang terkait dengan token baca. Untuk entri pada tab \textit{Write}, klik akan membawa aktor langsung menuju halaman formulir penulisan segmen RME, dengan konteks token tulis yang sesuai.

Tampilan dari halaman \textit{dashboard} untuk peran \texttt{MedicalPersonnel} dapat dilihat pada Gambar \ref{fig:halaman-dashboard-medical-personnel}.

\subsubsection{Halaman Delegasi Akses}
\label{subsubsec:halaman-delegasi-akses-hospital-client}

Halaman delegasi akses digunakan oleh aktor yang memiliki hak untuk mendelegasikan token aksesnya kepada personel lain di fasyankes. Pada atas halaman ini terdapat ringkasan hak akses yang saat ini dimiliki oleh aktor, berupa daftar pasien. Di bawahnya adalah \textit{form} delegasi yang aktif apabila terdapat setidaknya satu akses pasien yang tersedia. Pada formulir ini, aktor dapat memilih pasien sebagai hak akses yang akan didelegasikan, menentukan mode delegasi berupa \textit{Read}, \textit{Write}, atau keduanya, memilih \textit{preset} \textit{role} penerima delegasi seperti dokter atau perawat, serta menetapkan tanggal dan waktu kedaluwarsa delegasi. Sistem secara otomatis membatasi waktu kedaluwarsa maksimum token induk yang sedang digunakan. Formulir juga menyediakan pratinjau cakupan baca dan tulis berupa daftar \textit{dataset category} dan \textit{function category} yang dapat dipilih atau dihilangkan secara individual sesuai kapabilitas token induk. Terakhir, aktor memilih penerima delegasi dari daftar kandidat personel yang tersedia di fasyankes, kemudian menekan tombol \textit{Delegate} untuk memulai proses delegasi.

Tampilan halaman delegasi akses dapat dilihat pada Gambar \ref{fig:halaman-delegasi-akses-bagian-1} dan Gambar \ref{fig:halaman-delegasi-akses-bagian-2}.

\subsubsection{Halaman Daftar RME Pasien}
\label{subsubsec:halaman-daftar-emr-pasien-hospital-client}

Halaman daftar RME pasien menampilkan seluruh RME yang dimiliki oleh pasien tertentu dan dapat diakses oleh aktor sesuai token yang dibawa. Bagian atas halaman menampilkan nama pasien sebagai judul halaman. Di bawahnya, terdapat daftar RME dalam bentuk \textit{card} yang masing-masing memuat ID RME, tanggal pembuatannya, serta kumpulan label \textit{badge} yang merepresentasikan kategori dataset yang terdapat dalam RME tersebut. Setiap \textit{card} RME dapat diklik dan akan mengarahkan aktor menuju halaman detail RME dari RME yang dipilih, dengan konteks token akses yang diteruskan melalui \textit{query string}. Apabila tidak terdapat RME yang dapat diakses, halaman menampilkan pesan kosong.

Tampilan halaman daftar RME pasien dapat dilihat pada Gambar \ref{fig:halaman-daftar-rme-pasien-hospital-client}.

\subsubsection{Halaman Baca Detail RME}
\label{subsubsec:halaman-baca-detail-rme-hospital-client}

Halaman detail RME menampilkan isi lengkap dari sebuah RME yang dipilih oleh aktor. Bagian atas halaman memuat ID RME sebagai identifikasi, diikuti oleh tab navigasi yang merepresentasikan masing-masing \textit{dataset category} yang terdapat dalam RME tersebut. Aktor dapat berpindah antarklasifikasi dataset dengan mengklik tab yang sesuai. Pada setiap tab dataset, ditampilkan daftar kartu segmen yang memuat nama \textit{function category} dari segmen tersebut. Setiap kartu segmen dilengkapi dengan tombol ikon informasi yang dapat diklik untuk menampilkan atau menyembunyikan metadata. Konten payload segmen dimuat secara \textit{on-demand} dan ditampilkan dalam bentuk \textit{textarea} hanya-baca, atau dalam format tabel data
khusus apabila segmen termasuk \textit{function category} \texttt{ADMINISTRATIVE\_GENERAL}.

Tampilan halaman detail RME dapat dilihat pada Gambar \ref{fig:halaman-detail-rme-hospital-client}.

\subsubsection{Halaman Tulis Segmen RME}
\label{subsubsec:halaman-penulisan-segmen-rme-hospital-client}

Halaman penulisan segmen RME digunakan ketika aktor memiliki kapabilitas tulis terhadap RME pasien tertentu. Halaman ini dapat diakses dari daftar pasien pada tab \textit{Write} di halaman \textit{dashboard}.

Bagian atas halaman menampilkan nama pasien sebagai konteks identifikasi. Di bawahnya, terdapat menu pilihan \textit{dataset category} yang hanya memuat kategori yang diizinkan oleh token tulis yang aktif. Halaman secara otomatis menampilkan daftar formulir \textit{textarea} yang bersesuaian dengan \textit{function category} yang diizinkan dalam dataset tersebut, di mana setiap \textit{textarea} merepresentasikan satu jenis data klinis yang dapat diisikan. Aktor hanya perlu mengisi \textit{textarea} untuk \textit{function category} yang ingin dituliskan, karena hanya entri yang terisi yang akan dikirimkan. Setelah mengisi data yang diperlukan, aktor menekan tombol \textit{Write} untuk menyimpan segmen-segmen baru ke dalam RME pasien.

Tampilan halaman penulisan segmen RME dapat dilihat pada Gambar \ref{fig:halaman-penulisan-segmen-rme-hospital-client}.  \section{Pengujian Sistem}
\label{sec:pengujian-sistem}

Pengujian sistem dilakukan untuk memastikan implementasi memenuhi kebutuhan yang sudah didefinisikan pada Subbab \ref{subsec:kebutuhan-sistem}.

\subsection{Strategi Pengujian}
\label{subsec:strategi-pengujian}

Strategi pengujian disusun berdasarkan kebutuhan sistem pada Tabel \ref{tab:kebutuhan-fungsional-tambahan-decmed} dan \ref{tab:kebutuhan-nonfungsional-tambahan-decmed}. Pengujian otomatis diberi ID \texttt{PU-TA}, sedangkan pengujian skenario manual diberi ID \texttt{PS-TA}.

Pengujian otomatis digunakan untuk memeriksa fungsi teknis komponen / modul tertentu, sedangkan pengujian skenario manual digunakan untuk memeriksa perilaku sistem dari sudut pandang aktor layanan kesehatan. Pengujian otomatis mencakup \texttt{cargo test} pada komponen yang relevan. Pengujian skenario manual disusun dengan pendekatan positif. Skenario positif memastikan akses yang sah dapat diproses dengan baik sesuai ekspektasi.

Kriteria keberhasilan pengujian ditentukan berdasarkan hasil dari pengujian otomatis dan skenario manual.
 \subsection{Pengujian Skenario Akses dan Delegasi}
\label{subsec:pengujian-skenario-akses-delegasi}

Pengujian skenario akses dan delegasi dilakukan secara manual untuk memastikan perilaku sistem sesuai dengan rancangan alur akses dan delegasi pada Subbab \ref{subsec:rancangan-alur-akses-dan-delegasi-macaroon}.

Ringkasan pengujian skenario akses dan delegasi ditunjukkan pada Tabel \ref{tab:pengujian-skenario-akses-delegasi}. Struktur skenario pada tabel tersebut dibagi sebagai berikut: \texttt{PS-TA-01} sampai \texttt{PS-TA-11} menguji alur Rawat Jalan dan \texttt{PS-TA-12} sampai \texttt{PS-TA-22} menguji alur Rawat Inap.
Pengujian untuk skenario Rawat Jalan dan Rawat Inap, disimulasikan sebagai satu episode pelayanan. Perbedaan antara kedua alur adalah keberadaan segmen data \texttt{PERENCANAAN\_PEMULANGAN} yang hanya terdapat pada alur Rawat Inap. Delegasi akses yang dilakukan mengikuti \textit{preset} yang sudah terdefinisi, menghasilkan token turunan dengan \textit{caveats} sesuai dengan Lampiran \ref{lampiran:contoh-caveat-macaroon}.

\begingroup
\small
\renewcommand{\arraystretch}{1.2}
\begin{longtable}{|p{0.12\textwidth}|p{0.43\textwidth}|p{0.40\textwidth}|}
	\caption{Pengujian Skenario Akses dan Delegasi pada Alur Rawat Jalan dan Rawat Inap}
	\label{tab:pengujian-skenario-akses-delegasi}                                                                                \\
	\hline
	\textbf{ID}                                                                        & \textbf{Skenario} & \textbf{Ekspektasi} \\ \hline
	\endfirsthead

	\caption[]{Pengujian Skenario Akses dan Delegasi pada Alur Rawat Jalan dan Rawat Inap (lanjutan)}                            \\
	\hline
	\textbf{ID}                                                                        & \textbf{Skenario} & \textbf{Ekspektasi} \\ \hline
	\endhead

	\texttt{PS-TA-01}                                                                  &
	Pasien memberikan akses awal Rawat Jalan kepada petugas RM.                        &
	Token awal Rawat Jalan berhasil diterbitkan.                                                                                 \\ \hline
	\texttt{PS-TA-02}                                                                  &
	Petugas RM mendelegasikan akses pemeriksaan awal Rawat Jalan kepada perawat.       &
	Token turunan perawat berhasil dibuat.                                                                                       \\ \hline

	\texttt{PS-TA-03}                                                                  &
	Petugas RM mendelegasikan akses klinis Rawat Jalan kepada dokter.                  &
	Token turunan dokter berhasil dibuat.                                                                                        \\ \hline


	\texttt{PS-TA-04}                                                                  &
	Perawat mengisi hasil pemeriksaan awal Rawat Jalan.                                &
	Data pemeriksaan awal berhasil disimpan.                                                                                     \\ \hline

	\texttt{PS-TA-05}                                                                  &
	Dokter membaca pemeriksaan awal dan mengisi data klinis Rawat Jalan.               &
	Data klinis Rawat Jalan berhasil dibaca dan disimpan.                                                                        \\ \hline

	\texttt{PS-TA-06}                                                                  &
	Dokter mendelegasikan akses laboratorium Rawat Jalan.                              &
	Token turunan laboratorium berhasil dibuat.                                                                                  \\ \hline

	\texttt{PS-TA-07}                                                                  &
	Petugas laboratorium membaca pemeriksaan penunjang dan mengisi hasil laboratorium. &
	Data pemeriksaan penunjang berhasil dibaca dan hasil laboratorium disimpan.                                                  \\ \hline

	\texttt{PS-TA-08}                                                                  &
	Dokter membaca hasil laboratorium Rawat Jalan.                                     &
	Hasil laboratorium berhasil dibaca oleh dokter.                                                                              \\ \hline

	\texttt{PS-TA-09}                                                                  &
	Dokter mendelegasikan akses apotek Rawat Jalan.                                    &
	Token turunan apoteker berhasil dibuat.                                                                                      \\ \hline

	\texttt{PS-TA-10}                                                                  &
	Apoteker membaca resep dan mengisi status pelayanan obat.                          &
	Data resep dan pelayanan obat berhasil diproses.                                                                             \\ \hline

	\texttt{PS-TA-11}                                                                  &
	Pasien membuka riwayat delegasi episode Rawat Jalan.                               &
	Riwayat delegasi episode Rawat Jalan ditampilkan.                                                                            \\ \hline

	\texttt{PS-TA-12}                                                                  &
	Pasien memberikan akses awal Rawat Inap kepada petugas RM.                         &
	Token awal Rawat Inap berhasil diterbitkan.                                                                                  \\ \hline

	\texttt{PS-TA-13}                                                                  &
	Petugas RM mendelegasikan akses pemeriksaan awal Rawat Inap kepada perawat.        &
	Token turunan perawat berhasil dibuat.                                                                                       \\ \hline

	\texttt{PS-TA-14}                                                                  &
	Petugas RM mendelegasikan akses klinis Rawat Inap kepada dokter.                   &
	Token turunan dokter berhasil dibuat.                                                                                        \\ \hline

	\texttt{PS-TA-15}                                                                  &
	Perawat mengisi hasil pemeriksaan awal Rawat Inap.                                 &
	Data pemeriksaan awal berhasil disimpan.                                                                                     \\ \hline

	\texttt{PS-TA-16}                                                                  &
	Dokter membaca pemeriksaan awal dan mengisi data klinis Rawat Inap.                &
	Data klinis Rawat Inap berhasil dibaca dan disimpan.                                                                         \\ \hline

	\texttt{PS-TA-17}                                                                  &
	Dokter mendelegasikan akses laboratorium Rawat Inap.                               &
	Token turunan laboratorium berhasil dibuat.                                                                                  \\ \hline

	\texttt{PS-TA-18}                                                                  &
	Petugas laboratorium membaca pemeriksaan penunjang dan mengisi hasil laboratorium. &
	Data pemeriksaan penunjang berhasil dibaca dan hasil laboratorium disimpan.                                                  \\ \hline

	\texttt{PS-TA-19}                                                                  &
	Dokter membaca hasil laboratorium Rawat Inap.                                      &
	Hasil laboratorium berhasil dibaca oleh dokter.                                                                              \\ \hline

	\texttt{PS-TA-20}                                                                  &
	Dokter mendelegasikan akses apotek Rawat Inap.                                     &
	Token turunan apoteker berhasil dibuat.                                                                                      \\ \hline

	\texttt{PS-TA-21}                                                                  &
	Apoteker membaca resep dan mengisi status pelayanan obat.                          &
	Data resep dan pelayanan obat berhasil diproses.                                                                             \\ \hline

	\texttt{PS-TA-22}                                                                  &
	Pasien membuka riwayat delegasi episode Rawat Inap.                                &
	Riwayat akses dan delegasi Rawat Inap ditampilkan.                                                                           \\ \hline







\end{longtable}
\endgroup \subsection{Pengujian Unit dan Integrasi}
\label{subsec:pengujian-unit-integrasi}

Pengujian unit dan integrasi dilakukan terhadap komponen yang menangani token, segmentasi data, validasi akses, PRE. Pengujian unit dipetakan terhadap kebutuhan terkait. Ringkasan target pengujian unit dan integrasi ditunjukkan pada Tabel \ref{tab:target-pengujian-unit-integrasi}.

\begingroup
\small
\begin{longtable}{|p{0.12\textwidth}|p{0.14\textwidth}|p{0.54\textwidth}|p{0.18\textwidth}|}
	\caption{Target Pengujian Unit dan Integrasi}
	\label{tab:target-pengujian-unit-integrasi}                                                                                                                                                                                                                                                                                                                                       \\
	\hline
	\textbf{ID}                                                                                                                                                                                                                                                                                  & \textbf{Jenis} & \textbf{Fokus Uji dan Kondisi Lulus} & \textbf{Kebutuhan Terkait} \\ \hline
	\endfirsthead

	\caption[]{Target Pengujian Unit dan Integrasi (lanjutan)}                                                                                                                                                                                                                                                                                                                        \\
	\hline
	\textbf{ID}                                                                                                                                                                                                                                                                                  & \textbf{Jenis} & \textbf{Fokus Uji dan Kondisi Lulus} & \textbf{Kebutuhan Terkait} \\ \hline
	\endhead

	\texttt{PU-TA-01}                                                                                                                                                                                                                                                                            &
	Unit                                                                                                                                                                                                                                                                                         &
	Pembentukan dan verifikasi Macaroon valid. Token yang dibentuk dengan \textit{caveats} yang sesuai dapat diverifikasi dan diterima.                                                                                                                                                          &
	\texttt{KF-TA-01}, \texttt{KNF-TA-02}                                                                                                                                                                                                                                                                                                                                             \\ \hline

	\texttt{PU-TA-02}                                                                                                                                                                                                                                                                            &
	Unit                                                                                                                                                                                                                                                                                         &
	Integritas token Macaroon. Token atau \textit{caveats} yang dimodifikasi di luar prosedur Macaroon harus gagal diverifikasi.                                                                                                                                                                 &
	\texttt{KNF-TA-02}, \texttt{KNF-TA-03}                                                                                                                                                                                                                                                                                                                                            \\ \hline

	\texttt{PU-TA-03}                                                                                                                                                                                                                                                                            &
	Unit                                                                                                                                                                                                                                                                                         &
	Perhitungan \textit{effective capability}. Cakupan token turunan dihitung sebagai irisan dari cakupan token induk dan batasan tambahan pada token turunan.                                                                                                                                   &
	\texttt{KF-TA-04}, \texttt{KF-TA-06}, \texttt{KNF-TA-02}                                                                                                                                                                                                                                                                                                                          \\ \hline

	\texttt{PU-TA-04}                                                                                                                                                                                                                                                                            &
	Unit                                                                                                                                                                                                                                                                                         &
	Validasi atenuasi akses. Token turunan yang memperluas \textit{dataset}, \textit{function category}, atau aksi melebihi token induk harus ditolak.                                                                                                                                           &
	\texttt{KF-TA-04}, \texttt{KF-TA-06}, \texttt{KNF-TA-03}                                                                                                                                                                                                                                                                                                                          \\ \hline

	\texttt{PU-TA-05}                                                                                                                                                                                                                                                                            &
	Unit                                                                                                                                                                                                                                                                                         &
	Validasi masa berlaku token. Token dengan nilai \texttt{expires\_before} yang sudah terlewati harus ditolak.                                                                                                                                                                                 &
	\texttt{KNF-TA-03}                                                                                                                                                                                                                                                                                                                                                                \\ \hline

	\texttt{PU-TA-06}                                                                                                                                                                                                                                                                            &
	Unit                                                                                                                                                                                                                                                                                         &
	Validasi kedalaman delegasi. Token turunan yang melewati batas \texttt{max\_delegation\_depth} harus ditolak.                                                                                                                                                                                &
	\texttt{KF-TA-04}, \texttt{KF-TA-06}, \texttt{KNF-TA-03}                                                                                                                                                                                                                                                                                                                          \\ \hline

	\texttt{PU-TA-07}                                                                                                                                                                                                                                                                            &
	Unit                                                                                                                                                                                                                                                                                         &
	Verifikasi bukti kepemilikan \textit{wallet}. \textit{Signature} yang dibuat oleh \textit{wallet} aktor aktif terhadap konteks permintaan yang sesuai harus diterima.                                                                                                                        &
	\texttt{KF-TA-10}, \texttt{KNF-TA-06}                                                                                                                                                                                                                                                                                                                                             \\ \hline

	\texttt{PU-TA-08}                                                                                                                                                                                                                                                                            &
	Unit                                                                                                                                                                                                                                                                                         &
	Validasi keterikatan \textit{wallet proof} terhadap konteks permintaan. Bukti kepemilikan \textit{wallet} harus gagal apabila \texttt{token\_id}, \texttt{patient\_address}, \texttt{related\_rme\_id}, \texttt{operation}, atau \texttt{segment\_id} diubah setelah proses penandatanganan. &
	\texttt{KF-TA-10}, \texttt{KNF-TA-03}, \texttt{KNF-TA-06}                                                                                                                                                                                                                                                                                                                         \\ \hline

	\texttt{PU-TA-09}                                                                                                                                                                                                                                                                            &
	Unit                                                                                                                                                                                                                                                                                         &
	Validasi metadata segmen RME. Metadata segmen harus memiliki \texttt{dataset\_category} dan \texttt{function\_category} yang valid agar dapat digunakan sebagai dasar pencocokan akses.                                                                                                      &
	\texttt{KF-TA-05}, \texttt{KF-TA-07}--\texttt{KF-TA-09}, \texttt{KNF-TA-01}                                                                                                                                                                                                                                                                                                       \\ \hline

	\texttt{PU-TA-10}                                                                                                                                                                                                                                                                            &
	Integrasi                                                                                                                                                                                                                                                                                    &
	Penegakan cakupan akses pada \textit{Server} PRE. Permintaan akses ke segmen di luar cakupan \textit{caveats} harus ditolak sebelum data terenkripsi diambil dan sebelum proses \textit{re-encryption}.                                                                                      &
	\texttt{KF-TA-05}, \texttt{KF-TA-07}--\texttt{KF-TA-09}, \texttt{KNF-TA-01}, \texttt{KNF-TA-03}                                                                                                                                                                                                                                                                                   \\ \hline

	\texttt{PU-TA-11}                                                                                                                                                                                                                                                                            &
	Integrasi                                                                                                                                                                                                                                                                                    &
	Validasi keterikatan token terhadap pasien dan episode RME. Token valid untuk pasien atau episode RME tertentu harus ditolak ketika digunakan untuk pasien atau episode RME yang berbeda.                                                                                                    &
	\texttt{KF-TA-10}, \texttt{KNF-TA-03}, \texttt{KNF-TA-06}                                                                                                                                                                                                                                                                                                                         \\ \hline

	\texttt{PU-TA-12}                                                                                                                                                                                                                                                                            &
	Integrasi                                                                                                                                                                                                                                                                                    &
	Penegakan jenis aksi. Token dengan hak baca tidak dapat digunakan untuk operasi tulis, dan token dengan cakupan aksi terbatas tidak dapat digunakan di luar aksi yang tercantum pada \textit{caveats}.                                                                                       &
	\texttt{KF-TA-05}, \texttt{KF-TA-07}--\texttt{KF-TA-09}, \texttt{KNF-TA-01}, \texttt{KNF-TA-03}                                                                                                                                                                                                                                                                                   \\ \hline

	\texttt{PU-TA-13}                                                                                                                                                                                                                                                                            &
	Integrasi                                                                                                                                                                                                                                                                                    &
	Validasi \textit{subrole} terhadap \textit{dataset} dan \textit{function category}. Aktor dengan \textit{subrole} yang tidak sesuai harus ditolak meskipun aktor tersebut termasuk kategori tenaga medis.                                                                                    &
	\texttt{KF-TA-05}, \texttt{KF-TA-07}--\texttt{KF-TA-09}, \texttt{KNF-TA-01}, \texttt{KNF-TA-03}                                                                                                                                                                                                                                                                                   \\ \hline

	\texttt{PU-TA-14}                                                                                                                                                                                                                                                                            &
	Integrasi                                                                                                                                                                                                                                                                                    &
	Validasi revokasi token. Token yang sudah dicabut harus ditolak oleh \textit{Server} PRE ketika status revokasi atau \textit{revocation key} terdeteksi.                                                                                                                                     &
	\texttt{KF-TA-02}, \texttt{KNF-TA-05}                                                                                                                                                                                                                                                                                                                                             \\ \hline

	\texttt{PU-TA-15}                                                                                                                                                                                                                                                                            &
	Integrasi                                                                                                                                                                                                                                                                                    &
	Validasi revokasi berantai. Token turunan harus ikut ditolak apabila token induk atau akses awal yang menjadi asal rantai delegasi telah dicabut.                                                                                                                                            &
	\texttt{KF-TA-02}, \texttt{KNF-TA-03}, \texttt{KNF-TA-05}                                                                                                                                                                                                                                                                                                                         \\ \hline

	\texttt{PU-TA-16}                                                                                                                                                                                                                                                                            &
	Integrasi                                                                                                                                                                                                                                                                                    &
	Validasi pencatatan metadata akses dan delegasi. Proses pemberian akses, delegasi, dan revokasi harus menghasilkan metadata yang dapat digunakan untuk menelusuri \textit{delegator}, \textit{delegatee}, aksi, waktu, status, dan relasi delegasi.                                          &
	\texttt{KF-TA-03}, \texttt{KNF-TA-04}                                                                                                                                                                                                                                                                                                                                             \\ \hline
\end{longtable}
\endgroup
  \section{Evaluasi}
\label{sec:evaluasi}

Evaluasi dilakukan dengan membandingkan hasil implementasi dan pengujian terhadap kebutuhan sistem serta rumusan masalah Tugas Akhir. Evaluasi ini digunakan untuk menilai apakah pendekatan Macaroons, segmentasi data RME, dan penegakan akses pada \textit{Server} PRE mampu menjawab keterbatasan kontrol akses pada DecMed awal.

\subsection{Evaluasi terhadap Kebutuhan Sistem}
\label{subsec:evaluasi-kebutuhan-sistem}

Berdasarkan kebutuhan sistem pada Subbab \ref{subsec:kebutuhan-sistem}, implementasi dievaluasi terhadap kebutuhan fungsional dan nonfungsional yang menjadi dasar pengembangan DecMed. Evaluasi menghubungkan kebutuhan sistem pada Subbab \ref{subsec:kebutuhan-sistem} dengan implementasi dan pengujian yang telah dilakukan. Setiap kebutuhan diperiksa berdasarkan komponen yang merealisasikannya dan skenario pengujian yang memverifikasi perilakunya. Pemetaan ini menunjukkan keterhubungan antara kebutuhan, implementasi, dan pengujian tanpa memisahkan ketiganya sebagai daftar yang berdiri sendiri.

Kebutuhan kontrol akses granular dipenuhi melalui penggunaan Macaroons dengan \textit{caveats}. Token tidak hanya menyatakan aktor penerima akses, tetapi juga membatasi pasien, RME, kategori dataset, kategori fungsi, hak baca atau tulis, masa berlaku, dan kedalaman delegasi. Kebutuhan segmentasi data RME dipenuhi dengan memecah data \textit{off-chain} menjadi segmen berdasarkan \textit{dataset category} dan \textit{function category}. Kebutuhan delegasi berantai dipenuhi melalui atenuasi token pada sisi \textit{client}, sedangkan penegakan akses dilakukan pada \textit{Server} PRE sebelum data terenkripsi diambil dan sebelum proses \textit{re-encryption} dijalankan.

Ringkasan pemetaan kebutuhan sistem terhadap implementasi dan pengujian ditunjukkan pada Tabel \ref{tab:evaluasi-kebutuhan-sistem}.

\begingroup
\small
\begin{longtable}{|p{0.14\textwidth}|p{0.32\textwidth}|p{0.30\textwidth}|p{0.14\textwidth}|}
	\caption{Evaluasi Kebutuhan Sistem terhadap Implementasi dan Pengujian}
	\label{tab:evaluasi-kebutuhan-sistem}                                                                                                                                                                                                                                                          \\
	\hline
	\textbf{Kebutuhan} & \textbf{Implementasi Terkait}                        & \textbf{Pengujian Terkait}                                                                                                                                                                   & \textbf{Status}     \\ \hline
	\endfirsthead
	\caption[]{Evaluasi Kebutuhan Sistem terhadap Implementasi dan Pengujian (lanjutan)}                                                                                                                                                                                                           \\
	\hline
	\textbf{Kebutuhan} & \textbf{Implementasi Terkait}                        & \textbf{Pengujian Terkait}                                                                                                                                                                   & \textbf{Status}     \\ \hline
	\endhead
	\texttt{KF-TA-01}  & \texttt{I-TA-02}, \texttt{I-TA-03}, \texttt{I-TA-05} & \texttt{PS-TA-01}                                                                                                                                                                            & Terpenuhi.          \\ \hline
	\texttt{KF-TA-02}  & \texttt{I-TA-03}, \texttt{I-TA-04}, \texttt{I-TA-05} & \texttt{PU-TA-10}, \texttt{PS-TA-15}                                                                                                                                                         & Terpenuhi.          \\ \hline
	\texttt{KF-TA-03}  & \texttt{I-TA-04}, \texttt{I-TA-05}                   & \texttt{PS-TA-01}, \texttt{PS-TA-02}, \texttt{PS-TA-06}, \texttt{PS-TA-09}, \texttt{PS-TA-15}                                                                                                & Terpenuhi sebagian. \\ \hline
	\texttt{KF-TA-04}  & \texttt{I-TA-02}, \texttt{I-TA-04}, \texttt{I-TA-06} & \texttt{PS-TA-02}, \texttt{PS-TA-05}                                                                                                                                                         & Terpenuhi.          \\ \hline
	\texttt{KF-TA-05}  & \texttt{I-TA-01}, \texttt{I-TA-03}, \texttt{I-TA-06} & \texttt{PS-TA-05}                                                                                                                                                                            & Terpenuhi.          \\ \hline
	\texttt{KF-TA-06}  & \texttt{I-TA-02}, \texttt{I-TA-04}, \texttt{I-TA-06} & \texttt{PS-TA-06}, \texttt{PS-TA-09}                                                                                                                                                         & Terpenuhi.          \\ \hline
	\texttt{KF-TA-07}  & \texttt{I-TA-01}, \texttt{I-TA-03}, \texttt{I-TA-06} & \texttt{PS-TA-03}, \texttt{PS-TA-04}                                                                                                                                                         & Terpenuhi.          \\ \hline
	\texttt{KF-TA-08}  & \texttt{I-TA-01}, \texttt{I-TA-03}, \texttt{I-TA-06} & \texttt{PS-TA-07}, \texttt{PS-TA-08}                                                                                                                                                         & Terpenuhi.          \\ \hline
	\texttt{KF-TA-09}  & \texttt{I-TA-01}, \texttt{I-TA-03}, \texttt{I-TA-06} & \texttt{PS-TA-10}, \texttt{PS-TA-11}                                                                                                                                                         & Terpenuhi.          \\ \hline
	\texttt{KNF-TA-01} & \texttt{I-TA-01}, \texttt{I-TA-03}, \texttt{I-TA-06} & \texttt{PU-TA-09}, \texttt{PS-TA-04}, \texttt{PS-TA-08}, \texttt{PS-TA-11}                                                                                                                   & Terpenuhi.          \\ \hline
	\texttt{KNF-TA-02} & \texttt{I-TA-02}, \texttt{I-TA-03}                   & \texttt{PU-TA-01}, \texttt{PU-TA-02}, \texttt{PU-TA-03}, \texttt{PU-TA-04}                                                                                                                   & Terpenuhi.          \\ \hline
	\texttt{KNF-TA-03} & \texttt{I-TA-02}, \texttt{I-TA-03}, \texttt{I-TA-06} & \texttt{PU-TA-02}, \texttt{PU-TA-05}, \texttt{PU-TA-06}, \texttt{PU-TA-07}, \texttt{PU-TA-09}, \texttt{PU-TA-10}, \texttt{PS-TA-04}, \texttt{PS-TA-08}, \texttt{PS-TA-11}--\texttt{PS-TA-15} & Terpenuhi.          \\ \hline
	\texttt{KNF-TA-04} & \texttt{I-TA-04}, \texttt{I-TA-05}, \texttt{I-TA-06} & \texttt{PS-TA-01}, \texttt{PS-TA-02}, \texttt{PS-TA-06}, \texttt{PS-TA-09}, \texttt{PS-TA-15}                                                                                                & Terpenuhi sebagian. \\ \hline
	\texttt{KNF-TA-05} & \texttt{I-TA-03}, \texttt{I-TA-05}                   & \texttt{PU-TA-10}, \texttt{PS-TA-15}                                                                                                                                                         & Terpenuhi.          \\ \hline
	\texttt{KNF-TA-06} & \texttt{I-TA-03}, \texttt{I-TA-06}                   & \texttt{PU-TA-07}, \texttt{PS-TA-13}                                                                                                                                                         & Terpenuhi sebagian. \\ \hline
\end{longtable}
\endgroup





   \chapter{Penutup}
\label{chap:penutup}

\section{Kesimpulan}

Berdasarkan hasil analisis, perancangan, implementasi, dan evaluasi yang telah dilakukan, dapat disimpulkan beberapa hal sebagai berikut:

\begin{enumerate}
	\item Merancang dan mengimplementasikan kontrol akses granular pada DecMed dengan Macaroons sebagai token otorisasi. Batasan akses direpresentasikan melalui \textit{caveats}, sedangkan penegakan otoritas dilakukan pada \textit{Server} PRE sebelum data diakses.

	\item Segmentasi data RME diterapkan dengan memecah data \textit{off-chain} berdasarkan \textit{dataset category} dan \textit{function category}. Dengan struktur ini, data RME tidak lagi diperlakukan sebagai satu \textit{encrypted blob} besar, tetapi sebagai kumpulan segmen klinis. Metadata setiap segmen digunakan untuk menentukan segmen yang dapat diakses oleh aktor sesuai cakupan token.

	\item Atenuasi \textit{offline} diterapkan untuk mendukung pendelegasian akses antar tenaga kesehatan tanpa penerbitan token baru oleh pasien pada setiap perpindahan akses. Pemegang token membentuk token turunan secara \textit{offline} dengan menambahkan \textit{caveats} yang mempersempit cakupan token induk. Implementasi ini juga mencakup pencatatan delegasi, pembatasan kedalaman delegasi, validasi agar token turunan tidak memperluas hak akses token induk, serta revokasi berantai.
\end{enumerate}

\section{Saran}

Berdasarkan batasan dan hasil evaluasi Tugas Akhir ini, beberapa saran untuk pengembangan berikutnya adalah sebagai berikut:

\begin{enumerate}
	\item Cakupan skenario layanan klinis hingga pelayanan antar fasyankes, memungkinkan skenario nyata seperti rujukan antar fasyankes.

	\item Cakupan skenario layanan klinis hingga pelayanan gawat darurat. Mengimplementasikan data IGD dan mekanisme akses ketika pasien tidak sadar, memungkinkan skenario nyata gawat darurat.

	\item Integrasi dengan standar interoperabilitas data kesehatan dapat diteliti lebih lanjut. Segmentasi RME dapat dikembangkan untuk integrasi dengan SATUSEHAT atau pemetaan ke standar FHIR.

	\item Integrasi dan Pengujian dengan Sistem Rumah Sakit untuk mendukung penggunaan sistem dalam konteks nyata. Disesuaikan dengan tata laksana atau dokumen Pedoman Nasional Pelayanan Kesehatan.
\end{enumerate} 

\printbibliography

\titlespacing*{\chapter}{0pt}{0pt}{0pt}
\titlespacing*{\section}{0pt}{0pt}{*1}

\titleformat*{\section}{\bfseries}

\restoredefaultfloatplacements
\appendix

\chapter{Metadata Rekam Medis Elektronik}
\label{lampiran:metadata-rme}

Tabel berikut menjabarkan metadata variabel RME yang diadopsi dari KMK No. HK.01.07/MENKES/1423/2022 untuk dataset Rawat Jalan, Rawat Inap, Laboratorium, dan Apotek \autocite{kemenkes1423-2022}. Klasifikasi dokumentasi pada tabel berikut diperoleh melalui interpretasi terhadap definisi operasional metadata yang tercantum pada dokumen KMK tersebut. Dokumentasi administratif tidak dibagi ke dalam segmentasi yang lebih rinci karena tidak menjadi fokus utama pada rancangan segmentasi data dalam Tugas Akhir ini.

\begin{longtable}{|p{4cm}|p{5.2cm}|p{3.2cm}|}
	\caption{Metadata Variabel RME Rawat Jalan}
	\label{tab:metadata-rme-rawat-jalan}                                                                  \\
	\hline
	\textbf{Kelompok Data}        &
	\textbf{Subkelompok Metadata} &
	\textbf{Klasifikasi Dokumentasi}                                                                      \\ \hline
	\endfirsthead

	\hline
	\textbf{Kelompok Data}        &
	\textbf{Subkelompok Metadata} &
	\textbf{Klasifikasi Dokumentasi}                                                                      \\ \hline
	\endhead

	Lembar Identitas
	                              & Identitas pasien dan data dasar rekam medis
	                              & Administratif                                                         \\ \hline

	Cara Pembayaran
	                              & Metode atau skema pembayaran pelayanan
	                              & Administratif                                                         \\ \hline

	General Consent / Persetujuan Umum
	                              & Persetujuan umum pasien terhadap pelayanan
	                              & Administratif                                                         \\ \hline

	\multirow{3}{4cm}{Formulir Umum / Asesmen Awal Rawat Jalan}
	                              & Anamnesis
	                              & Klinis                                                                \\ \cline{2-3}

	                              & Pemeriksaan Fisik
	                              & Klinis                                                                \\ \cline{2-3}

	                              & Pemeriksaan Psikologis, Sosial Ekonomi, Spiritual
	                              & Klinis                                                                \\ \hline

	\multirow{9}{4cm}{Pemeriksaan Spesialistik}
	                              & Riwayat Penggunaan Obat
	                              & Klinis                                                                \\ \cline{2-3}

	                              & Rencana Rawat
	                              & Klinis                                                                \\ \cline{2-3}

	                              & Instruksi Medik dan Keperawatan
	                              & Klinis                                                                \\ \cline{2-3}

	                              & Pemeriksaan Penunjang
	                              & Klinis                                                                \\ \cline{2-3}





	                              & Diagnosis
	                              & Klinis                                                                \\ \cline{2-3}

	                              & Persetujuan Tindakan / Penolakan Tindakan \textit{(Informed Consent)}
	                              & Klinis                                                                \\ \cline{2-3}

	                              & Terapi: Tindakan dan Obat / Peresepan
	                              & Klinis                                                                \\ \hline
\end{longtable}

\begin{longtable}{|p{4cm}|p{5.2cm}|p{3.2cm}|}
	\caption{Metadata Variabel RME Rawat Inap}
	\label{tab:metadata-rme-rawat-inap}                                                                   \\
	\hline
	\textbf{Kelompok Data}        &
	\textbf{Subkelompok Metadata} &
	\textbf{Klasifikasi Dokumentasi}                                                                      \\ \hline
	\endfirsthead

	\hline
	\textbf{Kelompok Data}        &
	\textbf{Subkelompok Metadata} &
	\textbf{Klasifikasi Dokumentasi}                                                                      \\ \hline
	\endhead

	Lembar Identitas
	                              & Identitas pasien dan data dasar rekam medis
	                              & Administratif                                                         \\ \hline

	Cara Pembayaran
	                              & Metode atau skema pembayaran pelayanan
	                              & Administratif                                                         \\ \hline

	General Consent / Persetujuan Umum
	                              & Persetujuan umum pasien terhadap pelayanan
	                              & Administratif                                                         \\ \hline

	\multirow{3}{4cm}{Formulir Rawat Inap}
	                              & Anamnesis
	                              & Klinis                                                                \\ \cline{2-3}

	                              & Pemeriksaan Fisik
	                              & Klinis                                                                \\ \cline{2-3}

	                              & Pemeriksaan Psikologis, Sosial Ekonomi, Spiritual
	                              & Klinis                                                                \\ \hline

	\multirow{10}{4cm}{Pemeriksaan Spesialistik}
	                              & Riwayat Penggunaan Obat
	                              & Klinis                                                                \\ \cline{2-3}

	                              & Perencanaan Pemulangan Pasien
	                              & Klinis                                                                \\ \cline{2-3}

	                              & Rencana Rawat
	                              & Klinis                                                                \\ \cline{2-3}

	                              & Instruksi Medik dan Keperawatan
	                              & Klinis                                                                \\ \cline{2-3}

	                              & Pemeriksaan Penunjang
	                              & Klinis                                                                \\ \cline{2-3}





	                              & Diagnosis
	                              & Klinis                                                                \\ \cline{2-3}

	                              & Persetujuan Tindakan / Penolakan Tindakan \textit{(Informed Consent)}
	                              & Klinis                                                                \\ \cline{2-3}

	                              & Terapi
	                              & Klinis                                                                \\ \hline
\end{longtable}

\begin{longtable}{|p{4cm}|p{5.2cm}|p{3.2cm}|}
	\caption{Metadata Variabel RME Laboratorium}
	\label{tab:metadata-rme-laboratorium}                             \\
	\hline
	\textbf{Kelompok Data}        &
	\textbf{Subkelompok Metadata} &
	\textbf{Klasifikasi Dokumentasi}                                  \\ \hline
	\endfirsthead

	\hline
	\textbf{Kelompok Data}        &
	\textbf{Subkelompok Metadata} &
	\textbf{Klasifikasi Dokumentasi}                                  \\ \hline
	\endhead

	\multirow{7}{4cm}{Laboratorium}
	                              & Identitas Laboratorium dan Pasien
	                              & Administratif                     \\ \cline{2-3}

	                              & Permintaan Pemeriksaan
	                              & Klinis                            \\ \cline{2-3}

	                              & Spesimen Klinis
	                              & Klinis                            \\ \cline{2-3}

	                              & Fiksasi dan Pengolahan Spesimen
	                              & Klinis                            \\ \cline{2-3}

	                              & Hasil Pemeriksaan
	                              & Klinis                            \\ \cline{2-3}

	                              & Interpretasi dan Validasi Hasil
	                              & Klinis                            \\ \hline
\end{longtable}

\begin{longtable}{|p{4cm}|p{5.2cm}|p{3.2cm}|}
	\caption{Metadata Variabel RME Apotek}
	\label{tab:metadata-rme-apotek}                                        \\
	\hline
	\textbf{Kelompok Data}        &
	\textbf{Subkelompok Metadata} &
	\textbf{Klasifikasi Dokumentasi}                                       \\ \hline
	\endfirsthead

	\hline
	\textbf{Kelompok Data}        &
	\textbf{Subkelompok Metadata} &
	\textbf{Klasifikasi Dokumentasi}                                       \\ \hline
	\endhead

	\multirow{7}{4cm}{Peresepan}
	                              & Identitas Pasien
	                              & Administratif                          \\ \cline{2-3}

	                              & Data Resep dan Obat
	                              & Klinis                                 \\ \cline{2-3}

	                              & Riwayat Penggunaan Obat
	                              & Klinis                                 \\ \cline{2-3}

	                              & Riwayat Alergi
	                              & Klinis                                 \\ \cline{2-3}

	                              & Asal Resep
	                              & Klinis                                 \\ \cline{2-3}

	                              & Dokter Penulis Resep
	                              & Klinis                                 \\ \cline{2-3}

	                              & Status dan Pengkajian Resep
	                              & Klinis                                 \\ \hline

	\multirow{5}{4cm}{Dispensing}
	                              & Status Resep
	                              & Klinis                                 \\ \cline{2-3}

	                              & Waktu Penyiapan Obat
	                              & Klinis                                 \\ \cline{2-3}

	                              & Waktu Penyerahan Obat
	                              & Klinis                                 \\ \cline{2-3}

	                              & Nama Petugas yang Melakukan Dispensing
	                              & Klinis                                 \\ \cline{2-3}

	                              & Etiket
	                              & Klinis                                 \\ \hline
\end{longtable} \chapter{Alur Pemberian Akses Medis pada DecMed}
\label{lampiran:alur-pemberian-akses-medis}

Alur pemberian akses medis dari pasien kepada tenaga medis pada implementasi DecMed saat ini dapat dilihat pada Gambar \ref{fig:patient-create-access-medical-personnel}.

\begin{figure}[ht]
	\centering
	\includegraphics[width=\textwidth]{resources/chapter-3/analisis/patient_create_access_for_medical_personnel_svg.png}
	\caption{\textit{Sequence Diagram} Alur Pemberian Akses Medis dari Pasien kepada Tenaga Medis \autocite{sadewa2025decmed}}
	\label{fig:patient-create-access-medical-personnel}
\end{figure}
 \chapter{Kebutuhan Fungsional DecMed}
\label{lampiran:kebutuhan-fungsional-decmed}

\begin{longtable}{|p{1.2cm}|p{11cm}|}
	\caption{Kebutuhan Fungsional DecMed \autocite{sadewa2025decmed}}
	\label{tab:kebutuhan-fungsional-decmed}                                                                                                                      \\
	\hline
	\textbf{ID} & \textbf{Deskripsi}                                                                                                                             \\
	\hline
	\endfirsthead

	\multicolumn{2}{c}{\tablename\ \thetable{} -- Lanjutan}                                                                                                      \\
	\hline
	\textbf{ID} & \textbf{Deskripsi}                                                                                                                             \\
	\hline
	\endhead

	\hline
	\endfoot

	F01         & Pasien dapat melakukan registrasi akun pada sistem                                                                                             \\ \hline
	F02         & Pasien dapat masuk ke akun menggunakan credentials yang dimiliki                                                                               \\ \hline
	F03         & Pasien dapat keluar dari akun                                                                                                                  \\ \hline
	F04         & Pasien dapat melakukan scan kode QR yang dimiliki oleh masing-masing tenaga kesehatan dan petugas rekam medis                                  \\ \hline
	F05         & Pasien dapat membaca seluruh RME yang ia miliki                                                                                                \\ \hline
	F06         & Pasien dapat memberikan akses baca bagian administratif dari data RME miliknya kepada petugas rekam medis untuk rentang waktu tertentu         \\ \hline
	F07         & Pasien dapat memberikan akses baca bagian administratif dan klinis dari data RME miliknya kepada tenaga kesehatan untuk rentang waktu tertentu \\ \hline
	F08         & Pasien dapat melihat daftar log dari akses terhadap data RME yang dimiliki                                                                     \\ \hline
	F09         & Pasien dapat mencabut akses yang telah diberikan. Akses tersebut berupa akses baca maupun tulis terhadap data RME.                             \\ \hline
	F10         & Tenaga kesehatan dapat melakukan aktivasi client                                                                                               \\ \hline
	F11         & Tenaga kesehatan dapat melakukan registrasi akun pada sistem                                                                                   \\ \hline
	F12         & Tenaga kesehatan dapat masuk ke akun menggunakan credentials yang dimiliki                                                                     \\ \hline
	F13         & Tenaga kesehatan dapat keluar dari akun                                                                                                        \\ \hline
	F14         & Tenaga kesehatan dapat membaca bagian administratif dan klinis dari data RME milik pasien setelah diberikan akses dalam waktu tertentu         \\ \hline
	F15         & Tenaga kesehatan dapat menambahkan entri RME baru                                                                                              \\ \hline
	F16         & Tenaga kesehatan dapat mengedit bagian klinis dari entri RME baru yang dibuatnya pada rentang waktu yang telah ditentukan                      \\ \hline
	F17         & Petugas rekam medis dapat melakukan aktivasi client                                                                                            \\ \hline
	F18         & Petugas rekam medis dapat melakukan registrasi akun pada sistem                                                                                \\ \hline
	F19         & Petugas rekam medis dapat masuk ke akun menggunakan credentials yang dimiliki                                                                  \\ \hline
	F20         & Petugas rekam medis dapat keluar dari akun                                                                                                     \\ \hline
	F21         & Petugas rekam medis dapat membaca bagian administratif dari data RME milik pasien setelah diberikan akses dalam waktu tertentu                 \\ \hline
	F22         & Admin fasyankes dapat melakukan aktivasi client                                                                                                \\ \hline
	F23         & Admin fasyankes dapat melakukan registrasi akun pada sistem                                                                                    \\ \hline
	F24         & Admin fasyankes dapat masuk ke akun menggunakan credentials yang dimiliki                                                                      \\ \hline
	F25         & Admin fasyankes dapat keluar dari akun                                                                                                         \\ \hline
	F26         & Admin fasyankes dapat menambahkan personel fasyankes baru                                                                                      \\ \hline
	F27         & Admin fasyankes dapat membuat kunci aktivasi baru untuk personel fasyankes                                                                     \\ \hline
	F28         & Kementerian kesehatan dapat mendaftarkan fasyankes dan admin untuk fasyankes tersebut                                                          \\ \hline
	F29         & Kementerian kesehatan dapat membuat kunci aktivasi baru untuk admin fasyankes                                                                  \\ \hline
\end{longtable}  \chapter{DFD DecMed}
\label{lampiran:dfd-decmed}

Hasil pemodelan DFD DecMed Level 0 dan Level 1 dapat dilihat pada Gambar \ref{fig:dfd-decmed-level-0} dan \ref{fig:dfd-decmed-level-1}, sedangkan legenda pemodelan kedua DFD tersebut dapat dilihat pada Gambar \ref{fig:legenda-dfd}.

\begin{figure}[ht]
	\centering
	\includegraphics[width=1\textwidth]{resources/appendix/legenda-dfd.png}
	\caption{Legenda \textit{Data Flow Diagram} (Dokumentasi Penulis)}
	\label{fig:legenda-dfd}
\end{figure}

\begin{figure}[ht]
	\centering
	\includegraphics[width=1\textwidth]{resources/chapter-3/analisis/DFD-DecMed-Level-0.png}
	\caption{DFD DecMed Level 0 (Dokumentasi Penulis)}
	\label{fig:dfd-decmed-level-0}
\end{figure}

\begin{landscape}
	\begin{figure}[ht]
		\centering
		\includegraphics[height=0.80\textheight]{resources/chapter-3/analisis/DFD-DecMed-Level-1.png}
		\caption{DFD DecMed Level 1 (Dokumentasi Penulis)}
		\label{fig:dfd-decmed-level-1}
	\end{figure}
\end{landscape}

\begin{table}[!ht]
	\centering
	\caption{Daftar \textit{External Entities} pada DFD DecMed}
	\label{tab:dfd-external-entities}
	\small
	\begin{tabular}{|p{0.10\textwidth}|p{0.22\textwidth}|p{0.58\textwidth}|}
		\hline
		\textbf{ID} & \textbf{Nama}          & \textbf{Interaksi utama}                                                  \\ \hline
		E1          & Patient                & PIN, QR delegasi akses, \textit{revoke access}, lihat \textit{access log} \\ \hline
		E2          & Hospital Personnel     & PIN, baca/\textit{update} rekam medis, data administratif                 \\ \hline
		E3          & Ministry Administrator & Registrasi rumah sakit, \textit{activation key}                           \\ \hline
		E4          & IOTA Network           & RPC, \textit{publish/execute transaction}, \textit{read state}            \\ \hline
		E5          & IOTA Gas Station       & \texttt{reserve\_gas}, \texttt{execute\_tx}                               \\ \hline
	\end{tabular}
\end{table}

\begin{table}[!ht]
	\centering
	\caption{Daftar \textit{Process} DFD DecMed}
	\label{tab:dfd-process}
	\small
	\begin{tabular}{|p{0.06\textwidth}|p{0.5\textwidth}|}
		\hline
		\textbf{ID} & \textbf{Nama}                    \\ \hline
		P0          & Decmed (seluruh sistem)          \\ \hline
		P1          & Hospital Client                  \\ \hline
		P2          & Patient Client                   \\ \hline
		P3          & Ministry Client                  \\ \hline
		P4          & Proxy Re-encryption API          \\ \hline
		P5          & Authentication \& Access Control \\ \hline
		P6          & Medical Record Management        \\ \hline
		P7          & Proxy Re-encryption Engine       \\ \hline
		P8          & IOTA Transaction Client          \\ \hline
	\end{tabular}
\end{table}

\begin{table}[!ht]
	\centering
	\caption{Daftar \textit{Data Store} pada DFD DecMed}
	\label{tab:dfd-data-store}
	\small
	\begin{tabular}{|p{0.08\textwidth}|p{0.28\textwidth}|p{0.52\textwidth}|}
		\hline
		\textbf{ID} & \textbf{Nama}        & \textbf{Data sensitif}                                                                    \\ \hline
		D1          & Redis Cache          & \textit{Nonce challenge}, \texttt{AccessKeys} (k-frag, PRE public keys, seed terenkripsi) \\ \hline
		D2          & IPFS Cluster Storage & Blob rekam medis terenkripsi (CID direferensikan \textit{on-chain})                       \\ \hline
		D3          & Client OS Keyring    & IOTA keypair, PRE keys, PIN-derived secrets, \textit{activation key}                      \\ \hline
		D4          & IOTA On-Chain State  & Metadata, \textit{access grants}, \textit{access log}, \textit{hospital registry}         \\ \hline
	\end{tabular}
\end{table}

\begin{table}[!ht]
	\centering
	\caption{Daftar Data Flow pada DFD DecMed (Level 0)}
	\label{tab:dfd-dataflow-level-0}
	\small
	\begin{tabular}{|p{0.06\textwidth}|p{0.39\textwidth}|p{0.24\textwidth}|}
		\hline
		\textbf{ID} & \textbf{Label}                   & \textbf{Arah}       \\ \hline
		F1          & Patient Identity \& Keys         & E1~$\rightarrow$~P0 \\ \hline
		F2          & Access Delegation Grant          & E1~$\rightarrow$~P0 \\ \hline
		F3          & Hospital Personnel Credentials   & E2~$\rightarrow$~P0 \\ \hline
		F4          & Encrypted Medical Record Request & E2~$\rightarrow$~P0 \\ \hline
		F5          & Hospital Registration Request    & E3~$\rightarrow$~P0 \\ \hline
		F6          & On-Chain Transaction             & P0~$\rightarrow$~E4 \\ \hline
		F7          & On-Chain State Query             & P0~$\rightarrow$~E4 \\ \hline
		F10         & Re-encrypted Record Response     & P0~$\rightarrow$~E2 \\ \hline
		F11         & Decrypted Record View            & P0~$\rightarrow$~E1 \\ \hline
		F12         & Activation Key Delivery          & P0~$\rightarrow$~E2 \\ \hline
		F18         & Encrypted Blob                   & P0~$\rightarrow$~D2 \\ \hline
	\end{tabular}
\end{table}

\begin{table}[!ht]
	\centering
	\caption{Daftar Data Flow pada DFD DecMed (Level 1)}
	\label{tab:dfd-dataflow-level-1}
	\small
	\begin{tabular}{|p{0.08\textwidth}|p{0.43\textwidth}|p{0.31\textwidth}|}
		\hline
		\textbf{ID} & \textbf{Label}               & \textbf{Arah Ringkas} \\ \hline
		F1          & Patient PIN \& Session       & E1~$\rightarrow$~P2   \\ \hline
		F2          & QR Access Delegation Payload & E1~$\rightarrow$~P2   \\ \hline
		F3          & Hospital Personnel PIN       & E2~$\rightarrow$~P1   \\ \hline
		F4          & Ministry Admin Credentials   & E3~$\rightarrow$~P3   \\ \hline
		F5          & IOTA Key Material            & D3~$\rightarrow$~P1   \\ \hline
		F6          & IOTA Key Material            & D3~$\rightarrow$~P2   \\ \hline
		F7          & Proxy API Request (JWT)      & P1~$\rightarrow$~P4   \\ \hline
		F8          & Access Delegation API Call   & P2~$\rightarrow$~P4   \\ \hline
		F9          & On-Chain Access Grant Tx     & P2~$\rightarrow$~D4   \\ \hline
		F10         & Hospital Registration Tx     & P3~$\rightarrow$~D4   \\ \hline
		F11         & API Route Dispatch           & P4~$\rightarrow$~P5   \\ \hline
		F12         & API Route Dispatch           & P4~$\rightarrow$~P6   \\ \hline
		F13         & Authentication Nonce         & P5~$\rightarrow$~D1   \\ \hline
		F14         & Proxy Access Keys            & P5~$\rightarrow$~D1   \\ \hline
		F15         & JWT Access Token             & P5~$\rightarrow$~P1   \\ \hline
		F16         & Re-encryption Job            & P6~$\rightarrow$~P7   \\ \hline
		F17         & Proxy Access Keys            & D1~$\rightarrow$~P7   \\ \hline
		F18         & Encrypted Medical Blob       & P6~$\rightarrow$~D2   \\ \hline
		F19         & Encrypted Medical Blob       & D2~$\rightarrow$~P7   \\ \hline
		F20         & Ciphertext Fragment          & P7~$\rightarrow$~P1   \\ \hline
		F21         & On-Chain Metadata            & P8~$\rightarrow$~D4   \\ \hline
		F22         & On-Chain Metadata            & D4~$\rightarrow$~P6   \\ \hline
		F23         & Move Transaction             & P8~$\rightarrow$~E4   \\ \hline
		F24         & Gas Sponsorship Request      & P8~$\rightarrow$~E5   \\ \hline
		F25         & Gas Coin Reservation         & E5~$\rightarrow$~P8   \\ \hline
		F26         & Personnel Role Lookup        & P5~$\rightarrow$~P8   \\ \hline
		F27         & Medical Record Metadata      & P6~$\rightarrow$~P8   \\ \hline
		F28         & Client-Side Decryption       & P1~$\rightarrow$~E2   \\ \hline
		F29         & Access Log View              & D4~$\rightarrow$~P2   \\ \hline
	\end{tabular}
\end{table}
 \chapter{Analisis STRIDE pada Sistem DecMed}
\label{lampiran:analisis-stride-decmed}

{\small
\renewcommand{\arraystretch}{1.3}
\setlength{\tabcolsep}{3pt}
\begin{longtable}{|p{0.65cm}|p{2.2cm}|p{1.9cm}|p{4cm}|p{4cm}|}
	\caption{Analisis STRIDE pada Sistem DecMed}
	\label{tab:stride-decmed}                                                                                                                                                                                                                                                \\
	\hline
	\textbf{ID}                                                                                                                                                                 & \textbf{Komponen / Flow} & \textbf{Kategori} & \textbf{Skenario Ancaman} & \textbf{Dampak} \\
	\hline
	\endfirsthead

	\caption[]{Analisis STRIDE pada Sistem DecMed (lanjutan)}                                                                                                                                                                                                                \\
	\hline
	\textbf{ID}                                                                                                                                                                 & \textbf{Komponen / Flow} & \textbf{Kategori} & \textbf{Skenario Ancaman} & \textbf{Dampak} \\
	\hline
	\endhead

	T01                                                                                                                                                                         &
	P5, F7                                                                                                                                                                      &
	\textit{Spoofing}                                                                                                                                                           &
	Penyerang menggunakan token akses yang dicuri atau dipalsukan untuk data RME.                                                                                               &
	Akses tidak sah terhadap data rekam medis pasien.                                                                                                                                                                                                                        \\
	\hline

	T02                                                                                                                                                                         &
	P8, D4, E4, dan F9                                                                                                                                                          &
	\textit{Tampering}                                                                                                                                                          &
	Pengguna mengirim transaksi pemberian akses yang valid secara kriptografis, tetapi cakupan akses tidak dibatasi secara tepat oleh logika otorisasi \textit{smart contract}. &
	Metadata akses atau metadata RME dapat menunjuk pada cakupan data yang tidak semestinya, sehingga integritas terganggu.                                                                                                                                                  \\
	\hline

	T03                                                                                                                                                                         &
	P6, D2, dan F19                                                                                                                                                             &
	\textit{Tampering}                                                                                                                                                          &
	Referensi CID yang mengarah ke objek \textit{off-chain} dimanipulasi sehingga sistem mengambil \textit{blob} terenkripsi yang tidak sesuai.                                 &
	Data RME yang diambil tidak valid, tidak konsisten, atau gagal diverifikasi.                                                                                                                                                                                             \\
	\hline

	T04                                                                                                                                                                         &
	P2, P5, P8, D4, F8, F9, F21, dan F29                                                                                                                                        &
	\textit{Repudiation}                                                                                                                                                        &
	Pihak tertentu menyangkal pernah memberikan atau menggunakan akses.                                                                                                         &
	Kesulitan dalam audit dan investigasi penyalahgunaan akses.                                                                                                                                                                                                              \\
	\hline

	T05                                                                                                                                                                         &
	P5, P8, D4, F16, F19, F20, F21, dan F29                                                                                                                                     &
	\textit{Information Disclosure}                                                                                                                                             &
	Delegasi akses membuka cakupan data pasien yang lebih luas dari kebutuhan.                                                                                                  &
	Kebocoran data RME pasien.                                                                                                                                                                                                                                               \\
	\hline

	T06                                                                                                                                                                         &
	P8, D4, dan F22                                                                                                                                                             &
	\textit{Information Disclosure}                                                                                                                                             &
	Metadata transaksi mengungkap relasi pasien, fasyankes, atau data pasien.                                                                                                   &
	Privasi pasien dan informasi layanan klinis jadi terekspos.                                                                                                                                                                                                              \\
	\hline

	T07                                                                                                                                                                         &
	P4, P5, P6, P7, P8, D1, E4, dan F7                                                                                                                                          &
	\textit{Denial of Service}                                                                                                                                                  &
	Layanan PRE atau penyimpanan IPFS tidak dapat diakses akibat serangan atau gangguan layanan.                                                                                &
	Proses akses rekam medis menjadi lambat atau tidak tersedia.                                                                                                                                                                                                             \\
	\hline

	T08                                                                                                                                                                         &
	P2, P5, P8, dan D4, F8, F9,                                                                                                                                                 &
	\textit{Elevation of Privilege}                                                                                                                                             &
	Pengguna memperoleh hak akses lebih tinggi melalui delegasi yang tidak dibatasi.                                                                                            &
	Terjadi eskalasi hak akses di luar kebutuhan klinis yang sah.                                                                                                                                                                                                            \\
	\hline
\end{longtable}
}




















%
  \chapter{Legenda Diagram Arsitektur Solusi}
\label{lampiran:legenda-diagram-arsitektur-solusi}

\begin{figure}[ht]
	\centering
	\includegraphics[width=1\textwidth]{resources/appendix/decmed_architecture_Legenda.png}
	\caption{Legenda Diagram Arsitektur Solusi (Dokumentasi Penulis)}
	\label{fig:legenda-diagram-arsitektur-solusi}
\end{figure} \chapter{Skema Penyimpanan Data RME}
\label{lampiran:skema-penyimpanan-rme}

Lampiran ini memuat rincian skema penyimpanan data RME yang digunakan pada rancangan segmentasi data. Skema penyimpanan dibagi menjadi metadata segmen yang disimpan pada penyimpanan \textit{on-chain} dan isi segmen terenkripsi yang disimpan pada penyimpanan \textit{off-chain}.

\begin{longtable}{|>{\raggedright\arraybackslash}p{2.8cm}|>{\raggedright\arraybackslash}p{4.0cm}|>{\raggedright\arraybackslash}p{5.4cm}|}
	\caption{Skema Metadata Segmen RME \textit{On-Chain}} \label{tab:onchain-schema}                                                              \\
	\hline
	\textbf{Komponen} & \textbf{Key JSON}                                                                                   & \textbf{Keterangan} \\
	\hline
	\endfirsthead
	\hline
	\textbf{Komponen} & \textbf{Key JSON}                                                                                   & \textbf{Keterangan} \\
	\hline
	\endhead

	\multirow{4}{2.8cm}{\textbf{Identitas Segmen}}
	                  & \texttt{segment\_id}
	                  & UUID unik untuk satu segmen data RME.                                                                                     \\ \cline{2-3}
	                  & \texttt{related\_rme\_id}
	                  & ID RME atau episode pelayanan yang menjadi induk segmen, dengan format \verb|RME-{6_digit}|.                              \\ \cline{2-3}
	                  & \texttt{patient\_address}
	                  & Alamat \textit{wallet} pasien pemilik RME.                                                                                \\ \cline{2-3}
	                  & \texttt{hospital\_id}
	                  & ID fasilitas pelayanan kesehatan ketika menerima pelayanan.                                                               \\ \hline

	\multirow{2}{2.8cm}{\textbf{Kategori Data}}
	                  & \texttt{dataset\_category}
	                  & Kategori dataset. Sesuai dengan tabel \ref{tab:mapping-dataset-category}.                                                 \\ \cline{2-3}
	                  & \texttt{function\_category}
	                  & Kategori fungsi klinis atau administratif. Sesuai dengan tabel \ref{tab:mapping-function-category}.                       \\ \hline

	\multirow{2}{2.8cm}{\textbf{Pointer dan Integritas}}
	                  & \texttt{ipfs\_cid}
	                  & CID IPFS dari segmen data terenkripsi.                                                                                    \\ \cline{2-3}
	                  & \texttt{integrity\_hash}
	                  & Hash SHA-256 dari segmen terenkripsi.                                                                                     \\ \hline

	\multirow{2}{2.8cm}{\textbf{Kapsul Kunci}}
	                  & \texttt{capsule}
	                  & Kapsul PRE yang diperlukan untuk proses \textit{re-encryption}.                                                           \\ \cline{2-3}
	                  & \texttt{enc\_key\_and\_nonce}
	                  & Kunci simetris dan \textit{nonce} yang sudah terenkripsi untuk membuka segmen data.                                       \\ \hline

	\multirow{2}{2.8cm}{\textbf{Audit Metadata}}
	                  & \texttt{created\_at}
	                  & Waktu pembuatan metadata segmen.                                                                                          \\ \cline{2-3}
	                  & \texttt{author\_address}
	                  & Alamat \textit{wallet} aktor yang membuat atau memperbarui segmen.                                                        \\ \hline
\end{longtable}

\begin{longtable}{|>{\raggedright\arraybackslash}p{2.8cm}|>{\raggedright\arraybackslash}p{4.0cm}|>{\raggedright\arraybackslash}p{5.4cm}|}
	\caption{Skema Segmen Data RME \textit{Off-Chain}} \label{tab:offchain-schema}                                                                                     \\
	\hline
	\textbf{Komponen} & \textbf{Key JSON}                                                                                                        & \textbf{Keterangan} \\
	\hline
	\endfirsthead
	\hline
	\textbf{Komponen} & \textbf{Key JSON}                                                                                                        & \textbf{Keterangan} \\
	\hline
	\endhead

	\multirow{5}{2.8cm}{\textbf{Header Segmen}}
	                  & \texttt{segment\_id}
	                  & UUID yang sama dengan \texttt{segment\_id} pada metadata \textit{on-chain}.                                                                    \\ \cline{2-3}
	                  & \texttt{related\_rme\_id}
	                  & ID RME atau episode pelayanan yang menjadi induk segmen, dengan format \verb|RME-{6_digit}|.                                                   \\ \cline{2-3}
	                  & \texttt{dataset\_category}
	                  & Kategori dataset sesuai Tabel \ref{tab:mapping-dataset-category}.                                                                              \\ \cline{2-3}
	                  & \texttt{function\_category}
	                  & Kategori fungsi sesuai dengan Tabel \ref{tab:mapping-function-category}.                                                                       \\ \hline

	\multirow{4}{2.8cm}{\textbf{Konteks Medis}}
	                  & \texttt{patient\_address}
	                  & Alamat IOTA pasien pemilik RME.                                                                                                                \\ \cline{2-3}
	                  & \texttt{service\_date}
	                  & Tanggal dan waktu ketika segmen tercatat dengan format \texttt{YYYY-MM-DD HH:MM:SS} beserta informasi \textit{timezone}.                       \\ \cline{2-3}
	                  & \texttt{author\_address}
	                  & Alamat IOTA aktor yang menulis isi segmen.                                                                                                     \\ \hline

	\multirow{2}{2.8cm}{\textbf{Isi Segmen}}
	                  & \texttt{payload}
	                  & Objek segmen data RME dengan tipe data \texttt{Value}. Pada implementasi ini, isi segmen dalam bentuk \texttt{text}.                           \\ \cline{2-3}
	                  & \texttt{payload\_hash}
	                  & \textit{Hash} isi \texttt{payload}.                                                                                                            \\ \hline
\end{longtable}
 \chapter{Contoh \textit{Caveat} Token Macaroon untuk Setiap Tahapan}
\label{lampiran:contoh-caveat-macaroon}

Lampiran ini menyajikan contoh \textit{caveat} token Macaroon pada setiap tahapan alur pemberian akses, mulai dari penerbitan token awal oleh pasien hingga rangkaian delegasi yang dilakukan oleh tenaga medis. Nilai alamat \textit{wallet} (\texttt{0xPASIEN}, \texttt{0xPETUGAS\_RM}, \texttt{0xDOKTER}, \texttt{0xPERAWAT}, \texttt{0xLAB}, \texttt{0xAPOTEKER}), \texttt{hospital-001}, dan \texttt{RME-001} hanya bersifat ilustratif. Pada implementasi, nilai-nilai tersebut akan berbeda tergantung pada alamat \textit{wallet} pemegang token serta ID RME yang di-\textit{generate}. Simbol titik tiga (\texttt{...}) adalah bentuk ilustrasi yang digunakan untuk menunjukkan bahwa token dapat memiliki lebih banyak nilai lainnya selain yang dituliskan.

\section*{Penerbitan Awal: Pasien $\rightarrow$ Petugas Rekam Medis}
\label{subsec:contoh-caveat-token-macaroon-penerbitan-awal}

\begin{lstlisting}[caption={Contoh \textit{Caveat} Token \texttt{Read} Awal untuk Petugas Rekam Medis}, label={lst:contoh-token-baca-awal-rekam-medis}]
purpose = Read
patient_address = 0xPASIEN
root_subject = 0xPETUGAS_RM
hospital_cid = hospital-001

read_dataset_in = [RAWAT_JALAN, RAWAT_INAP, LABORATORIUM, APOTEK]

read_function_in = [
  ADMINISTRATIVE_GENERAL,
  ANAMNESIS,
  PEMERIKSAAN_FISIK,
  PEMERIKSAAN_PSIKOLOGIS,
  ...
]  // seluruh 15 fungsi

expires_before = 2026-05-16T18:00:00
max_delegation_depth = 2
\end{lstlisting}

\begin{lstlisting}[caption={Contoh \textit{Caveat} Token \texttt{Update} Awal untuk Petugas Rekam Medis}, label={lst:contoh-token-pembaruan-awal-rekam-medis}]
purpose = Update
patient_address = 0xPASIEN
root_subject = 0xPETUGAS_RM
related_rme_id = RME-001
hospital_cid = hospital-001

write_dataset_in = [RAWAT_JALAN, LABORATORIUM, APOTEK]

write_function_in = [
  ... // union fungsi dari ketiga dataset di atas
]

expires_before = 2026-05-16T18:00:00
max_delegation_depth = 2
\end{lstlisting}

\section*{Delegasi: Petugas Rekam Medis $\rightarrow$ Dokter}
\label{subsec:contoh-caveat-token-macaroon-delegasi-dokter}

\begin{lstlisting}[caption={Contoh \textit{Caveat} Delegasi Petugas Rekam Medis $\rightarrow$ Dokter}, label={lst:contoh-token-delegasi-dokter}]
parent_token_hash = <sha256(parent)>
delegated_by = 0xPETUGAS_RM
delegated_to = 0xDOKTER

read_dataset_in = [RAWAT_JALAN, RAWAT_INAP, LABORATORIUM, APOTEK]
write_dataset_in = [RAWAT_JALAN, LABORATORIUM, APOTEK]

read_function_in = [
  ... // semua 15 fungsi
]
write_function_in = [
  ... // 14 fungsi, ADMINISTRATIVE_GENERAL dihapus
]

related_rme_id = RME-001
expires_before = 2026-05-16T18:00:00
max_delegation_depth = 1
\end{lstlisting}

\section*{Delegasi: Petugas Rekam Medis $\rightarrow$ Perawat}
\label{subsec:contoh-caveat-token-macaroon-delegasi-perawat}

\begin{lstlisting}[caption={Contoh \textit{Caveat} Delegasi Petugas Rekam Medis $\rightarrow$ Perawat}, label={lst:contoh-token-delegasi-perawat}]
parent_token_hash = <sha256(parent)>
delegated_by = 0xPETUGAS_RM
delegated_to = 0xPERAWAT

read_dataset_in = [RAWAT_JALAN]
write_dataset_in = [RAWAT_JALAN]

read_function_in = [
  ADMINISTRATIVE_GENERAL, 
  ANAMNESIS, 
  PEMERIKSAAN_FISIK]
write_function_in = [ANAMNESIS, PEMERIKSAAN_FISIK]

expires_before = 2026-05-16T18:00:00
max_delegation_depth = 1
\end{lstlisting}

\section*{Delegasi: Dokter $\rightarrow$ Petugas Laboratorium}
\label{subsec:contoh-caveat-token-macaroon-delegasi-lab}

\begin{lstlisting}[caption={Contoh \textit{Caveat} Delegasi Dokter $\rightarrow$ Petugas Laboratorium}, label={lst:contoh-token-delegasi-lab}]
parent_token_hash = <sha256(token_dokter)>
delegated_by = 0xDOKTER
delegated_to = 0xLAB

read_dataset_in = [LABORATORIUM]
write_dataset_in = [LABORATORIUM]

read_function_in = [
  ADMINISTRATIVE_GENERAL, 
  PEMERIKSAAN_PENUNJANG, 
  LABORATORIUM]
write_function_in = [LABORATORIUM]

expires_before = 2026-05-16T18:00:00
max_delegation_depth = 0
\end{lstlisting}

\section*{Delegasi: Dokter $\rightarrow$ Apoteker}
\label{subsec:contoh-caveat-token-macaroon-delegasi-apoteker}

\begin{lstlisting}[caption={Contoh \textit{Caveat} Delegasi Dokter $\rightarrow$ Apoteker}, label={lst:contoh-token-delegasi-apoteker}]
parent_token_hash = <sha256(token_dokter)>
delegated_by = 0xDOKTER
delegated_to = 0xAPOTEKER

read_dataset_in = [RAWAT_JALAN, APOTEK]
write_dataset_in = [APOTEK]

read_function_in = [
  ADMINISTRATIVE_GENERAL,
  RIWAYAT_PENGGUNAAN_OBAT,
  TERAPI,
  PERESEPAN,
  DISPENSING
]
write_function_in = [PERESEPAN, DISPENSING]

expires_before = 2026-05-16T18:00:00
max_delegation_depth = 0
\end{lstlisting}
 \chapter{Daftar Dependensi DecMed}
\label{lampiran:daftar-dependensi-perangkat-lunak}

Daftar ini hanya mencakup dependensi langsung yang dideklarasikan pada masing-masing komponen, bukan seluruh dependensi transitif. Dependensi internal yang ditandai dengan \texttt{path} merupakan \textit{crate} lokal yang dikembangkan sebagai bagian dari sistem.

\begin{longtable}{|>{\raggedright\arraybackslash}p{0.55\textwidth}|>{\raggedright\arraybackslash}p{0.38\textwidth}|}
	\caption{Dependensi Umum yang Digunakan pada Beberapa Komponen}
	\label{tab:dependensi-umum}                                       \\
	\hline
	\textbf{\textit{Crate}}               & \textbf{Versi/Sumber}     \\
	\hline
	\endfirsthead
	\hline
	\textbf{\textit{Crate}}               & \textbf{Versi/Sumber}     \\
	\hline
	\endhead

	\texttt{serde}                        & \texttt{1}                \\ \hline
	\texttt{serde\_json}                  & \texttt{1}                \\ \hline
	\texttt{sha2}                         & \texttt{0.10}             \\ \hline
	\texttt{uuid}                         & \texttt{1}                \\ \hline
	\texttt{bip39}                        & \texttt{2}                \\ \hline
	\texttt{iota-sdk}                     & Git rev. \texttt{3d7c588} \\ \hline
	\texttt{iota-keys}                    & Git rev. \texttt{3d7c588} \\ \hline
	\texttt{iota-config}                  & Git rev. \texttt{3d7c588} \\ \hline
	\texttt{iota-types}                   & Git rev. \texttt{3d7c588} \\ \hline
	\texttt{iota-json-rpc-types}          & Git rev. \texttt{3d7c588} \\ \hline
	\texttt{shared\_crypto}               & Git rev. \texttt{3d7c588} \\ \hline
	\texttt{move-core-types}              & Git rev. \texttt{3d7c588} \\ \hline
	\texttt{dotenv}                       & \texttt{0.15}             \\ \hline
	\texttt{bcs}                          & \texttt{0.1.4}            \\ \hline
	\texttt{schemars}                     & \texttt{0.8}              \\ \hline
	\texttt{rand}                         & \texttt{0.9}              \\ \hline
	\texttt{bincode}                      & \texttt{2}                \\ \hline
	\texttt{umbral-pre}                   & \texttt{0.11}             \\ \hline
	\texttt{aes-gcm}                      & \texttt{0.10}             \\ \hline
	\texttt{argon2}                       & \texttt{0.5.3}            \\ \hline
	\texttt{hex}                          & \texttt{0.4.3}            \\ \hline
	\texttt{regex}                        & \texttt{1.11.1}           \\ \hline
	\texttt{generic-array}                & \texttt{1.2.0}            \\ \hline
	\texttt{k256}                         & \texttt{0.13.4}           \\ \hline
	\texttt{elliptic-curve}               & \texttt{0.13.8}           \\ \hline
	\texttt{chrono}                       & \texttt{0.4}              \\ \hline
	\texttt{base64}                       & \texttt{0.22}             \\ \hline
	\texttt{strum}/\texttt{strum\_macros} & \texttt{0.27}             \\ \hline
	\texttt{anyhow}                       & \texttt{1.0}              \\ \hline
	\texttt{thiserror}                    & \texttt{2.0}              \\ \hline
\end{longtable}

\begin{longtable}{|>{\raggedright\arraybackslash}p{0.55\textwidth}|>{\raggedright\arraybackslash}p{0.38\textwidth}|}
	\caption{Dependensi PRE}
	\label{tab:dependensi-backend-pre}                    \\
	\hline
	\textbf{\textit{Crate}}       & \textbf{Versi/Sumber} \\
	\hline
	\endfirsthead
	\hline
	\textbf{\textit{Crate}}       & \textbf{Versi/Sumber} \\
	\hline
	\endhead

	\texttt{axum}                 & \texttt{0.8.4}        \\ \hline
	\texttt{axum-extra}           & \texttt{0.10.1}       \\ \hline
	\texttt{tower}                & \texttt{0.5.2}        \\ \hline
	\texttt{tokio}                & \texttt{1.44.2}       \\ \hline
	\texttt{reqwest}              & \texttt{0.12}         \\ \hline
	\texttt{redis}                & \texttt{0.32}         \\ \hline
	\texttt{r2d2}                 & \texttt{0.8}          \\ \hline
	\texttt{dotenvy}              & \texttt{0.15.7}       \\ \hline
	\texttt{decmed-rme-segment}   & \texttt{path}         \\ \hline
	\texttt{decmed-macaroon-auth} & \texttt{path}         \\ \hline
\end{longtable}

\begin{longtable}{|>{\raggedright\arraybackslash}p{0.55\textwidth}|>{\raggedright\arraybackslash}p{0.38\textwidth}|}
	\caption{Dependensi \textit{Client}}
	\label{tab:dependensi-client-tauri}                       \\
	\hline
	\textbf{\textit{Crate}}           & \textbf{Versi/Sumber} \\
	\hline
	\endfirsthead
	\hline
	\textbf{\textit{Crate}}           & \textbf{Versi/Sumber} \\
	\hline
	\endhead

	\texttt{tauri}                    & \texttt{2}            \\ \hline
	\texttt{tauri-plugin-opener}      & \texttt{2}            \\ \hline
	\texttt{tauri-plugin-http}        & \texttt{2}            \\ \hline
	\texttt{tauri-plugin-fs}          & \texttt{2}            \\ \hline
	\texttt{keyring}                  & \texttt{3}            \\ \hline
	\texttt{rustls-platform-verifier} & \texttt{0.5}          \\ \hline
	\texttt{rqrr}                     & \texttt{0.9}          \\ \hline
	\texttt{image}                    & \texttt{0.25}         \\ \hline
	\texttt{tokio-stream}             & \texttt{0.1.17}       \\ \hline
	\texttt{decmed-rme-segment}       & \texttt{path}         \\ \hline
	\texttt{decmed-macaroon-auth}     & \texttt{path}         \\ \hline
	\texttt{tauri-build}              & \texttt{2}            \\ \hline
\end{longtable}

\begin{longtable}{|>{\raggedright\arraybackslash}p{0.55\textwidth}|>{\raggedright\arraybackslash}p{0.38\textwidth}|}
	\caption{Dependensi \texttt{decmed-rme-segment}}
	\label{tab:dependensi-decmed-rme-segment}       \\
	\hline
	\textbf{\textit{Crate}} & \textbf{Versi/Sumber} \\
	\hline
	\endfirsthead
	\hline
	\textbf{\textit{Crate}} & \textbf{Versi/Sumber} \\
	\hline
	\endhead

	\texttt{base64}         & \texttt{0.22}         \\ \hline
	\texttt{serde}          & \texttt{1}            \\ \hline
	\texttt{serde\_json}    & \texttt{1}            \\ \hline
	\texttt{sha2}           & \texttt{0.10}         \\ \hline
	\texttt{uuid}           & \texttt{1}            \\ \hline
\end{longtable}

\begin{longtable}{|>{\raggedright\arraybackslash}p{0.55\textwidth}|>{\raggedright\arraybackslash}p{0.38\textwidth}|}
	\caption{Dependensi \texttt{decmed-macaroon-auth}}
	\label{tab:dependensi-decmed-macaroon-auth}         \\
	\hline
	\textbf{\textit{Crate}}     & \textbf{Versi/Sumber} \\
	\hline
	\endfirsthead
	\hline
	\textbf{\textit{Crate}}     & \textbf{Versi/Sumber} \\
	\hline
	\endhead

	\texttt{chrono}             & \texttt{0.4}          \\ \hline
	\texttt{serde}              & \texttt{1}            \\ \hline
	\texttt{serde\_json}        & \texttt{1}            \\ \hline
	\texttt{thiserror}          & \texttt{2.0}          \\ \hline
	\texttt{uuid}               & \texttt{1}            \\ \hline
	\texttt{sha2}               & \texttt{0.10}         \\ \hline
	\texttt{hex}                & \texttt{0.4}          \\ \hline
	\texttt{macaroon}           & \texttt{path}         \\ \hline
	\texttt{decmed-rme-segment} & \texttt{path}         \\ \hline
\end{longtable}
 \chapter{Legenda Package Diagram}
\label{lampiran:legenda-package-diagram}

\begin{figure}[ht]
	\centering
	\includegraphics[width=1\textwidth]{resources/appendix/Package Diagram Legend.png}
	\caption{Legenda Package Diagram (Dokumentasi Penulis)}
	\label{fig:legenda-package-diagram}
\end{figure} \chapter{Modul dan Implementasi \textit{Crate} \texttt{decmed-rme-segment}}
\label{lampiran:implementasi-crate-decmed-rme-segment}

\begingroup
\small
\begin{longtable}{|>{\raggedright\arraybackslash}p{0.22\textwidth}|>{\raggedright\arraybackslash}p{0.68\textwidth}|}
	\caption{Daftar Modul dan Implementasi \textit{Crate} \texttt{decmed-rme-segment}}
	\label{tab:modul-crate-decmed-rme-segment}                                                                                                                                                                                                                                                                                                                                                                                                                        \\
	\hline
	\textbf{Modul} & \textbf{Implementasi}                                                                                                                                                                                                                                                                                                                                                                                                                            \\
	\hline
	\endfirsthead
	\hline
	\textbf{Modul} & \textbf{Implementasi}                                                                                                                                                                                                                                                                                                                                                                                                                            \\
	\hline
	\endhead

	\texttt{types}
	               & \texttt{CreateRmeSegmentRequest} (\textit{struct}), \texttt{RmeSegmentData} (\textit{struct}) dengan \texttt{new()} dan \texttt{validate()}; \texttt{ClientEncryptedRmeSegment} (\textit{struct}) dengan \texttt{validate()} dan \texttt{into\_metadata()}; \texttt{RmeSegmentMetadata} (\textit{struct}) dengan \texttt{validate()}; \texttt{CreateRmeSegmentResponse} (\textit{struct}) dengan implementasi \texttt{From<RmeSegmentMetadata>}.
	\\ \hline

	\texttt{category}
	               & \texttt{DatasetCategory} (\textit{enum}), \texttt{FunctionCategory} (\textit{enum}), \texttt{ALL\_DATASET\_CATEGORIES} (\textit{const}), \texttt{ALL\_FUNCTION\_CATEGORIES} (\textit{const}).
	\\ \hline

	\texttt{crypto}
	               & \texttt{canonical\_json()}, \texttt{payload\_hash()}, \texttt{sha256\_hex()}, \texttt{ciphertext\_integrity\_hash\_from\_base64()}.
	\\ \hline

	\texttt{error}
	               & \texttt{SegmentValidationError} (\textit{enum}).
	\\ \hline

	\texttt{administrative\newline \_payload}
	               & \texttt{AdministrativeGeneralPayload} (\textit{struct}) dengan \texttt{validate()} dan \texttt{to\_payload\_value()}; \texttt{administrative\_general\_payload\_from\_value()}, \texttt{administrative\_general\_payload\_from\_fields()}.
	\\ \hline

	\texttt{validation}
	               & \texttt{assert\_valid\_correction\_request()}, \texttt{assert\_valid\_correction\_metadata()}, \texttt{get\_allowed\_function\_categories()}, \texttt{is\_valid\_segment\_category()}, \texttt{assert\_valid\_segment\_category()}, \texttt{assert\_segment\_pair\_consistent()}, \texttt{assert\_no\_plaintext\_medical\_fields()}.
	\\ \hline
\end{longtable}
\endgroup
 \chapter{Modul dan Implementasi \textit{Crate} \texttt{decmed-macaroon-auth}}
\label{lampiran:implementasi-crate-decmed-macaroon-auth}

\begingroup
\small
\begin{longtable}{|>{\raggedright\arraybackslash}p{0.18\textwidth}|>{\raggedright\arraybackslash}p{0.72\textwidth}|}
	\caption{Daftar Modul dan Implementasi \textit{Crate} \texttt{decmed-macaroon-auth}}
	\label{tab:modul-crate-decmed-macaroon-auth}                                                                                                                                                                                                                                                                                                      \\
	\hline
	\textbf{Modul} & \textbf{Implementasi}                                                                                                                                                                                                                                                                                                            \\
	\hline
	\endfirsthead
	\hline
	\textbf{Modul} & \textbf{Implementasi}                                                                                                                                                                                                                                                                                                            \\
	\hline
	\endhead

	\texttt{attenuation}
	               & \texttt{DelegationAttenuationParams} (\textit{struct}), \texttt{attenuate\_macaroon()}.
	\\ \hline

	\texttt{caveats}
	               & \texttt{CaveatKey} (\textit{enum}), \texttt{CaveatValue} (\textit{enum}), \texttt{DecmedCaveat} (\textit{struct}), \texttt{ParsedCaveats} (\textit{struct}), \texttt{parse\_caveat\_line()}, \texttt{format\_dataset\_list()}, \texttt{format\_function\_list()}, \texttt{caveat\_line()}, \texttt{add\_caveat\_to\_macaroon()}.
	\\ \hline

	\texttt{delegation}
	               & \texttt{DelegationStep} (\textit{struct}), \texttt{DelegationChain} (\textit{struct}) dengan \texttt{from\_parsed()} dan \texttt{delegation\_depth()}.
	\\ \hline

	\texttt{delegation\_\newline proof}
	               & \texttt{DelegationProofContext} (\textit{struct}) dengan
	\texttt{from\_token()}, \texttt{from\_last\_delegation\_step()},
	dan \texttt{canonical\_message()}; \texttt{hash\_macaroon\_token()}.
	\\ \hline

	\texttt{effective}
	               & \texttt{AccessMode} (\textit{enum}: \texttt{Read}/\texttt{Write}), \texttt{EffectiveCapability} (\textit{struct}) dengan \texttt{from\_parsed()}, \texttt{allows\_dataset()}, \texttt{allows\_function()}, dan \texttt{is\_expired()}.
	\\ \hline

	\texttt{errors}
	               & \texttt{CaveatVerificationError} (\textit{enum}).
	\\ \hline

	\texttt{issuance}
	               & \texttt{AdminTokenKind} (\textit{enum}: \texttt{Read}/\texttt{Update}), \texttt{InitialAdminPersonnelTokenParams} (\textit{struct}), \texttt{admin\_all\_datasets()}, \texttt{admin\_all\_functions()}, \texttt{admin\_write\_datasets()}, \texttt{admin\_write\_functions()}, \texttt{issue\_admin\_personnel\_token()}.
	\\ \hline

	\texttt{revocation}
	               & \texttt{hash\_token()}, \texttt{token\_revocation\_key()}, \texttt{compute\_revocation\_keys()}.
	\\ \hline

	\texttt{rme\_id}
	               & \texttt{format\_related\_rme\_id()}.
	\\ \hline

	\texttt{verify}
	               & \texttt{SegmentAccessContext} (\textit{struct}), \texttt{TokenVerificationContext} (\textit{struct}), \texttt{VerifiedDecmedToken} (\textit{struct}), \texttt{verify\_macaroon\_signature()}, \texttt{verify\_decmed\_token()}, \texttt{verify\_segment\_access()}.
	\\ \hline

	\texttt{wallet\_proof}
	               & \texttt{WalletProofContext} (\textit{struct}) dengan \texttt{canonical\_message()}, \texttt{WalletSignatureVerifier} (\textit{trait}).
	\\ \hline
\end{longtable}
\endgroup
 \chapter{Alur Implementasi Penegakan Otoritas pada Server PRE}
\label{lampiran:alur-implementasi-penegakan-otoritas-pre}

Alur penegakan otoritas pada \textit{Server} PRE diilustrasikan dalam dua buah flow diagram yang saling berkelanjutan, sebagaimana ditunjukkan pada Gambar \ref{fig:alur-penegakan-otoritas-pre-a} dan \ref{fig:alur-penegakan-otoritas-pre-b}. Gambar \ref{fig:legenda-flowdiag-penegakan-otoritas-pre} memuat legenda bentuk pada \textit{Flow Diagram}.

\begin{figure}[ht]
	\centering
	\includegraphics[width=0.65\textwidth,keepaspectratio]{resources/appendix/legenda flowdiag.png}
	\caption{Legenda bentuk pada \textit{Flow Diagram}}
	\label{fig:legenda-flowdiag-penegakan-otoritas-pre}
\end{figure}

\begin{figure}[ht]
	\centering
	\includegraphics[height=0.8\textheight,keepaspectratio]{resources/appendix/flowdiag-penegakan otorisasi a.png}
	\caption{Alur implementasi penegakan otoritas pada \textit{Server} PRE bagian autentikasi token, \textit{effective capability}, \textit{delegation chain}, dan revokasi}
	\label{fig:alur-penegakan-otoritas-pre-a}
\end{figure}

\begin{figure}[ht]
	\centering
	\includegraphics[height=0.8\textheight,keepaspectratio]{resources/appendix/flowdiag-penegakan otorisasi b.png}
	\caption{Alur implementasi penegakan otoritas pada \textit{Server} PRE bagian validasi \textit{role} dan \textit{subrole}, validasi segmen, \textit{wallet proof}, serta eksekusi PRE, IPFS, dan IOTA}
	\label{fig:alur-penegakan-otoritas-pre-b}
\end{figure}
 \chapter{Implementasi REST API Server PRE}
\label{lampiran:implementasi-rest-api-pre}
{\small
	\begin{longtable}{|p{0.10\textwidth}|p{0.30\textwidth}|p{0.34\textwidth}|p{0.16\textwidth}|}
		\caption{Daftar Endpoint REST API pada Server PRE}
		\label{tab:daftar-endpoint-rest-api-pre}                                                                         \\
		\hline
		\textbf{Method} & \textbf{Endpoint}                                          & \textbf{Fungsi} & \textbf{Status} \\ \hline
		\endfirsthead

		\hline
		\textbf{Method} & \textbf{Endpoint}                                          & \textbf{Fungsi} & \textbf{Status} \\ \hline
		\endhead

		\texttt{GET}    & \texttt{/}
		                & \textit{Health check}.
		                & Eksisting                                                                                      \\ \hline

		\texttt{GET}    & \texttt{/medical-\allowbreak records}
		                & Mengambil daftar rekam medis.
		                & Dimodifikasi                                                                                   \\ \hline

		\texttt{GET}    & \texttt{/medical-\allowbreak record}
		                & Mengambil detail rekam medis.
		                & Eksisting                                                                                      \\ \hline

		\texttt{POST}   & \texttt{/medical-\allowbreak record}
		                & Membuat rekam medis.
		                & Eksisting                                                                                      \\ \hline

		\texttt{PUT}    & \texttt{/medical-\allowbreak record}
		                & Memperbarui rekam medis.
		                & Eksisting                                                                                      \\ \hline

		\texttt{POST}   & \texttt{/medical-\allowbreak record-\allowbreak segment}
		                & Membuat segmen rekam medis.
		                & Baru                                                                                           \\ \hline

		\texttt{GET}    & \texttt{/medical-\allowbreak record-\allowbreak update}
		                & Mengambil data untuk pembaruan RME.
		                & Baru                                                                                           \\ \hline

		\texttt{GET}    & \texttt{/administrative}
		                & Mengambil data administratif pasien.
		                & Eksisting                                                                                      \\ \hline



		\texttt{GET}    & \texttt{/gen/\allowbreak macaroon-\allowbreak key}
		                & Menghasilkan \texttt{RootKey} Macaroon.
		                & Baru                                                                                           \\ \hline

		\texttt{GET}    & \texttt{/gen/sig}
		                & Menghasilkan \textit{signature}.
		                & Eksisting                                                                                      \\ \hline

		\texttt{GET}    & \texttt{/gen/\allowbreak proxy-\allowbreak address}
		                & Menghasilkan dan mendaftarkan alamat proxy.
		                & Eksisting                                                                                      \\ \hline

		\texttt{POST}   & \texttt{/nonce}
		                & Mengambil nonce.
		                & Eksisting                                                                                      \\ \hline

		\texttt{POST}   & \texttt{/keys}
		                & Menyimpan kunci.
		                & Eksisting                                                                                      \\ \hline

		\texttt{POST}   & \texttt{/revocations/\allowbreak patient}
		                & Mencabut hak akses dan turunannya.
		                & Baru                                                                                           \\ \hline

		\texttt{POST}   & \texttt{/delegations/\allowbreak attenuate}
		                & Melakukan atenuasi token Macaroon sesuai \textit{request}.
		                & Baru                                                                                           \\ \hline
	\end{longtable}
} \chapter{Implementasi Smart Contract Sistem DecMed}
\label{lampiran:implementasi-smart-contract-sistem-decmed}

\let\textunderscorewithoutbreak\_
\newcommand{\breakabletextunderscore}{\textunderscorewithoutbreak\allowbreak{}}
\newcolumntype{Y}[1]{>{\let\_\breakabletextunderscore\raggedright\arraybackslash}p{#1}<{\let\_\textunderscorewithoutbreak}}

\begin{longtable}{|Y{0.36\textwidth}|Y{0.56\textwidth}|}
	\caption{Daftar Module Smart Contract Sistem DecMed}
	\label{tab:daftar-module-smart-contract-decmed}                                                                                                                                                    \\
	\hline
	\textbf{Nama Module}                 & \textbf{Deskripsi}                                                                                                                                          \\
	\hline
	\endfirsthead
	\hline
	\textbf{Nama Module}                 & \textbf{Deskripsi}                                                                                                                                          \\
	\hline
	\endhead

	\texttt{decmed::admin}               & Modul global admin untuk registrasi rumah sakit, aktivasi admin fasyankes, dan daftar rumah sakit.                                                          \\ \hline
	\texttt{decmed::hospital\_personnel} & Modul akun tenaga fasyankes, aktivasi, akses baca/tulis, \textit{delegation}, dan \textit{audit delegation}.                                                \\ \hline
	\texttt{decmed::patient}             & Modul akun pasien, \textit{grant} akses ke \texttt{AdministrativePersonnel}, log akses, \textit{audit delegation}, dan pembacaan metadata RME milik pasien. \\ \hline
	\texttt{decmed::proxy}               & Modul yang dipakai \textit{server} PRE untuk baca/tulis metadata RME.                                                                                       \\ \hline
	\texttt{decmed::shared}              & Modul penyedia \texttt{GlobalAdminCap} dan \texttt{ProxyCap}                                                                                                \\ \hline
\end{longtable}

\begin{longtable}{|Y{0.45\textwidth}|Y{0.47\textwidth}|}
	\caption{Daftar Method pada Module \texttt{decmed::admin}}
	\label{tab:daftar-method-decmed-admin}                                                                          \\
	\hline
	\textbf{Nama Method}             & \textbf{Deskripsi}                                                           \\
	\hline
	\endfirsthead
	\hline
	\textbf{Nama Method}             & \textbf{Deskripsi}                                                           \\
	\hline
	\endhead

	\texttt{create\_activation\_key} & Membuat data rumah sakit dan activation key pertama untuk admin rumah sakit. \\ \hline
	\texttt{update\_activation\_key} & Memperbarui activation key dan metadata admin rumah sakit.                   \\ \hline
	\texttt{get\_hospitals}          & Mengambil daftar rumah sakit untuk client kementerian/global admin.          \\ \hline
\end{longtable}

\begin{longtable}{|Y{0.45\textwidth}|Y{0.47\textwidth}|}
	\caption{Daftar Method pada Module \texttt{decmed::hospital\_personnel}}
	\label{tab:daftar-method-decmed-hospital-personnel}                                                                                \\
	\hline
	\textbf{Nama Method}                      & \textbf{Deskripsi}                                                                     \\
	\hline
	\endfirsthead
	\hline
	\textbf{Nama Method}                      & \textbf{Deskripsi}                                                                     \\
	\hline
	\endhead

	\texttt{create\_activation\_key}          & Admin fasyankes membuat activation key untuk tenaga administratif atau medis.          \\ \hline
	\texttt{use\_activation\_key}             & Menandai activation key tenaga fasyankes sebagai sudah digunakan.                      \\ \hline
	\texttt{signup}                           & Registrasi akun tenaga fasyankes ke address IOTA dan menyimpan metadata administratif. \\ \hline
	\texttt{is\_account\_registered}          & Mengecek apakah akun tenaga fasyankes sudah aktif/terdaftar.                           \\ \hline
	\texttt{get\_account\_state}              & Menentukan status \textit{onboarding} tenaga fasyankes.                                \\ \hline
	\texttt{get\_account\_info}               & Mengambil metadata akun, \textit{role}, \textit{subrole}, dan metadata rumah sakit.    \\ \hline
	\texttt{update\_account\_activation\_key} & Admin fasyankes mengganti activation key personel.                                     \\ \hline
	\texttt{update\_administrative\_metadata} & Memperbarui profil administratif tenaga fasyankes.                                     \\ \hline
	\texttt{get\_hospital\_personnels}        & Mengambil daftar personel di bawah admin fasyankes.                                    \\ \hline
	\texttt{get\_delegatee\_candidates}       & Mengambil kandidat personel yang bisa menerima delegasi akses.                         \\ \hline
	\texttt{cleanup\_read\_access}            & Membersihkan akses baca yang sudah kedaluwarsa.                                        \\ \hline
	\texttt{cleanup\_update\_access}          & Membersihkan akses \textit{update} yang sudah kedaluwarsa.                             \\ \hline
	\texttt{get\_read\_access}                & Mengambil daftar \textit{capability} baca milik tenaga fasyankes.                      \\ \hline
	\texttt{get\_update\_access}              & Mengambil daftar \textit{capability} \texttt{update} milik tenaga fasyankes.           \\ \hline
	\texttt{create\_delegated\_access}        & Membuat \textit{delegated access} untuk personel lain.                                 \\ \hline
\texttt{create\_medical\_record}          & Membuat entri RME blob besar. (\textit{legacy})                                        \\ \hline
\end{longtable}

\begin{longtable}{|Y{0.45\textwidth}|Y{0.47\textwidth}|}
	\caption{Daftar Method pada Module \texttt{decmed::patient}}
	\label{tab:daftar-method-decmed-patient}                                                                                   \\
	\hline
	\textbf{Nama Method}                      & \textbf{Deskripsi}                                                             \\
	\hline
	\endfirsthead
	\hline
	\textbf{Nama Method}                      & \textbf{Deskripsi}                                                             \\
	\hline
	\endhead

	\texttt{signup}                           & Registrasi akun pasien dan metadata administratif terenkripsi.                 \\ \hline
	\texttt{is\_account\_registered}          & Mengecek apakah \textit{address} pasien  terdaftar.                            \\ \hline
	\texttt{get\_account\_state}              & Menentukan status \textit{onboarding} pasien.                                  \\ \hline
	\texttt{get\_account\_info}               & Mengambil metadata administratif pasien.                                       \\ \hline
	\texttt{update\_administrative\_metadata} & Memperbarui metadata administratif pasien.                                     \\ \hline
	\texttt{get\_hospital\_personnel\_info}   & Mengambil info publik tenaga fasyankes saat pasien memproses QR.               \\ \hline
	\texttt{create\_access}                   & Memberi akses baca/tulis kepada tenaga fasyankes.                              \\ \hline
	\texttt{get\_access\_log}                 & Mengambil log \textit{grant} akses pasien.                                     \\ \hline
	\texttt{get\_delegation\_audit\_log}      & Mengambil log \textit{audit trail} \textit{delegation} pasien.                 \\ \hline
	\texttt{get\_medical\_records}            & Mengambil daftar metadata RME pasien.                                          \\ \hline
	\texttt{get\_medical\_record}             & Mengambil satu metadata RME berdasarkan indeks.                                \\ \hline
	\texttt{revoke\_access}                   & Mencabut akses tenaga fasyankes dan memicu \textit{cascade revoke delegation}. \\ \hline
\end{longtable}

\begin{longtable}{|Y{0.45\textwidth}|Y{0.47\textwidth}|}
	\caption{Daftar Method pada Module \texttt{decmed::proxy}}
	\label{tab:daftar-method-decmed-proxy}                                                                                                          \\
	\hline
	\textbf{Nama Method}                          & \textbf{Deskripsi}                                                                              \\
	\hline
	\endfirsthead
	\hline
	\textbf{Nama Method}                          & \textbf{Deskripsi}                                                                              \\
	\hline
	\endhead

	\texttt{create\_capability}                   & Membuat \texttt{ProxyCap} untuk address service proxy.                                          \\ \hline
	\texttt{create\_medical\_record\_segment}     & Menambahkan metadata RME tersegmentasi ke daftar medical metadata pasien.                       \\ \hline
	\texttt{get\_administrative\_data}            & Mengambil metadata administratif pasien untuk kebutuhan akses/dekripsi.                         \\ \hline
	\texttt{get\_hospital\_personnel\_role}       & Mengecek role tenaga fasyankes saat setup token/key access.                                     \\ \hline
	\texttt{get\_hospital\_personnel\_auth\_info} & Mengecek \textit{role} dan \textit{subrole}.                                                    \\ \hline
	\texttt{get\_medical\_records}                & Mengambil daftar metadata RME pasien dan menyaringnya sebagai segmen berdasarkan cakupan akses. \\ \hline
	\texttt{get\_medical\_record}                 & Mengambil satu metadata segmen dan data pendukung untuk pembacaan \textit{payload}.             \\ \hline
	\texttt{is\_patient\_registered}              & Memvalidasi bahwa pasien terdaftar                                                              \\ \hline
	\texttt{create\_medical\_record}              & Membuat entri RME blob besar. (\textit{legacy})                                                 \\ \hline
	\texttt{update\_medical\_record}              & Melakukan edit entri RME blob besar. (\textit{legacy})                                          \\ \hline
	\texttt{get\_medical\_record\_update}         & Mengambil data RME untuk proses edit entri blob besar. (\textit{legacy})                        \\ \hline
	\texttt{get\_delegation\_access\_snapshot}    & Mengambil \textit{snapshot} metadata akses delegasi.                                            \\ \hline
\end{longtable}
 \chapter{Tangkapan Layar Implementasi \textit{Patient Client}}
\label{lampiran:implementasi-patient-client}

Semua tangkapan layar berikut merupakan dokumentasi penulis.

\begin{figure}[ht]
	\centering
	\includegraphics[width=0.35\textwidth]{resources/appendix/client/patient/home.png}
	\caption{Halaman \textit{Home} daftar RME.}
	\label{fig:halaman-home-daftar-rme}
\end{figure}

\begin{figure}[p]
	\centering
	\includegraphics[width=0.35\textwidth]{resources/appendix/client/patient/home detail RME.png}
	\caption{Halaman detail RME.}
	\label{fig:halaman-detail-rme}
\end{figure}

\begin{figure}[p]
	\centering
	\includegraphics[width=0.35\textwidth]{resources/appendix/client/patient/delegation audit.png}
	\caption{Halaman audit delegasi RME.}
	\label{fig:halaman-detail-audit-delegasi}
\end{figure}



\begin{figure}[p]
	\centering
	\includegraphics[width=0.35\textwidth]{resources/appendix/client/patient/scan page.png}
	\caption{Halaman \textit{scan} QR.}
	\label{fig:halaman-detail-scan-qr}
\end{figure}

\begin{figure}[p]
	\centering
	\includegraphics[width=0.35\textwidth]{resources/appendix/client/patient/scan page confirm access.png}
	\caption{Halaman konfirmasi pemberian akses awal.}
	\label{fig:halaman-konfirmasi-akses-awal}
\end{figure}
 \chapter{Tangkapan Layar Implementasi \textit{Hospital Client}}
\label{lampiran:implementasi-hospital-client}

Semua tangkapan layar berikut merupakan dokumentasi penulis.

\begin{figure}[htbp]
	\centering
	\includegraphics[width=0.95\textwidth]{resources/appendix/client/hospital/dashboard.png}
	\caption{Halaman \textit{dashboard} \textit{MedicalPersonnel}.}
	\label{fig:halaman-dashboard-medical-personnel}
\end{figure}

\begin{figure}[htbp]
	\centering
	\includegraphics[width=0.95\textwidth]{resources/appendix/client/hospital/delegate 1.png}
	\caption{Halaman delegasi akses bagian 1.}
	\label{fig:halaman-delegasi-akses-bagian-1}
\end{figure}

\begin{figure}[htbp]
	\centering
	\includegraphics[width=0.95\textwidth]{resources/appendix/client/hospital/delegate 2.png}
	\caption{Halaman delegasi akses bagian 2.}
	\label{fig:halaman-delegasi-akses-bagian-2}
\end{figure}

\begin{figure}[htbp]
	\centering
	\includegraphics[width=0.95\textwidth]{resources/appendix/client/hospital/dashboard pasien.png}
	\caption{Halaman daftar RME pasien pada \textit{Hospital Client}.}
	\label{fig:halaman-daftar-rme-pasien-hospital-client}
\end{figure}

\begin{figure}[htbp]
	\centering
	\includegraphics[width=0.95\textwidth]{resources/appendix/client/hospital/dashboard detail rme.png}
	\caption{Halaman detail RME pada \textit{Hospital Client}.}
	\label{fig:halaman-detail-rme-hospital-client}
\end{figure}

\begin{figure}[htbp]
	\centering
	\includegraphics[width=0.95\textwidth]{resources/appendix/client/hospital/akses tulis.png}
	\caption{Halaman alur akses tulis pada \textit{Hospital Client}.}
	\label{fig:halaman-penulisan-segmen-rme-hospital-client}
\end{figure}
  
\end{document}
